		\documentclass[ALCO,ThmDefs,Unicode,epreuves]{cedram}
\OneNumberAllTheorems

%\def\exemname{Example}

\usepackage{float}
\usepackage{ytableau}
\usepackage{young}


\newenvironment{Case}[1]{\smallskip\noindent\emph{Case #1}\endgraf
 \nobreak\noindent\ignorespaces}{\endgraf\smallskip}


\DeclareMathOperator{\orm}{o}
\DeclareMathOperator{\cp}{cp}
\DeclareMathOperator{\Ini}{Ini}
\DeclareMathOperator{\perm}{tblr}
\DeclareMathOperator{\lex}{lex}
\DeclareMathOperator{\row}{row}
\DeclareMathOperator{\col}{col}
\DeclareMathOperator{\word}{btlr}




\newcommand{\opp}{^{\orm}}
\newcommand{\leftright}{_{\mathsf{LR}}}

\newcommand{\lside}{_{\mathsf{L}}}
\newcommand{\rside}{_{\mathsf{R}}}
\newcommand{\lessthan}[2]{#1\,{_{<\,#2}}}
\newcommand{\morethan}[2]{#1\,{_{>\,#2}}}
\newcommand{\Lessthan}[2]{#1\,{_{\leqslant\,#2}}}
\newcommand{\Morethan}[2]{#1\,{_{\geqslant\,#2}}}
\newcommand{\jdt}{\mathrm{jdt}}
\newcommand{\citex}[2]{\cite[#1]{#2}}

\def\ttab(#1){\text{\tiny\ytableaushort{#1}}}

\DeclareMathOperator{\CSTD}{CStd}
\DeclareMathOperator{\STD}{Std}
\DeclareMathOperator{\shape}{Shape}
\DeclareMathOperator{\neb}{-neb}

\DeclareMathOperator{\ind}{ind}
\DeclareMathOperator{\SD}{SD}
\DeclareMathOperator{\SA}{SA}
\DeclareMathOperator{\WD}{WD}
\DeclareMathOperator{\WA}{WA}
\DeclareMathOperator{\D}{D}
\DeclareMathOperator{\A}{A}
\DeclareMathOperator{\RS}{RS}
\DeclareMathOperator{\RSP}{P}
\DeclareMathOperator{\RSQ}{Q}
\DeclareMathOperator{\Tab}{Tab}



%%%%%%%%% a placer avant \begin{document}


%%%%%%%%%%%%%%%%%%%%%%%%%%%%%%%%%%%%%
\newcommand*{\mk}{\mkern -1mu}
\newcommand*{\Mk}{\mkern -2mu}
\newcommand*{\mK}{\mkern 1mu}
\newcommand*{\MK}{\mkern 2mu}

\hypersetup{urlcolor=purple, linkcolor=blue, citecolor=red}


\newcommand*{\romanenumi}{\renewcommand*{\theenumi}{\roman{enumi}}}
\newcommand*{\Romanenumi}{\renewcommand*{\theenumi}{\Roman{enumi}}}
\newcommand*{\alphenumi}{\renewcommand*{\theenumi}{\alph{enumi}}}
\newcommand*{\Alphenumi}{\renewcommand*{\theenumi}{\Alph{enumi}}}
%%%%%%%%%%%%%%%%%%%%%%%%%%%%%%%%%%



%%%%%%%%%%%%%%%%%%%%%%%%%%%%%%%%%%%%%%%%%%%
%%%%% Auteur

\author{\firstname{Van} \middlename{Minh} \lastname{Nguyen}}

\address{School of Mathematics and Statistics\\
University of Sydney \\
NSW 2006, Australia}

\email{Van.Nguyen@sydney.edu.au}

%%%%% Sujet

\subjclass{05E10, 20C08}

\keywords{Coxeter groups, Hecke algebras, $W$-graphs, Kazhdan--Lusztig polynomials, cells}


%%%%% Gestion

\DOI{10.5802/alco.91}
\datereceived{2018-07-19}
\daterevised{2019-06-06}
\dateaccepted{2019-07-21}


%%%%% Titre et résumé

\title[Type $A$ admissible cells are Kazhdan--Lusztig]{Type $A$ admissible cells are Kazhdan--Lusztig}

\begin{abstract}
Admissible $W\!$-graphs were defined and combinatorially characterized by
Stembridge in \cite{stem:addwgraph}. The theory of admissible $W\!$-graphs
was motivated by the need to construct $W\!$-graphs for Kazhdan--Lusztig
cells, which play an important role in the representation theory of Hecke algebras,
without computing Kazhdan--Lusztig polynomials. In this paper, we shall show
that type $A$-admissible $W\!$-cells are Kazhdan--Lusztig as conjectured
by Stembridge in his original paper.

\end{abstract}

%%%%%%%%%%%%%%%%%%%%%%%%%%%%%%%%%%%%%%%%%%%




\begin{document}


\maketitle

\section{Introduction}
\label{intro}
\looseness-1
Let $(W,S)$ be a Coxeter system and $\mathcal{H}(W)$ its Hecke algebra over
$\mathbb{Z}[q,q^{-1}]$, the ring of Laurent polynomials in the indeterminate $q$.
We are interested in representations of $W$ and $\mathcal{H}(W)$ that can be described
by combinatorial objects, namely~$W\!$-graphs. In particular, we are interested in~$W\!$-graphs
corresponding to Kazhdan--Lusztig left cells.

In principle, when computing left cells one encounters the problem of having to compute a large number of
Kazhdan--Lusztig polynomials before any explicit description of their~$W\!$-graphs can be given.
In \cite{stem:addwgraph}, Stembridge introduced \emph{admissible\/} $W\!$-graphs; these can
be described combinatorially and can be constructed without calculating Kazhdan--Lusztig polynomials.
Moreover, the $W\!$-graphs corresponding to Kazhdan--Lusztig left cells are
admissible. Stembridge showed in~\cite{stem:finitewgraph} that for any given finite
$W$ there are only finitely many stongly connected admissible~$W\!$-graphs. It
was conjectured by Stembridge that in type $A$ all strongly connected admissible
$W$-graphs are isomorphic to Kazhdan--Lusztig left cells. In this paper we
complete the proof of Stembridge's conjecture.

We shall work with \emph{$S$-coloured graphs} (as defined in
Section~3 below), of which $W\!$-graphs are examples. These graphs
have both edges (bi-directional) and arcs (uni-directional). A
\emph{cell} in such a graph $\Gamma$ is by definition a strongly
connected component of~$\Gamma$, and a \emph{simple part} is a
connected component of the graph obtained by removing all arcs that are
not edges and edges of weight greater than~1. A \emph{simple component}
is the full subgraph spanned by a simple part. Stembridge introduced
$W\!$-\emph{molecular graphs}, defined by conditions that are weaker
than those for admissible $W\!$-graphs. Simple components of molecular
graphs are called \emph{molecules}.

By an \emph{admissible $W\!$-cell} we mean a cell in an admissible
$W\!$-graph (rather than a cell in a $W\!$-molecular graph).

In~\cite{ChuMic:typeAMol}, Chmutov established the first step towards
the proof of Stembridge's conjecture, showing that if $(W,S)$ is of
type $A_{n-1}$ then the simple part of a $W\!$-molecule is isomorphic
to the simple part of a Kazhdan--Lusztig left cell. The proof made use of the
axiomatisation of dual equivalence graphs on standard tableaux generated
by dual Knuth equivalence relations, given in an earlier paper by
Assaf~\cite{assaf:dualequigraphs}.

\looseness-1
By Chmutov's result each molecule of an arbitrary admissible $W\!$-graph of
type~$A_{n-1}$ is associated with a partition of $n$, and the vertices
within each molecule are parametrized by the set of standard tableaux of
the corresponding shape. The Bruhat order on $W$ induces a partial order on
these tableaux, and this extends naturally to a partial order on the set of all
$n$-box standard tableaux (see Section~\ref{sec:8} below), thus giving rise
to a partial order on the vertex set of the graph. We are able to use the
combinatorics of tableaux to show that type $A$ admissible $W\!$-graphs
are \emph{ordered}, in the sense of Definition~\ref{orderedWgraphdef}: if an
arc has tail corresponding to a standard tableau~$t$ and head corresponding
to a standard tableau $u$ then either $u<t$, or else head and tail
belong to the same molecule and $u=st$ for some simple transposition~$s$.
This property of admissible $W\!$-graphs is the key to our proof of the
conjecture of Stembridge.

The proof of Proposition~\ref{monomolecularadmcellsareKL} furnishes an
algorithm for computation of $W\!$-graphs for left cells in type~$A_{n-1}$
(avoiding the computation of Kazhdan--Lusztig polynomials). This has been
implemented in Magma and checked for a variety of partitions $\lambda$
with $n\le 16$, and module dimension up to 171600 (for $\lambda=(5,5,3,3)$).
The Magma code is available on request.

We organize the paper in the following sections. Section~\ref{sec:2} and
Section~\ref{sec:2'} deal
with the background on Coxeter groups and the corresponding Hecke algebras.
In Section~\ref{sec:3} the definition and properties of $W\!$-graphs are
recalled.
In Section~\ref{sec:5}, we recall the definitions of admissible $W\!$-graphs
and molecules and how these can be characterized combinatorially.
Section~\ref{sec:6} presents combinatorics of tableaux and the relationship
between Kazhdan--Lusztig left cells, dual Knuth equivalence classes and
admissible molecules. We introduce the
paired dual Knuth equivalence relation in Section~\ref{sec:8}. In
Section~\ref{sec:8ex}, we prove the first main result, namely that for
type $A_{n-1}$, all admissible $W\!$-graphs are ordered. The proof that
type $A$ admissible $W\!$-cells are isomorphic to Kazhdan--Lusztig left
cells is completed in Section~\ref{sec:9}.

\section{Coxeter groups}
\label{sec:2}
Let $(W,S)$ be a Coxeter system and $l$ the length function on $W$.
The Coxeter group $W$ comes equipped with the left weak
order, the right weak order and the Bruhat order, respectively denoted by
$\leqslant\lside$, $\leqslant\rside$ and $\leqslant$, and defined as follows.

\begin{defi}\label{Orders_on_W}
The \emph{left weak order} is the partial order on $W$ generated by
the relations $x\leqslant\lside y$ for all $x,\,y\in W$
with $l(x)<l(y)$ and $yx^{-1}\in S$.

The \emph{right weak order} is the partial order on $W$ generated by
the relations $x\leqslant\rside y$ for all $x\,y\in W$
with $l(x)<l(y)$ and $x^{-1}y\in S$.

The \emph{Bruhat order} is the partial order on $W$ generated by
the relations $x\leqslant y$ for all $x,\,y\in W$
with $l(x)<l(y)$ and $yx^{-1}$ conjugate to an element of~$S$.
\end{defi}

Observe that the weak orders are characterized by the property that
$x\leqslant\rside xy$ and $y\leqslant\lside xy$ whenever $l(xy)=l(x)+l(y)$.

For each $J\subseteq S$ let $W_{J}$ be the (standard parabolic)
subgroup of $W$ generated by~$J$, and let $D_{J}$ the set of
distinguished (or minimal) representatives of the left cosets of
$W_{J}$ in~$W$. Thus each $w \in W$ has a unique factorization
$w = du$ with $d \in D_{J}$ and $u \in W_{J}$, and $l(du) = l(d) +l(u)$
holds for all $d \in D_{J}$ and $u \in W_{J}$. It is easily seen
that $D_J$ is an ideal of $(W,\leqslant\lside)$, in the sense that if
$w\in D_J$ and $v\in W$ with $v\leqslant\lside w$ then $v\in D_J$.

If $W_J$ is finite then we denote the longest element of $W_{J}$
by~$w_{J}$. By \cite[Lemma~2.2.1]{gecpfei:charHecke}, if $W$ is finite then
$D_{J} = \{d\in W\mid d\leqslant\lside d_J\}$, where $d_{J}$ is
the unique element in $D_{J} \cap w_{S}W_{J}$.

\section{Hecke algebras}
\label{sec:2'}

Let $\mathcal{A} = \mathbb{Z}[q, q^{-1}]$, the ring of Laurent
polynomials with integer coefficients in the indeterminate $q$,
and let $\mathcal{A}^{+} = \mathbb{Z}[q]$.
The Hecke algebra of a Coxeter system $(W, S)$, denoted by
$\mathcal{H}(W)$ or simply by $\mathcal{H}$,
is an associative $\mathcal{A}$-algebra with $\mathcal{A}$-basis
$\{\,H_w\mid w\in W\,\}$ satisfying
\begin{align*}
H^{2}_{s} &= 1 + (q -q^{-1})H_{s} \quad \text{for all $s \in S$},\\
H_{xy}&= H_{x}H_{y} \quad
\text{for all $x,\,y \in W$ with $l(xy)=l(x)+l(y)$.}
\end{align*}
We let $a \mapsto \overline{a}$ be the involutory automorphism of
$\mathcal{A}=\mathbb{Z}[q,q^{-1}]$ defined by $\overline{q}=q^{-1}$.
It is well known that this extends to an involutory automorphism of
$\mathcal{H}$ satisfying
\begin{equation*}
\overline{H_{s}} = H_{s}^{-1} = H_{s} - (q - q^{-1}) \quad
\text{for all $s \in S$}.
\end{equation*}

If $J\subseteq S$ then $\mathcal{H}(W_{J})$, the Hecke algebra
associated with the Coxeter system $(W_{J},J)$, is isomorphic to
the subalgebra of $\mathcal{H}(W)$ generated by
$\{\,H_{s} \mid s \in J\,\}$. We shall identify $\mathcal{H}(W_{J})$
with this subalgebra.

\section{\texorpdfstring{$W$}{W}-graphs}
\label{sec:3}

Given a set $S$, we define an \emph{$S$-coloured graph} to be
a triple $\Gamma=(V, \mu, \tau)$ consisting of a set $V\!$, a
function $\mu\colon V \times V\to \mathbb{Z}[q,q^{-1}]$ and a function
$\tau$ from $V$ to $\mathcal{P}(S)$, the power set of $S$.
The elements of $V$ are the \emph{vertices} of $\Gamma$, and if
$v\in V$ then $\tau(v)$ is the \emph{colour} of the vertex.
To interpret $\Gamma$ as a (directed) graph, we adopt the convention that
if $v,\,u\in V$ then $(v,u)$ is an arc from $v$ to $u$ of $\Gamma$
if and only if $\mu(u,v) \neq 0$ and $\tau(u)\nsubseteq\tau(v)$, and
$\{v,u\}$ is an edge of $\Gamma$ if and only if $(v,u)$ and
$(u,v)$ are both arcs. We call $\mu(u,v)$ the \emph{weight}
of the arc~$(v,u)$. An edge $\{u,v\}$ is said to be \emph{symmetric}
if $\mu(u,v)=\mu(v,u)$, and \emph{simple} if $\mu(u,v)=\mu(v,u)=1$.

An $S$-coloured graph is \emph{reduced} if $\mu(u,v)=0$ whenever
$\tau(u)\subseteq\tau(v)$. Except when stated otherwise, all $S$-coloured
graphs we consider are assumed to be reduced.

If $(W,S)$ is a Coxeter system, then a \emph{$W\!$-graph} is
an $S$-coloured graph $\Gamma=(V, \mu, \tau)$ such that the free
$\mathcal{A}$-module with basis~$V$ admits an $\mathcal{H}$-module
structure satisfying
\begin{equation}\label{wgraphdef}
 H_{s}v = \begin{cases}
 -q^{-1}v \quad &\text{if $s \in \tau(v)$}\\
 qv + \sum_{\{u \in V \mid s \in \tau(u)\}}\mu(u,v)u
 \quad &\text{if $s \notin \tau(v)$},
 \end{cases}
\end{equation}
for all $s \in S$ and $v \in V\!$.
We write $M_\Gamma$ for the $\mathcal{H}$-module given by
the $W\!$-graph~$\Gamma$ in this way.

Note that $\mu(u,v)$ appears in Eq.~\eqref{wgraphdef} only if there is an $s$
with $s\in\tau(u)$ and $s\notin\tau(v)$. Thus redefining $\mu(u,v)=0$
whenever $\tau(u)\subseteq \tau(v)$ does not alter~$M_\Gamma$.
So a non-reduced $S$-coloured graph is a $W$-graph if and only if the
corresponding reduced $S$-coloured graph is a $W$-graph.

Since $M_\Gamma$
is $\mathcal{A}$-free with basis~$V$ it admits an
$\mathcal{A}$-semilinear involution $\alpha \mapsto \overline{\alpha}$,
uniquely determined by the condition
that $\overline v=v$ for all $v\in V$. We call this the
\emph{bar involution} on $M_\Gamma$. It is a consequence of
\eqref{wgraphdef} that $\overline{h\alpha}=\overline{h}\overline{\alpha}$
for all $h\in \mathcal{H}$ and $\alpha \in M_\Gamma$.

We shall sometimes write $\Gamma(V)$ for the $W\!$-graph with vertex
set~$V\!$, if the functions $\mu$ and~$\tau$ are clear from the context.

If $\Gamma=(V,\mu,\tau)$ is an $S$-coloured graph and $J\subseteq S$ then
the $W_J$-restriction of $\Gamma$ is defined to be the $J$-coloured graph
$\Gamma_{J} = (V,\mu_J,\tau_J)$
where $\tau_J(v)=\tau(v)\cap J$ for all $v\in V$ and
\[
\mu_J(u,v)=
\begin{cases}
\mu(u,v)&\text{ if $\tau_J(u)\nsubseteq \tau_J(v)$,}\\
0&\text{ if $\tau_J(u)\subseteq \tau_J(v)$.}
\end{cases}
\]
It is clear that if $J\subseteq S$ and $\Gamma=(V,\mu,\tau)$ is a
$W\!$-graph then $\Gamma_J=(V,\mu_J,\tau_J)$ is a $W_J$-graph.

Following~\cite{kazlus:coxhecke}, define a preorder
$\leqslant_{\Gamma}$ on $V$ as follows: $u \leqslant_{\Gamma} v$
if there exists a sequence of vertices $u = x_{0},x_{1},
\ldots, x_{m} = v$ such that $\tau(x_{i-1}) \nsubseteq
\tau(x_{i})$ and $\mu(x_{i-1},x_{i}) \neq 0$ for all $i \in
[1,m]$. That is, $u \leqslant_{\Gamma} v$ if there is a directed path
from $v$ to $u$ in~$\Gamma$. Let $\sim_{\Gamma}$ be
the equivalence relation determined by this preorder.
The equivalence classes with respect to $\sim_{\Gamma}$
are called the \emph{cells} of $\Gamma$. That is, the cells are
the strongly connected components of the directed graph $\Gamma$.
Each equivalence class,
regarded as a full subgraph of $\Gamma$, is itself a $W\!$-graph,
with the $\mu$ and $\tau$ functions being the restrictions
of those for $\Gamma$. The preorder $\leqslant_{\Gamma}$ induces a partial
order on the set of cells: if $\mathcal{C}$ and $\mathcal{C}'$ are
cells, then $\mathcal{C} \leqslant_{\Gamma} \mathcal{C}'$ if $u \leqslant_{\Gamma} v$
for some $u \in \mathcal{C}$ and $v \in \mathcal{C}'$.

It follows readily from \eqref{wgraphdef} that a subset of
$V$ spans a $\mathcal{H}(W)$-submodule of $M_{\Gamma}$ if
and only if it is $\Gamma$-closed, in the sense that for every vertex
$v$ in the subset, each $u\in V$ satisfying $\mu(u,v)\ne 0$ and
$\tau(u)\nsubseteq\tau(v)$ is also in the subset.
Thus $U\subseteq V$ is a $\Gamma$-closed subset of $V$ if and
only if $U=\bigcup_{v\in U}\{\,u\in V\mid u\leqslant_{\Gamma}v\,\}$.
Clearly, a subset of $V$ is $\Gamma$-closed if and only if it is the
union of cells that form an ideal with respect to the partial order
$\leqslant_\Gamma$ on the set of cells.

Suppose that $U$ is a $\Gamma$-closed subset of $V$, and let
$\Gamma(U)$ and $\Gamma(V \setminus U)$ be the full subgraphs of
$\Gamma$ induced by $U$ and $V \setminus U$, with arc weights
and vertex colours inherited from $\Gamma$. Then $\Gamma(U)$ and
$\Gamma(V \setminus U)$ are themselves $W\!$-graphs, and
\[
M_{\Gamma(V \setminus U)} \cong M_{\Gamma(V)}/M_{\Gamma(U)}
\]
as $\mathcal{H}(W)$-modules.

We end this section by recalling the original Kazhdan--Lusztig $W\!$-graph
for the regular representation of $\mathcal{H}(W)$.
For each $w \in W$, define the sets
\begin{align*}
\mathcal{L}(w) &= \{s \in S \mid l(sw) < l(w)\}\\
\intertext{and}
\mathcal{R}(w) & = \{s \in S \mid l(ws) < l(w)\},
\end{align*}
the elements of which are called the left descents of $w$ and the right descents
of $w$, respectively. Kazhdan and Lusztig give a recursive procedure that defines
polynomials $P_{y,w}$ whenever $y,\,w \in W$ and $y < w$. These polynomials satisfy
$\deg P_{y,w}\leqslant\frac12(l(w)-l(y)-1)$, and $\mu_{y,w}$ is defined to be
the leading coefficient of $P_{y,w}$ if the degree is $\frac12(l(w)-l(y)-1)$,
or 0 otherwise.

Define $W\opp=\{\,w\opp\mid w\in W\,\}$ to be the group opposite to $W\!$, and
observe that $(W\times W\opp\!,\,S\sqcup S\opp)$ is a Coxeter system.
Kazhdan and Lusztig show that if $\mu$ and $\bar\tau$ are defined by the formulas
\begin{align*}
\mu(w,y)=\mu(y,w) &= \begin{cases}
 \mu_{y,w} &\quad \text{if $y < w$}\\
 \mu_{w,y} &\quad \text{if $w < y$}
 \end{cases}\\
\bar\tau(w) &= \mathcal{L}(w)\sqcup\mathcal{R}(w)\opp
\end{align*}
then $(W,\mu,\bar\tau)$ is a $(W\times W\opp)$-graph (usually not reduced).
Thus $M=\mathcal{A}W$ may be regarded as an
$(\mathcal{H},\mathcal{H})$-bimodule. Furthermore, the construction
produces an explicit $(\mathcal{H},\mathcal{H})$-bimodule isomorphism
$M\cong\mathcal{H}$.

It follows easily from the definition of $\mu_{y,w}$ that $\mu(y,w)\neq 0$ only
if $l(w)-l(y)$ is odd; thus $(W,\mu,\bar\tau)$ is a bipartite graph. The non-negativity
of all coefficients of the Kazhdan--Lusztig polynomials, conjectured in
\cite{kazlus:coxhecke}, has been proved by Elias
and Williamson in~\cite{bengeogedi:howgesoerge}.

Since $W$ and $W\opp$ are standard parabolic subgroups of~$W\times W\opp$, it follows
that $\Gamma=(W,\mu,\tau)$ is a $W\!$-graph and $\Gamma\opp=(W,\mu,\tau\opp)$ is a
$W\opp\!$-graph, where $\tau$ and $\tau\opp$ are defined by
$\tau(w)=\mathcal{L}(w)$ and $\tau\opp(w)=\mathcal{R}(w)\opp$, for all $w\in W$.

In accordance with the theory described above, there are preorders on~$W$
determined by the $(W\times W\opp)$-graph structure, the $W\!$-graph structure
and the $W\opp\!$-graph structure. We call these the \emph{two-sided preorder}
(denoted by $\preceq\leftright$), the \emph{left preorder} ($\preceq\lside$)
and the \emph{right preorder} ($\preceq\rside$). The
corresponding cells are the \emph{two-sided cells}, the \emph{left cells}
and the \emph{right cells}.

\section{Admissible \texorpdfstring{$W$}{W}-graphs}
\label{sec:5}

Let $(W,S)$ be a Coxeter system, not necessarily finite. For $s,\,t\in S$,
let $m(s,t)$ be the order of $st$ in $W\!$. Thus $\{s,t\}$ is a
bond in the Coxeter diagram if and only if $m(s,t)>2$.

\begin{defi}[{\cite[Definition~2.1]{stem:addwgraph}}]\label{admissibleWg}
An $S$-coloured graph $\Gamma = (V,\mu,\tau)$ is
\emph{admissible} if the following three conditions are satisfied:
\begin{enumerate}[label=(\roman*)]
\item $\mu(V\times V) \subseteq \mathbb{N}$;
\item $\Gamma$ is symmetric, that is, $\mu(u,v) = \mu(v,u)$ if $\tau(u) \nsubseteq \tau(v)$
and $\tau(v) \nsubseteq \tau(u)$;
\item $\Gamma$ has a bipartition.
\end{enumerate}
\end{defi}

\begin{rema}\label{KLgraphisadmissibl}
As we have seen in Sec.~\ref{sec:3}, the Kazhdan--Lusztig
$W\!$-graph $\Gamma=\Gamma(W,\mu,\tau)$ is admissible.
So its cells are admissible.
\end{rema}
Let $(W,S)$ be a braid finite Coxeter system. (That is, $m(s,t)<\infty$ for all
$s,t\in S$.)

\begin{defi}[{\citex{Definition~2.1}{stem:morewgraph}}]\label{compatibility}
An $S$-coloured graph $\Gamma = (V,\mu,\tau)$ is said to satisfy the
\emph{Compatibility Rule}
if for all $u,\,v\in V$ with $\mu(u,v)\ne 0$,
each $s \in \tau(u)\setminus \tau(v)$ and each $t\in\tau(v)\setminus\tau(u)$
are joined by a bond in the Coxeter diagram of $W$.
\end{defi}
By \cite[Proposition~4.1]{stem:addwgraph}, every $W\!$-graph satisfies the
Compatibility Rule.

\begin{defi}[{\citex{Definition~2.3}{stem:morewgraph}}]\label{simplicity}
An $S$-coloured graph $\Gamma = (V,\mu,\tau)$ is said to satisfy the
\emph{Simplicity Rule}
if for all $u,\,v\in V$ with $\mu(u,v) \neq 0$,
either $\tau(v) \subsetneqq \tau(u)$ and $\mu(v,u) = 0$ or else
$\tau(u)$ and $\tau(v)$ are not comparable and $\mu(u,v) = \mu(v,u) = 1$.
\end{defi}
The Simplicity Rule implies that if $\mu(u,v) \neq 0$ and $\mu(v,u) \neq 0$
then $\mu(u,v)=\mu(v,u)=1$. That is, all edges are simple. Furthermore if
$\{u,v\}$ is an edge then $\tau(u)$ and $\tau(v)$ are not comparable,
so that there exist at least one $s\in \tau(u)\setminus\tau(v)$ and
at least one $t\in \tau(v)\setminus\tau(u)$. If the Compatibility Rule is
also satisfied, then $\{s,t\}$ must be a bond in the Coxeter
diagram.

If $(W,S)$ is simply-laced then every $W\!$-graph with non-negative
integer arc weights satisfies the Simplicity Rule, even if it fails to be
admissible: see \cite[Remark~4.3]{stem:addwgraph}.

\begin{defi}[{\citex{Definition~2.4}{stem:morewgraph}}]\label{bonding}
An admissible $S$-coloured graph $\Gamma = (V,\mu,\tau)$ satisfies the
\emph{Bonding Rule}
if for all $s,t\in S$ with $m(s,t)>2$,
the vertices $v$ of $\Gamma$ satisfying either $s\in\tau(v)$ and $t\notin\tau(v)$
or $s\notin\tau(v)$ and $t\in\tau(v)$,
together with edges of $\{u,v\}$ of $\Gamma$ such that
$s\in\tau(u)\setminus\tau(v)$ and $t\in\tau(v)\setminus\tau(u)$,
form a disjoint union of Dynkin diagrams of types $A$, $D$ or $E$
with Coxeter numbers that divide~$m(s,t)$.
\end{defi}

Equivalently, an admissible $S$-coloured graph $\Gamma$ satisfies the
Bonding Rule if and only if for all $s,\,t\in S$ with $m(s,t)>2$ the
graphs of the cells in the $W_{\{s,t\}}$-restriction of $\Gamma$ are all Dynkin
diagrams of types $A$, $D$ or $E$ with Coxeter numbers that divide~$m(s,t)$.
(Note that in the $W_{\{s,t\}}$-restriction of $\Gamma$ each edge joins
a vertex of colour $\{s\}$ to a vertex of colour~$\{t\}$.)

\begin{rema}\label{molecularbondingrule}
In the case $m(s,t)=3$, the
Bonding Rule
becomes the
\emph{Simply-Laced Bonding Rule}:
for every vertex $u$ such that
$s\in\tau(u)$ and $t\notin\tau(u)$, there exists a unique
adjacent vertex $v$ such that $t\in\tau(v)$ and $s\notin\tau(v)$.
\end{rema}
By \cite[Proposition~4.4]{stem:addwgraph}, admissible $W\!$-graphs satisfy the
Bonding Rule.

Let $\Gamma=(V,\mu,\tau)$ be an $S$-coloured graph. Let $s,t\in S$ with
$m(s,t) = m \geqslant 2$. Suppose that $u,v\in V$ with $s,t\notin \tau(u)$
and $s,t\in\tau(v)$. For $2\leqslant k\leqslant m$, a directed path
$(u,v_1,\ldots,v_{k-1},v)$ in $\Gamma$ is said to be \emph{alternating of type
$(s,t)$} if $s\in \tau(v_i)$ and $t\notin\tau(v_i)$ for odd $i$ and
$t\in\tau(v _i)$ and $s\notin\tau(v_i)$ for even $i$. Define
\begin{equation}\label{altsums}
N^k_{s,t}(\Gamma;u,v)=
\sum_{v_1,\ldots,v_{k-1}}\mu(v,v_{k-1})\mu(v_{k-1}v_{k-2})\cdots \mu(v_2,v_1)\mu(v_1,u),
\end{equation}
where the sum extends over all paths $(u,v_1,\ldots,v_{k-1},v)$ that are
alternating of type $(s,t)$.

Note that if $\Gamma$ is admissible then all terms in \eqref{altsums} are positive.

\begin{defi}[{\citex{Definition~2.9}{stem:morewgraph}}]\label{polygon}
An admissible $S$-coloured graph $\Gamma = (V,\mu,\tau)$ satisfies the
\emph{Polygon Rule}
if for all $s,t\in S$ and all $u, v\in V$ such
that $s,t\in\tau(v)\setminus \tau(u)$, we have
\begin{equation*}
N_{s,t}^{k}(\Gamma;u,v) = N_{t,s}^{k}(\Gamma;u,v) \quad \text{for all
$r$ such that $2 \leqslant k \leqslant m(s,t)$}.
\end{equation*}
\end{defi}
By \cite[Proposition~4.7]{stem:addwgraph}, all $W\!$-graphs with integer arc
weights satisfy the Polygon Rule.

The next result provides a necessary and sufficient condition for an
admissible $S$-coloured graph to be a $W\!$-graph.

\begin{theo}[{\citex{Theorem~4.9}{stem:addwgraph}}]\label{combinatorialCharacterisation}
An admissible $S$-coloured graph $\Gamma = (V,\mu,\tau)$ is a $W\!$-graph
if and only if it satisfies the
Compatibility Rule, the Simplicity Rule, the Bonding Rule and the Polygon Rule.
\end{theo}
It is convenient to introduce a weakened version of the
Polygon Rule.
\begin{defi}[{\citex{Definition~2.9}{stem:morewgraph}}]\label{localpolygon}
An admissible $S$-coloured graph $\Gamma = (V,\mu,\tau)$ satisfies the
\emph{Local Polygon Rule}
if for all $s,t\in S$, all $k$ such that
$2\leqslant k\leqslant m(s,t)$, and all $u, v$ such that
$s,t\in\tau(v)\setminus \tau(u)$, we have
$N_{s,t}^{k}(\Gamma;u,v) = N_{t,s}^{k}(\Gamma;u,v)$
under any of the following conditions:
\begin{enumerate}[label=(\roman*),topsep=1 pt]
\item $k=2$, and $\tau(u)\setminus\tau(v)\ne\emptyset$;
\item $k=3$, and there exist $r,r'\in\tau(u)\setminus\tau(v)$
(not necessarily distinct) such that
$\{r,s\}$ and $\{r',t\}$ are not bonds in the Dynkin diagram of $W$;
\item $k\geqslant 4$, and there is $r\in\tau(u)\setminus\tau(v)$ such that
$\{r,s\}$ and $\{r,t\}$ are not bonds in the Dynkin diagram of $W$.
\end{enumerate}
\end{defi}

\begin{defi}[{\citex{Definition~3.3}{stem:morewgraph}}]\label{moleculargraph}
An admissible $S$-coloured graph is called a \emph{$W\!$-molecular graph} if it satisfies the
Compatibility Rule, the Simplicity Rule, the Bonding Rule and Local Polygon Rules.
\end{defi}
A \emph{simple part} of an $S$-coloured graph $\Gamma$ is a connected
component of the graph obtained by removing all arcs and all non-simple edges,
and a \emph{simple component} of $\Gamma$ is the full subgraph spanned by a
simple part.

\begin{defi}\label{molecule}
A \emph{$W\!$-molecule} is a $W\!$-molecular graph that has only
one simple part.
\end{defi}

\begin{rema}\label{WgraphisWmoleculargraph}
If $\Gamma$ is an admissible $W\!$-graph then its simple components are
$W\!$-molecules, by~\cite[Fact~3.1]{stem:morewgraph}.
More generally, by~\cite[Fact~3.2]{stem:morewgraph}, the full subgraph of
$\Gamma$ induced by any union of simple parts is a $W\!$-molecular graph.
\end{rema}

It is easy to check that if $\Gamma=(V,\mu,\tau)$ is a $W\!$-molecular graph
then its $W_J$-restriction $\Gamma_{J}$ is a $W_J$-molecular graph.
The $W_J$-molecules of $\Gamma_J$ are called $W_J$-submolecules
of~$\Gamma$.

\begin{prop}[{\citex{Lemma~1}{ChuMic:typeAMol}}]\label{arctransport}
Let $(W,S)$ be a Coxeter system and $M=(V,\mu,\tau)$
a $W\!$-molecular graph, and let $J=\{r,s,t\}\subseteq S$ with
$m(s,t)=3$ and $r\notin\{s,t\}$. Suppose that $v,v',u,u'\in V$,
and that $\{v,v'\}$ and $\{u,u'\}$ are simple edges with
\begin{align*}
\qquad&&\tau(v)\cap J &= \{s\},&\tau(u)\cap J &= \{s,r\},&&\qquad\\[-1.5pt]
\qquad&&\tau(v')\cap J &= \{t\},&\tau(u')\cap J &= \{t,r\}.&&\qquad
\end{align*}
Then $\mu(u,v) = \mu(u',v')$.
\end{prop}

\section{Tableaux, left cells and admissible molecules of type \texorpdfstring{$A$}{A}}
\label{sec:6}

For the remainder of this paper we shall focus attention on Coxeter systems
of type~$A$. For each positive integer $n$ we write $W_n$ for the symmetric
group on the set $\{1,2,\ldots,n\}$, and let $S_n=\{s_i\mid i \in [1,n-1]\}$,
where $s_i$ is the transposition that swaps $i$ and $i+1$. Then $(W_n,S_n)$
is a Coxeter system of type $A_{n-1}$. We write $\mathcal{H}_n$ for the
Hecke algebra of~$W_n$. If $1\leqslant h\leqslant k\leqslant n$ then we write
$W_{[h,k]}$ for the standard parabolic subgroup of $W_n$ generated by
$\{\,s_i\mid i\in[h,k-1]\,\}$. We adopt a left operator convention for permutations,
writing $wi$ for the image of $i$ under the permutation~$w$. We also define
$s_{ij}\in W_n$ to be the transposition that swaps $i$~and~$j$.

A sequence of nonnegative integers $\alpha =
(\alpha_{1},\alpha_{2} \ldots, \alpha_{k})$ is called a \emph{composition}
of~$n$ if $\sum_{i=1}^k\alpha_i=n$. The $\alpha_i$ are called the \emph{parts}
of $\alpha$. We adopt the convention that $\alpha_i=0$ for all $i>k$.
A composition $\lambda = (\lambda_1,\lambda_2,\ldots,\lambda_k)$ is called a
\emph{partition} of $n$ if $\lambda_{1} \ge \cdots \ge \lambda_{k}>0$.
We define $C(n)$ and $P(n)$ to be the sets
of all compositions of~$n$ and all partitions of~$n$, respectively.

Since some of the conventions we are about to adopt are slightly non-standard, it seems
appropriate to first make the following motivational remarks.

If $\Gamma=(V,\mu,\tau)$ is the $W_n\!$-graph associated with a Kazhdan--Lusztig
left cell in~$W_n$ then the $\mathcal{H}_n$-module $M(\Gamma)$ is irreducible,
and specializing to $q=1$ yields a $W_n$-module $M(\Gamma)_1$ that is isomorphic
to a Specht module. It follows that $V$ is in bijective correspondence with
the set of standard Young tableaux of some shape. For each $v\in V$ the set
$\tau(v)\subseteq S_n$ generates a parabolic subgroup of~$W_n$, which acts on the
1-dimensional subspace of $M(\Gamma)_1$ spanned by~$v$ via the sign character.
Now since a Young tableau is conventionally associated with the $1$-character of its row
group and the sign character of its column group, it becomes more natural for
our purposes to focus on columns rather than rows. Hence we make
the following definition.

\begin{defi}\label{diagram}
For each
$\alpha=(\alpha_{1},\ldots,\alpha_{k})\in C(n)$ we define
\[[\alpha]=
\{\,(i,j)\mid1\leqslant i\leqslant\alpha_{j}\text{ and }1\leqslant
j\leqslant k\,\},
\]
and call this as the \emph{diagram} of~$\alpha$.
\end{defi}

Pictorially
$[\alpha]$ is represented by a top-justified array of boxes
with $\alpha_{j}$ boxes in the $j$-{th} column; the pair
$(i,j)\in[\alpha]$ corresponds to the $i$-{th} box in the
$j$-{th} column. So the diagram of~$\alpha = (3,4,2)$ is
\[\vcenter{\hbox{\begin{Young}
 &&\cr
 &&\cr
 &\cr
 \omit&\cr
 \end{Young}.}}
\]
Thus if $\lambda$ is a partition, the diagram $[\lambda]$ in our sense
is the transpose of the usual Young diagram of $\lambda$.

If $\lambda\in P(n)$ then $\lambda^*$ denotes the \emph{conjugate}
of $\lambda$,
defined to be the partition whose diagram is the transpose of
$[\lambda]$; that is,
$[\lambda^*]=\{(j,i) \mid (i,j) \in \lambda\}$ (which is the Young
diagram of $\lambda$).

Let $\lambda\in P(n)$. If $(i,j)\in[\lambda]$ and
$[\lambda]\setminus\{(i,j)\}$ is still the diagram of a partition,
we say that the box $(i,j)$ is \emph{$\lambda$-removable}.
Similarly, if $(i,j)\notin[\lambda]$ and
$[\lambda]\cup\{(i,j)\}$ is again the diagram of a partition,
we say that the box $(i,j)$ is \emph{$\lambda$-addable}.

If $\alpha\in C(n)$ then an $\alpha$\emph{-tableau} is a bijection
$t\colon[\alpha] \rightarrow\mathcal{T}$, where $\mathcal{T}$ is a
totally ordered set with~$n$ elements. We call $\mathcal{T}$ the
\emph{target} of~$t$. In this paper the target will always be an
interval $[m+1,m+n]$, with $m=0$ unless otherwise specified.
The composition $\alpha$ is called the
\emph{shape} of $t$, and we write $\alpha=\shape(t)$.
For each $i\in[1,n]$ we define $\row_t(i)$ and $\col_t(i)$ to be the row index
and column index of $i$~in~$t$ (so that
$t^{-1}(i)=(\row_t(i),\col_t(i)$)). We define $\Tab_m(\alpha)$
to be the set of all $\alpha$-tableaux with target
$\mathcal{T}=[m+1,m+n]$, and $\Tab(\alpha)=\Tab_0(\alpha)$.
If $h\in\mathbb{Z}$ and $t\in\Tab_m(\alpha)$ then we define
$t+h\in\Tab_{m+h}(\alpha)$ to be the tableau obtained by adding
$h$ to all entries of~$t$.

We define $\tau_\alpha\in \Tab(\alpha)$ to be the
specific $\alpha$-tableau given by
$
\tau_\alpha(i,j)=i+\sum_{h=1}^{j-1}\alpha_h
$
for all $(i,j)\in[\alpha]$. That is, in $\tau_\alpha$ the numbers
$1,\,2,\,\dots,\,\alpha_1$ fill the first column of
$[\alpha]$ in order from top to bottom, then the numbers
$\alpha_1+1,\,\alpha_1+2,\,\dots,\,\alpha_1+\alpha_2$
similarly fill the second column, and so on. If $\lambda\in P(n)$ then
we also define $\tau^{\lambda}$ to be the $\lambda$-tableau that
is the transpose of $\tau_{\lambda^*}$.
Whenever $\lambda\in P(n)$ and $t\in\Tab_m(\lambda)$ we define
$t^*\in\Tab_m(\lambda^*)$ to be the transpose of~$t$.
For example, if $\alpha=(3,2)$ then 
\[
\tau_{\alpha} = \ttab(14, 25, 3),\qquad
\tau_\alpha^*=\ttab(123, 45)\qquad
 \text{and} \qquad
 \tau^{\alpha}=\ttab(12, 34, 5).
 \]

We call $\tau_\alpha^{-1}(i)$ the $i$-th box of $[\alpha]$
in the top-to-bottom-left-to-right reading order, or TBLR~order.
(Below we shall also have occasion to make use of the BTLR order.)

Let $\alpha\in C(n)$ and $t\in\Tab(\alpha)$. We say that
$t$ is \emph{column standard} if the entries increase down
each column. That is, $t$ is column standard if $t(i,j)<t(i+1,j)$
whenever $(i,j)$ and $(i+1,j)$ are both in~$[\alpha]$. We
define $\CSTD(\alpha)$ to be the set of all column standard
$\alpha$-tableaux. In the
case $\lambda\in P(n)$ we say that $t$ is \emph{row standard} if
its transpose is column standard (so that $t(i,j) < t(i,j+1)$ whenever
$(i,j)$ and $(i,j+1)$ are both in~$[\lambda]$), and we say that $t$
is \emph{standard} if it is both row standard and column standard.
For example, if
\[
t=\vcenter{\hbox{\ttab(13, 24, 5)}},\qquad u=\vcenter{\hbox{\ttab(21, 43, 5)}}
\]
then $t$ is standard, while $u$ is column standard but not row standard.

For each $\lambda\in P(n)$ we write $\STD(\lambda)$
for the set of all standard $\lambda$-tableaux (which is the set of all
standard Young tableaux associated with $\lambda^{*}$).
We also define $\STD(n)=\bigcup_{\lambda\in P(n)}\STD(\lambda)$.

It is clear that for any fixed composition $\alpha\in C(n)$ the group
$W_{n}$ acts on $\Tab(\alpha)$, via
$(wt)(i,j) = w(t(i,j))$ for all $(i,j)\in[\alpha]$, for all
$\alpha$-tableaux $t$ and all $w \in W_{n}$. Moreover, the
map from $W_n$ to $\Tab(\alpha)$ defined by $w\mapsto
w\tau_\alpha$ for all $w\in W_n$ is bijective. We define the map
$\perm\colon\Tab(\alpha) \mapsto W_n$ to be the inverse of
$w\mapsto w\tau_\alpha$, and use this to transfer the left weak order
and the Bruhat order from $W_n$ to $\Tab(\alpha)$. Thus if
$t_1$ and $t_2$ are arbitrary $\alpha$-tableaux, we write
$t_1\leqslant\lside t_2$ if and only if $\perm(t_1)\leqslant\lside\perm(t_2)$, and
$t_1\leqslant t_2$ if and only if $\perm(t_1)\leqslant\perm(t_2)$.
Similarly, we define the length of $t\in\Tab(\alpha)$
by $l(t)=l(\perm(t))$. Since the identity element of $W_n$ is the unique
minimal element of $(W_n,{\leqslant\lside})$ and also of $(W_n,{\leqslant})$,
it follows that $\tau_\alpha$ is the unique minimal element of
$(\Tab(\alpha),{\leqslant\lside})$ and of $(\Tab(\alpha),{\leqslant})$.

\begin{rema}\label{BruhatTab}
It follows from Definition~\ref{Orders_on_W} that the Bruhat order on
$\Tab(\alpha)$ is generated by the requirement that $t\leqslant s_{ij}t$
whenever $i<j$ and $i$ precedes $j$ in $t$ the TBLR reading order.
\end{rema}

We also define a bijection $\word\colon\Tab(\alpha)\to W_n$ by using the
BTLR reading order instead of the TBLR order. Thus if $t\in\Tab(\alpha)$
and $b_1,\ldots,b_n$ is the sequence given by reading $t$ in
bottom-to-top-left-to-right order, then $\word(t)\in W_n$ is given by
$i\mapsto b_i$ for all $i\in\{1,\ldots,n\}$.

\begin{rema}\label{readingword}
It is easily seen that if $\alpha\in C(n)$ and $t\in\Tab(\alpha)$ then
$\perm(t)=\word(t)w_\alpha^{-1}$, where $w_\alpha=\word(\tau_{\alpha})$.
\end{rema}

Given $\alpha\in C(n)$ we define $J_\alpha$ to be the subset
of $S$ consisting of those $s_i$ such that $i$ and $i+1$ lie in the same column
of~$\tau_\alpha$, and $W_\alpha$ to be the standard parabolic subgroup
of $W_n$ generated by~$J_\alpha$. Note that the longest element of
$W_\alpha$ is the element $w_{\alpha}=\word(\tau_{\alpha})$ defined in
Remark~\ref{readingword} above. We write
$D_{\alpha}$ for the set of minimal length representatives of the left
cosets of $W_{\alpha}$ in $W_n$. Since $l(ds_i)>l(d)$ if and
only if $di<d(i+1)$, it follows that
$D_{\alpha}=\{\,d\in W_n \mid di<d(i+1) \text{ whenever
$s_i\in W_{\alpha}$}\,\}$, and the set of column standard
$\alpha$-tableaux is precisely
$\{\,d\tau_{\alpha} \mid d\in D_{\alpha}\,\}$.

We shall also need to work with tableaux defined on skew diagrams.

\begin{defi}\label{skewpartition}\label{skewdiagram}
If $m$ and $n$ are nonnegative inegers and $\lambda\in P(m+n)$ and $\pi\in P(m)$
are such that $[\pi]\subseteq[\lambda]$ then we define
\[
[\lambda/\pi]=[\lambda]\setminus[\pi]=
\{(i,j)\mid (i,j)\in[\lambda] \text{ and } (i,j)\notin[\pi]\}
\]
and call this a \emph{skew diagram of shape $\lambda/\pi$}. We
also write $\lambda/\pi\vdash n$ and call $\lambda/\pi$ a
\emph{skew partition of~$n$}. In the case $m=0$ we identify
$\lambda/\pi$ with $\lambda$, and say that $\lambda/\pi$ is
a \emph{normal shape}.
\end{defi}

Thus the skew diagram corresponding to $\lambda/\pi$ is the
transpose of the skew Young diagram corresponding to $\lambda/\pi$.

\begin{defi}\label{skewtableau}
A \textit{skew tableau of shape $\lambda/\pi$}, or $(\lambda/\pi)$-tableau,
where $\lambda/\pi$ is a skew partition of $n$, is a bijective map
$t\colon[\lambda/\pi] \rightarrow\mathcal{T}$, where $\mathcal{T}$ is
a totally ordered set with~$n$ elements. We write $\Tab_m(\lambda/\pi)$
for the set of all $(\lambda/\pi)$-tableaux for which the target set
$\mathcal{T}$ is the interval $[m+1,m+n]$. We shall omit the subscript
$m$ if $m=0$.
\end{defi}

Let $\lambda/\pi$ be a skew partition of $n$. We define
$\tau_{\lambda/\pi}\in \Tab(\lambda/\pi)$
by
\begin{equation}\label{toptab}
\tau_{\lambda/\pi}(i,j)=i-\pi_j+\sum_{h=1}^{j-1}(\lambda_h-\pi_h)
\end{equation}
for all $(i,j)\in[\lambda/\pi]$, and define
$\tau^{\lambda/\pi}\in \Tab(\lambda/\pi)$ to be the transpose of
$\tau_{\smash{\lambda^*/\pi^*}}$.

If $\lambda/\pi\vdash n$ and $m\in\mathbb{Z}$ then $W_{[m+1,m+n]}$ acts
naturally on $\Tab_m(\lambda/\pi)$, and as for normal shapes we define
$\perm\colon\Tab_m(\lambda/\pi)\to W_{[m+1,m+n]}$ to be the inverse
of the map $w\mapsto w\tau_{\lambda/\pi}$. We transfer the Bruhat order
and the left weak order from $W_{[m+1,m+n]}$ to $\Tab_m(\lambda/\pi)$
via the bijection $\perm$, just as for normal shapes.

All of our notation and terminology for partitions and tableaux
extends naturally to skew partitions and tableaux, and will be used
without further comment.

Let $\alpha \in C(n)$ and $t$ a column standard
$\alpha$-tableau. For each $m\in\mathbb{Z}$ we define $\Lessthan tm$
to be the tableau obtained by removing from $t$ all boxes with entries
greater than~$m$.
Thus if $\beta=\shape(\Lessthan tm)$ then $\beta\in C(m)$ and
$[\beta]=\{\,b\in[\alpha]\mid t(b)\leqslant m\,\}$, and
$\Lessthan tm\colon[\beta]\to [1,m]$ is the restriction of~$t$.
It is clear that $\Lessthan tm$ is column standard. Moreover, if
$\alpha=\lambda\in P(n)$ and $t\in\STD(\lambda)$ then
$\beta=\pi\in P(m)$ and $\Lessthan tm\in\STD(\pi)$.

Similarly, if $\lambda\in P(n)$ and $t\in\STD(\lambda)$ then
for each $m\in\mathbb{Z}$ we define $\morethan tm$ to be the skew tableau
obtained by removing from $t$ all boxes with entries
less than or equal to $m$. Observe that
$\{\,b\in [\lambda]\mid t(b)\leqslant m\,\}$ is the diagram of
a partition $\nu\in P(n)$, and $\lambda/\nu$ is
a skew partition of $n-m$. Clearly $\morethan tm$ is the restriction of~$t$
to $[\lambda/\nu]$, and $\morethan tm\in \STD_m(\lambda/\nu)$.

We also let $\lessthan tm=\Lessthan t{m-1}$ and $\Morethan tm=\morethan t{m-1}$.
For example, if
$t=\vcenter{\hbox{\ttab(12,35,4)}}\in\STD((3,2))$~then
\begin{align*}
\lessthan t5=\Lessthan t4&=\vcenter{\hbox{\ttab(12,3,4)}}\in\STD((3,1)),\\
\Morethan t2=\morethan t1&=\vcenter{\hbox{\ttab(\none 2,35,4)}}\in\STD_1((3,2)/(1)).
\end{align*}

For the next two definitions we follow the conventions of~\cite{math:heckeA}
(although we use columns rather than rows).

\begin{defi}\label{Dominance Order}
Let $n$ be a nonnegative integer and $\alpha,\,\beta \in C(n) $. We say that
$\alpha$ \emph{dominates}~$\beta$,
and write $\alpha \trianglerighteq \beta $, if $\sum_{i=1}^{k} \alpha_i \geqslant
\sum_{i=1}^{k} \beta_i$ for each positive integer~$k$.
\end{defi}

\begin{defi}\label{CSTD dominance}
Let $\alpha\in C(n)$ and $t,\,u\in\CSTD(\alpha)$. We say that
$t$ \emph{dominates}~$u$, and write $t\trianglerighteq u$, if
$\shape(\Lessthan tm)\trianglerighteq\shape(\Lessthan um)$ for all $m\in[1,n]$.
\end{defi}

In~\cite{math:heckeA} Mathas defines a partial order $\trianglerighteq$ on
$W_n$ as follows: if $v,\,w\in W_n$ then $v\trianglerighteq w$ if and
only if $v$ has a reduced expression that is a subexpression of some reduced expression
for~$w$. By \cite[Theorem~1.1]{deodhar:bruhat} we see that this order is
the reverse of the usual Bruhat order on $W$ (as defined in Definition~\ref{Orders_on_W}),
in that $v\trianglerighteq w$ if and only if $v\leqslant w$. Hence restating
\cite[Theorem~3.8]{math:heckeA} gives the following theorem.

\begin{theo}\label{equidombruhatold}
Let $\alpha\in C(n)$, and let $t$ and $u$ be column standard
$\alpha$-tableaux. Then $t\trianglerighteq u$ if and only if
$\perm(t) \leqslant \perm(u)$.
\end{theo}

By our previous definitions, this says that if $t,\,u\in\CSTD(\alpha)$
then $t\trianglerighteq u$ if and only if $t\leqslant u$. Accordingly,
we make the following definition.
\begin{defi}\label{BruhatforC(n)}
If $\alpha,\,\beta\in C(n)$ we define $\alpha\leqslant\beta$
if and only if $\alpha\trianglerighteq\beta$. We call $\leqslant$ the
\emph{Bruhat order} on $C(n)$.
\end{defi}
We obtain the following variant of Definition~\ref{CSTD dominance}.

\begin{prop}\label{equidombruhat}
Let $\alpha\in C(n)$, and let $t$ and $u$ be column standard
$\alpha$-tableaux. Then $t\leqslant u$ if and only if
$\shape(\Lessthan tm)\leqslant\shape(\Lessthan um)$ for all $m\in[1,n]$.
\end{prop}

\begin{exam*}
Let $t,\,u\in C(1,2,1,1)$ be given by
\[
t=\vcenter{\hbox{\ttab(2145,\none3)}},\qquad
u=\vcenter{\hbox{\ttab(2341,\none5)}}
\]
and consider $\Lessthan um$ and $\Lessthan tm$ and their shapes for each $m\in [1,5]$.
\begin{align*}
 \Lessthan tm&: \vcenter{\hbox{\ttab(\none1\none\none)}}\qquad\! \vcenter{\hbox{\ttab(21\none\none)}}
 \qquad\! \vcenter{\hbox{\ttab(21\none\none,\none3)}}\qquad\! \vcenter{\hbox{\ttab(214\none,\none3)}}
 \qquad\! \vcenter{\hbox{\ttab(2145,\none3)}} \\
 \shape(\Lessthan tm)&: (0,1,0,0)\qquad\! (1,1,0,0)\qquad\! (1,2,0,0)\qquad\! (1,2,1,0)\!\qquad (1,2,1,1)\\
 \Lessthan um&: \vcenter{\hbox{\ttab(\none\none\none1)}}\!\qquad \vcenter{\hbox{\ttab(2\none\none1)}}
 \qquad\! \vcenter{\hbox{\ttab(23\none1)}}\qquad \vcenter{\hbox{\ttab(2341)}}
 \qquad \!\vcenter{\hbox{\ttab(2341,\none5)}}\\
 \shape(\Lessthan um)&: (0,0,0,1)\qquad\! (1,0,0,1)\qquad\! (1,1,0,1)\qquad\! (1,1,1,1)\qquad\! (1,2,1,1).
\end{align*}
We find that $\shape(\Lessthan tm)\leqslant\shape(\Lessthan um)$
for all $m\in[1,5]$, and so $t\leqslant u$ by Proposition~\ref{equidombruhat}.
\end{exam*}

Let $\alpha,\,\beta\in C(n)$ with $\alpha\ne\beta$, and suppose that
$\alpha\geqslant\beta$. Then $\sum_{i=1}^{k} \alpha_i \leqslant\sum_{i=1}^{k} \beta_i$
for all $k\in\mathbb{N}$, and so the least $k$ such that $\alpha_k\ne\beta_k$
must satisfy $\alpha_k<\beta_k$. We make the following definition.

\begin{defi}\label{Lexicographic Order}
Let $\alpha,\,\beta \in C(n) $. We write $\alpha >_{\lex} \beta $ (or
$\beta <_{\lex} \alpha$) if there exists a positive integer~$k$ such
that $\alpha_k<\beta_k$ and $\alpha_i=\beta_i$ for all $i<k$.
We write
$\alpha \geqslant_{\lex} \beta $ if $\alpha=\beta$ or $\alpha >_{\lex} \beta $.
\end{defi}

It is clear that $\geqslant_{\lex}$ is a total order on $C(n)$, and the
remarks preceding Definition~\ref{Lexicographic Order} have established
the next proposition, which says that $\geqslant_{\lex}$ is a refinement
of $\geqslant$.

\begin{prop}\label{domimplieslex}
If $\alpha,\beta\in C(n)$ with $\alpha \geqslant \beta$ then
$\alpha\geqslant_{\lex}\beta$.
\end{prop}

We call $\geqslant_{\lex}$ the \emph{lexicographic order} on $C(n)$.
Note, however, that our lexicographic order is the reverse of the one defined
in~\cite{math:heckeA}, which is defined as a refinement of $\trianglerighteq$
rather than $\geqslant$.

\begin{rema}\label{resdomrem}
Let $\gamma\in C(n)$ and $t,\,u\in \CSTD(\gamma)$ with $t\ne u$. Since
$\Lessthan t0 = \Lessthan u0$ and $\Lessthan tn \ne \Lessthan un$, we can
choose $i\in[0,n-1]$ satisfying $\Lessthan ti = \Lessthan ui$ and
$\Lessthan t{(i+1)}\neq\Lessthan u{(i+1)}$. We shall show that if
$t>u$ in the Bruhat order then $i+1$ occurs in a later column in
$t$ than in~$u$.

Note that $\Lessthan t{(i+1)}$ and $\Lessthan u{(i+1)}$ have different shapes,
since $\Lessthan t{(i+1)}$ is obtained by adding the number $i+1$ to the bottom
of some column of $\Lessthan ti = \Lessthan ui$, and clearly $\Lessthan u{(i+1)}$
must be obtained by adding $i+1$ to the bottom of a different column. Let
$\alpha=\shape(\Lessthan t{(i+1)})$ and $\beta=\shape(\Lessthan u{(i+1)})$, and
let $k=\col_t(i+1)$ and $l=\col_u(i+1)$. Then $k\ne l$, and $\alpha_j=\beta_j$
for all $j<m=\min(k,l)$. Furthermore, $\alpha_m=\beta_m+1$ if $m=k$, and
$\beta_m=\alpha_m+1$ if $m=l$. Thus $\alpha_m<\beta_m$ if and only if $m=l$.
So $\alpha>_{\lex}\beta$ if and only if $\min(k,l)=l$.

Now suppose that $t>u$. Since $\Lessthan ti = \Lessthan ui$
we have $\shape(\Lessthan th)=\shape(\Lessthan uh)$ for all $h\leq i$,
and by Proposition~\ref{equidombruhat}, we must have
$\shape(\Lessthan t{(i+1)}) \geqslant \shape(\Lessthan u{(i+1)})$. That
is, $\alpha\geqslant\beta$. By Proposition~\ref{domimplieslex} it
follows that $\alpha\geqslant_{\lex}\beta$, and so
$\col_t(i+1)=k>l=\col_u(i+1)$.
\end{rema}

Let $\alpha=(\alpha_1,\ldots,\alpha_k)\in C(n)$. For each
$t\in\Tab(\alpha)$, we define $\cp(t)$ to be the composition of
the number $\sum_{i=1}^{k}i\alpha_i$ given by $\cp(t)_i=\col_t(n+1-i)$,
the column index of $n+1-i$ in $t$, for all $i\in [1,n]$.
Thus, for example, putting
\[
t=\vcenter{\hbox{\ttab(1954,62\none7,\none3,\none8)}},\qquad
u=\vcenter{\hbox{\ttab(1954,32\none7,\none6,\none8)}}
\]
gives $\cp(t)=(2,2,4,1,3,4,2,2,1)$ and $\cp(u)=(2,2,4,2,3,4,1,2,1)$.

\begin{nota*}
Let $\alpha\in C(n)$ and $t,\,u\in\Tab(\alpha)$. We write
$t\geqslant_{\lex}u$ if $\cp(t)\geqslant_{\lex}\cp(u)$, and we
write $t>_{\lex}u$ if $\cp(t)>_{\lex}\cp(u)$.
\end{nota*}
In the example above, since $\cp(t)_{4}<\cp(u)_{4}$ and
$\cp(t)_{i}=\cp(u)_{i}$ for $i=1,2,3$, it follows that $\cp(t) >_{\lex}\cp(u)$,
whence $t>_{\lex}u$.

Observe that $\cp(t)=\cp(u)$ if and only if each column of $t$
contains the same numbers as the corresponding column of~$u$. Thus
$\geqslant_{\lex}$ is not a partial order on $\Tab(\alpha)$ but
merely a preorder. It is, however, a total preorder, in the sense that
any two elements of $\Tab(\alpha)$ are comparable, and it becomes a total order
when restricted to~$\CSTD(\alpha)$.

\begin{defi}\label{Lexicographic Order Tableaux} Given $\alpha\in C(n)$,
we call the restriction of $\geqslant_{\lex}$ to $\CSTD(\alpha)$
the \emph{lexicographic order} on~$\CSTD(\alpha)$.
\end{defi}

\begin{rema}\label{LexTableauEquiv}
If $t,\,u\in\Tab(\alpha)$ then $t>_{\lex}u$ if and only if
there exists an integer $k\in[1,n]$ such
that $\cp(t)_k<\cp(u)_k$ and $\cp(t)_l=\cp(u)_l$ for all $l<k$.
That is, $t>_{\lex}u$ if and only if there exists $k\in[1,n]$ such
that $\col_t(n+1-k)<\col_u(n+1-k)$ and $\col_t(n+1-l)=\col_u(n+1-l)$
for all $l<k$. Writing $j=n+1-k$, this says that $t>_{\lex}u$ if
and only if there exists $j\in[1,n]$ such that $\col_t(j)<\col_u(j)$
and $\col_t(h)=\col_u(h)$ for all $h\in[j+1,n]$. Observe that if
$u$ and $t$ are column standard then the latter condition
is equivalent to $\morethan tj =\morethan uj$.
\end{rema}

\begin{lemm}\label{bruhatimplieslextableau}
Let $\alpha\in C(n)$, and let $t,\,u\in\Tab(\alpha)$. If
$t\geqslant u$ then $t\geqslant_{\lex} u$.
\end{lemm}

\begin{proof}
By Remark~\ref{BruhatTab} it suffices to show that if $1\leqslant i<j\leqslant n$
and $i$ precedes $j$ in $u$ in the TBLR order, then
$t=(i,j)u\geqslant_{\lex}u$. We may assume that $\col_u(i)\ne\col_u(j)$, since
otherwise $\cp(t)=\cp(u)$ and $t\geqslant_{\lex}u$ certainly holds.
Since $i<j$ and $i$ precedes $j$ in $u$, we see that $\col_t(j)=\col_u(i)<\col_u(j)$
and $j$ is the maximum element of $\{\,k\mid\col_t(k)\ne\col_u(k)\,\}=\{i,j\}$,
and it follows from Remark~\ref{LexTableauEquiv} that $t>_{\lex}u$, as required.
\end{proof}

\begin{coro}\label{domimplieslextableau}
For each $\alpha\in C(n)$, the lexicographic order on $\CSTD(\alpha)$
is a total order that refines the Bruhat order.
\end{coro}

\begin{defi}\label{ascentsdescentsdef}
Let $\lambda\in P(n)$. For each $t\in\STD(\lambda)$ we define
\begin{align*}
\SA(t) &= \{i\in[1,n-1] \mid \row_t(i) > \row_t(i+1)\,\},\\
\SD(t) &= \{i\in[1,n-1] \mid \col_t(i) > \col_t(i+1)\},\\
\WA(t) &= \{i\in[1,n-1] \mid \row_t(i)=\row_t(i+1)\},\\
\WD(t) &= \{i\in[1,n-1] \mid \col_t(i)=\col_t(i+1)\},
\end{align*}
and call the elements of these (respectively) the \emph{strong ascents, strong descents, weak ascents}
and \emph{weak descents} of~$t$.
We also define $\A(t)=\SA(t)\cup\WA(t)$ and
$\D(t)=\SD(t)\cup\WD(t)$.
\end{defi}
\begin{rema}\label{equalsubsetsS}
It is easily checked that
$i\in\SA(t)$ if and only if $s_it\in\STD(\lambda)$ and $s_it > t$,
while $i\in\SD(t)$ if and only if $s_it < t$ (which implies that
$s_it\in\STD(\lambda)$).
Note also that if $w=\perm(t)$ then $i\in\D(t)$ if
and only if $s_i\in \mathcal{L}(ww_{\lambda})$; this is proved
in~\cite[Lemma~5.2]{nguyen:wgideals2}.
\end{rema}

\begin{lemm}\label{dessjtminust}
Let $\lambda\in P(n)$ and $t\in\STD(\lambda)$. If $i\in\SA(t)$
then $\D(s_it)\setminus\D(t)=\{i\}$.
\end{lemm}

\begin{proof}
By Remark~\ref{equalsubsetsS} the condition $i\in\SA(t)$ implies
$i\in\SD(s_it)$. Thus $i\in\D(s_it)\setminus\D(t)$.

Since $t$ and $s_it$ differ only in the positions of $i$ and $i+1$, it
is clear that if $j<i-1$ or $j>i+1$ then $j$ and $j+1$ occupy the same
positions in $t$ and in $s_it$, and it is immediate from
Definition~\ref{ascentsdescentsdef} that $j\in\D(s_it)$ if and only
if $j\in\D(t)$. So it remains to check that if $j\in\{i-1,i+1\}$ and
$j\in\D(s_it)$ then $j\in\D(t)$. If $i-1\in\D(s_it)$ then
 $\col_{s_{i}t}(i-1)\geqslant\col_{s_{i}t}(i)$, and
$\col_t(i-1)=\col_{s_{i}t}(i-1)\geqslant\col_{s_{i}t}(i)>\col_{s_{i}t}(i+1)=\col_t(i)$.
If $\col_{s_{i}t}(i+1)\geqslant\col_{s_{i}t}(i+2)$ then
$\col_t(i+1)=\col_{s_{i}t}(i)>\col_{s_{i}t}(i+1)\geqslant\col_{s_{i}t}(i+2)=\col_t(i+2)$,
and the result follows.
\end{proof}

\begin{rema}\label{minimal-tableau}
It is clear that if $\lambda/\pi\vdash n$ and
$m\in\mathbb{Z}$ then $m+\tau_{\lambda/\pi}$ is the unique minimal element of
$\STD_m(\lambda/\pi)$ with respect to the Bruhat order and the left weak order.
Accordingly, we call $m+\tau_{\lambda/\pi}$ the \emph{minimal element
of $\STD_m(\lambda/\pi)$}. It is easily shown that if $t\in\STD_m(\lambda/\pi)$
then $t=m+\tau_{\lambda/\pi}$ if and only if $\SD(t)=\emptyset$.
That is, $t$ is minimal if and only if $\D(t)=\WD(t)$.
\end{rema}
For technical reasons it is convenient to make the following definition.

\begin{defi}\label{m-critical}
Let $\lambda/\pi \vdash n > 1$ and $m\in\mathbb{Z}$. Let $i$ be minimal
such that $\lambda_i > \pi_i$, and assume that $\lambda_{i+1}>\pi_{i+1}$.
The $m$-critical tableau of shape $\lambda/\pi$ is the tableau
$t\in\STD_{m-1}(\lambda/\pi)$ such that $\col_t(m)=i$ and $\col_t(m+1)=i+1$,
and $\morethan t {(m+1)}$ is the minimal tableau of its shape.
\end{defi}
For example, if we put
\[
u=\vcenter{\hbox{\ttab(\none\none\none57,\none\none468,\none,\none3)}},\qquad
v=\vcenter{\hbox{\ttab(\none589,\none7,4,6)}}
\]
then $u$ is the $3$-critical tableau of shape $(4,4,2,2,2)/(4,3,1)$
and $v$ is the $4$-critical tableau of shape $(4,2,1,1)/(2)$.
(Note that the first column of $u$ is empty, as is its third row.)

If $t$ is the $m$-critical tableau of shape $\lambda/\pi$ then,
with $i$ as in the definition, column~$i$ of $t$ is the first
nonempty column, the number $m$ goes at the top of column~$i$ and
$m+1$ goes at the top of column~$i+1$, after which the numbers
$m+2,\,m+3,\,\ldots,\,m+n-1$ are inserted into the remaining places,
in TBLR order. Thus
$\col_t(m+2)=i$ if and only if $\lambda_i - \pi_i > 1$.

\begin{lemm}\label{critical-tableau}
Let $\lambda\in P(n)$ and $m\in\mathbb{Z}$, and let $t\in\STD(\lambda)$
satisfy $\col_t(m+1)=\col_t(m)+1$. Then
$\Morethan tm$ is $m$-critical if and only
if the following two conditions both hold:
\begin{enumerate}[label=(\arabic*),topsep=1 pt]
\item\label{lemma6.23_1} either $\col_t(m)=\col_t(m+2)$ or $m+1\notin\SD(t)$,
\item\label{lemma6.23_2} every $j\in\D(t)$ with $j>m+1$ is in $\WD(t)$.
\end{enumerate}
\end{lemm}

\begin{proof}
Let $\shape(\Morethan tm)=\lambda/\pi$, and put $i=\col_t(m)$. Note
that since $m+1$ is in column $i+1$ of $\Morethan tm$, it
follows that $\lambda_{i+1}>\pi_{i+1}$.

Given that $\col_t(m+1)=\col_t(m)+1$, the second alternative in condition~\ref{lemma6.23_1} is
equivalent to $\col_t(m)+1\leqslant\col_t(m+2)$. Hence condition~\ref{lemma6.23_1} is
equivalent to $\col_t(m)\leqslant\col_t(m+2)$. But by Remark~\ref{minimal-tableau},
condition~\ref{lemma6.23_2} holds if and only if $\morethan{t}{(m+1)}$ is minimal, which
in turn is equivalent to
$\col_t(m+2)\leqslant\col_t(m+3)\leqslant \, \cdots \, \leqslant \col_t(n)$.
So~\ref{lemma6.23_1} and~\ref{lemma6.23_2} both hold if and only if $\morethan{t}{(m+1)}$ is minimal
and $\col_t(j)\geqslant \col_t(m)$ for all $j\geqslant m$.

Since $\col_t(m+1)=i+1$, it follows from the definition that
$\Morethan tm$ is $m$-critical if and only if $\morethan{t}{(m+1)}$
is minimal and $i=\col_t(m)$ is equal to $\min\{\,j\mid \lambda_j>\pi_j\,\}$.
But this last condition holds if and only if $m$ is in the first nonempty
column of $\Morethan tm$, and since this holds if and only if
$\col_t(j)\geqslant \col_t(m)$ for all $j\geqslant m$, the result is
established.
\end{proof}

Recall that if $w\in W_{n}$ then applying the Robinson--Schensted algorithm
to the sequence $(w1,w2,\ldots,wn)$ produces a pair $\RS(w) = (\RSP(w),\RSQ(w))$,
where $\RSP(w),\,\RSQ(w)\in \STD(\lambda)$ for some $\lambda\in P(n)$. Details
of the algorithm can be found (for example) in~\cite[Section~3.1]{sag:sym}.
The first component of $\RS(w)$ is called the \emph{insertion}
tableau and the second component is called the \emph{recording} tableau.

The next two results are well-known.

\begin{theo}[{\cite[Theorem~3.1.1]{sag:sym}}]\label{rs}
The map $\RS\colon W_n\to\bigcup_{\lambda\in P(n)}\STD(\lambda)^2$ is
bijective.
\end{theo}

\begin{theo}[{\cite[Theorem~3.6.6]{sag:sym}}]\label{rswinv}
Let $w\in W_n$. If $\RS(w) = (t,x)$ then $\RS(w^{-1}) = (x,t)$.
\end{theo}
The next lemma will be used below in the discussion of dual Knuth
equivalence classes.

\begin{lemm}[{\citex{Lemma~6.3}{nguyen:wgideals2}}]\label{rstlambda}
Let $\lambda\in P(n)$ and let $w\in W_n$. Then
$\RS(w)=(t,\tau_\lambda)$ for some $t\in\STD(\lambda)$ if and
only if $w=vw_{\lambda}$ for some $v\in W_n$ such that
$v\tau_\lambda\in\STD(\lambda)$. When these conditions hold,
$t=v\tau_\lambda$.
\end{lemm}

\begin{defi}\label{dualKnuthequivalencerelation}
The \emph{dual Knuth equivalence relation} is the equivalence relation $\approx$
on $W_n$ generated by the requirements that for all $x\in W_n$ and
$k \in [1,n-2]$,
\begin{enumerate}[label=(\arabic*),topsep=1 pt]
\item\label{defi6.27_1} $x\approx s_{k+1}x$ whenever $\mathcal{L}(x)\cap\{s_{k},s_{k+1}\}=\{s_{k}\}$
and $\mathcal{L}(s_{k+1}x)\cap\{s_k,s_{k+1}\}=\{s_{k+1}\}$,
\item\label{defi6.27_2} $x\approx s_kx$ whenever $\mathcal{L}(x)\cap\{s_{k},s_{k+1}\} = \{s_{k+1}\}$
and $\mathcal{L}(s_kx)\cap\{s_k,s_{k+1}\}=\{s_{k}\}$.
\end{enumerate}
\end{defi}
The relations~\ref{defi6.27_1} and~\ref{defi6.27_2} above are known as the dual Knuth relations of
the first kind and second kind, respectively.

\begin{rema}\label{altDualKnuth}
It is not hard to check that~\ref{defi6.27_1} and~\ref{defi6.27_2} above can be combined to give an
alternative formulation of Definition~\ref{dualKnuthequivalencerelation},
as follows: $\approx$ is the equivalence relation on $W_n$ generated by
the requirement that $x\approx sx$ for all $x\in W_n$ and $s\in S_n$
such that $x<sx$ and $\mathcal{L}(x)\nsubseteq\mathcal{L}(sx)$.
In \cite{kazlus:coxhecke} Kazhdan and Lusztig show that whenever this
holds then $x$ and $sx$ are joined by a simple
edge in the Kazhdan--Lusztig $W$-graph $\Gamma=\Gamma(W_n)$. Furthermore,
they show that the dual Knuth equivalence classes coincide with the
left cells in $\Gamma(W_n)$.
\end{rema}
The following result is well-known.

\begin{theo}[{\citex{Theorem~3.6.10}{sag:sym}}]\label{knuth}
Let $x,y \in W_n$. Then $x\approx y$ if and only if $\RSQ(x)=\RSQ(y)$.
\end{theo}
Let $\lambda\in P(n)$, and for each $t\in\STD(\lambda)$
define $C(t)=\{\,w\in W_n \mid \RSQ(w)=t\,\}$. Theorem~\ref{knuth} says
that these sets are the dual Knuth equivalence classes in~$W_n$. It follows
from Lemma~\ref{rstlambda} that
$C(\tau_\lambda)=\{\,vw_{\lambda} \mid v\tau_\lambda\in\STD(\lambda)\,\}
=\{\,\perm(t)w_{\lambda}\mid t\in\STD(\lambda)\,\}
=\{\,\word(t)\mid t\in\STD(\lambda)\,\}$.

Let $t,\,u\in\STD(\lambda)$, and suppose that
$t = s_ku$ for some $k\in[2,n-1]$. By Remark~\ref{equalsubsetsS}
above, if $x=\word(u)$ then $\mathcal{L}(x)\cap\{s_{k-1},s_k\}=\{s_{k-1}\}$
and $\mathcal{L}(s_kx)\cap\{s_{k-1},s_k\}=\{s_k\}$ if and only
if $\D(u)\cap\{k-1,k\}=\{k-1\}$ and $\D(t)\cap\{k-1,k\}=\{k\}$.
Under these circumstances we write $u\to^{*1} t$, and say that there
is a \emph{dual Knuth move of the first kind} from $u$ to~$t$.
Similarly, if $t=s_ku$ for some $k\in[1,n-2]$ such that
$\D(u)\cap\{k,k+1\}=\{k+1\}$ and $ \D(t)\cap\{k,k+1\}=\{k\}$ then we
write $u\to^{*2}t$, and say that there is a \emph{dual Knuth
move of the second kind} from $u$~to~$t$.
Since $C(\tau_\lambda)$ is a single dual Knuth equivalence class,
we obtain the following result.

\begin{prop}\label{oneclass}
Let $\lambda\in P(n)$ and $t,\,u\in\STD(\lambda)$. Then $t$
can be transformed into $u$ by a sequence of dual Knuth moves or
inverse dual Knuth moves.
\end{prop}


We call the integer $k$ above the \emph{index} of the corresponding
dual Knuth move, and denote it by $\ind(u,t)$. For convenience, we shall
abbreviate ``dual Knuth Move'' to ``DKM''.

\begin{rema}
DKMs are also defined for standard skew tableaux; the definitions
are exactly the same as for tableaux of normal shape.
If $\lambda/\pi\vdash n$ and $u,\,t\in\STD(\lambda/\pi)$ then we
write $u\approx t$ if and only if $u$ and $t$ are related by a
sequence of DKMs.
\end{rema}
\begin{defi}\label{approxm-in-W}
For each $J\subseteq S_n$ let $\approx_J$ be the
equivalence relation on $W_n$ generated by the requirement that
$x\approx_J sx$ for all $s\in J$ and $x\in W_n$
such that $x<sx$ and $\mathcal L(x)\cap J\nsubseteq\mathcal L(sx)$.
\end{defi}

\begin{rema}\label{restriction-simple-edge}
Let $J\subseteq S_n$, let $(W,S)=(W_n,S_n)$ and let
$\Gamma$ be the regular Kazhdan--Lusztig $W$-graph. By the results of
Section~\ref{sec:3} we know that a simple
edge $\{x,y\}$ of $\Gamma$ remains a simple edge of
$\Gamma_J$ provided that
$\mathcal{L}(x)\cap J\nsubseteq\mathcal{L}(y)\cap J$ and
$\mathcal{L}(y)\cap J\nsubseteq\mathcal{L}(x)\cap J$. Recall that
the simple edges of $\Gamma$ all have the form $\{x,sx\}$, where
$s\in S$ and $x<sx\in W$. Given that $x<sx$, the condition
$\mathcal{L}(sx)\cap J\nsubseteq\mathcal{L}(x)\cap J$ holds if and
only if $s\in J$, and so $\{x,sx\}$ is a simple edge of
$\Gamma_J$ if and only if $s\in J$ and
$\mathcal L(x)\cap J\nsubseteq\mathcal L(sx)$. Thus $\approx_J$
is the equivalence relation on $W$ generated by the requirement
that $x\approx_J y$ whenever $\{x,y\}$ is a simple edge
of~$\Gamma_J$.
\end{rema}

\begin{defi}\label{approxm-dfn}
Let $\lambda\in P(n)$ and $1\leqslant m\leqslant n$. Let $\approx_m$
be the equivalence relation on $\STD(\lambda)$ defined by the
requirement that $u\approx_m t$ whenever there is a DKM
of index at most~$m-1$ from $u$ to~$t$ and
$\D(u)\cap[1,m-1]\nsubseteq\D(t)$. We shall call such a DKM a
\emph{$(\leqslant m)$-DKM}.
The $\approx_m$~equivalence classes in $\STD(\lambda)$ will be called the
\emph{$(\leqslant m)$-subclasses}
of~$\STD(\lambda)$, and we shall say that $u,\,t\in\STD(\lambda)$ are
\emph{$(\leqslant m)$ dual Knuth equivalent} whenever $u\approx_m t$.
\end{defi}

\begin{rema}\label{approxm-rmk}
Assume that $\lambda\in P(n)$ and $1\leqslant m\leqslant n$, and
let $u,\,t\in\STD(\lambda)$. If $u\to^{*2}t$ and $\ind(u,t)\leqslant m-1$
then $\D(u)\cap[1,m-1]\nsubseteq\D(t)$ if and only if $\ind(u,t)\in[1,m-2]$.
Clearly this holds if and only if $\morethan um=\morethan tm$ and
$\Lessthan um\to^{*2}\Lessthan tm$. If $u\to^{*1}t$ and $\ind(u,t)\leqslant m-1$
then $\ind(u,t)\in[2,m-1]$, and $\D(u)\cap[1,m-1]\nsubseteq\D(t)$ is
automatically satisfied. Clearly this holds if and only if
$\morethan um=\morethan tm$ and $\Lessthan um\to^{*1}\Lessthan tm$.
It follows that $u\approx_m t$ if and only if $\morethan um=\morethan tm$,
since $\shape(\Lessthan um)=\shape(\Lessthan tm)$ guarantees that
$\Lessthan um$ and $\Lessthan tm$ are related by a sequence of
DKMs. So in fact $u\approx_m t$ if and only if $t=wu$
for some $w\in W_m$.
\end{rema}
It is a consequence of Definitions~\ref{approxm-in-W} and \ref{approxm-dfn}
that if $u,\,t\in\STD(\lambda)$ then $u\approx_m t$ if and only if
$\word(u)\approx_J \word(t)$, where $J=S_m$. The set of all
$(\leqslant m)$-subclasses of $\STD(\lambda)$ is in bijective correspondence
with the set $\{\,v\in\STD_m(\lambda/\pi)\mid \pi\in P(m) \text{ and }
[\pi]\subseteq[\lambda]\,\}$, and each $(\leqslant m)$-subclass of
$\STD(\lambda)$ is in bijective correspondence with
$\STD(\pi)$ for some $\pi\in P(m)$ with $[\pi]\subseteq[\lambda]$.
If $t\in\STD(\lambda)$ then the $(\leqslant m)$-subclass that contains
$t$ is denoted by $C_m(t)$ and is given by
$C_m(t) = \{\,u\in\STD(\lambda) \mid\morethan um = \morethan tm\,\}$.

In view of Remark~\ref{altDualKnuth} and Theorem~\ref{knuth}, the following
theorem follows from the results of Kazhdan and
Lusztig~\cite[\S5]{kazlus:coxhecke}.

\begin{theo}\label{klthr}
With $\Gamma$ as in Remark~\ref{restriction-simple-edge}, if $t,\,t'\in\STD(n)$
then the $W_n$-graphs $\Gamma(C(t))$ and $\Gamma(C(t'))$
are isomorphic if and only if $\shape(t)=\shape(t')$.
In particular, if $\lambda\in P(n)$ then
$\Gamma(C(t))\cong \Gamma(C(\tau_{\lambda}))$ for every $t\in\STD(\lambda)$.
\end{theo}

\begin{coro}\label{leftcelllambda}
Let $\Gamma$ be the $W_n$-graph of a Kazhdan--Lusztig left cell of~$W_n$. Then
$\Gamma$ is isomorphic to $\Gamma(C(\tau_\lambda))$ for some
$\lambda\in P(n)$.
\end{coro}
Clearly for each $\lambda\in P(n)$ the bijection $t\mapsto\word(t)$
from $\STD(\lambda)$ to $C(\tau_\lambda)$ can be used to create a
$W_n$-graph isomorphic to $\Gamma(C(\tau_\lambda))$ with $\STD(\lambda)$
as the vertex set.

\begin{nota}\label{std-tableau-W-graph}
For each $\lambda\in P(n)$ we write
$\Gamma_\lambda=\Gamma(\STD(\lambda),\mu^{(\lambda)},\tau^{(\lambda)})$ for the
$W_n$-graph just described.
\end{nota}

\begin{rema}\label{Gamma-lambda-submolecules}
Let $\lambda\in P(n)$ and let $J=S_m\subseteq S_n$. It follows from
Remark~\ref{restriction-simple-edge} and Definition~\ref{approxm-dfn}
that the $J$-submolecules of $\Gamma_\lambda$ are spanned by the
$(\leqslant m)$-subclasses of~$\STD(\lambda)$.
\end{rema}
Now let $\lambda\in P(n)$ and $1\leqslant m\leqslant n$, and put
$J=S_n\setminus S_m$. The $J$-submolecules of $\Gamma_\lambda$ can
be determined by an analysis similar to that used above. We define
$\approx^m$ to be the equivalence relation on $\STD(\lambda)$ generated
by the requirement that $u\approx^m t$ whenever there is a DKM
of index at least~$m$ from $u$ to~$t$ and
$\D(u)\cap[m,n-1]\nsubseteq\D(t)$. The $\approx^m$ equivalence
classes in $\STD(\lambda)$ will be called the
\emph{$(\geqslant m)$-subclasses} of~$\STD(\lambda)$.
If $u,\,t\in\STD(\lambda)$ then $u\approx^m t$ if and only if
$\word(u)\approx_J \word(t)$, with $J=S_n\setminus S_m$. An
equivalent condition is that $\lessthan um=\lessthan tm$
and $\Morethan um\approx \Morethan tm$. It follows that if
$t\in\STD(\lambda)$ then the $(\geqslant m)$-subclass that contains
$t$ is the set
$C^m(t) = \{\,u\in\STD(\lambda) \mid\lessthan um = \lessthan tm
\text{ and } \Morethan um\approx \Morethan tm\,\}$.
\begin{rema}\label{Gamma-lambda-submolecules+}
Let $\lambda\in P(n)$ and $m\in[1,n]$, and put $J=S_n\setminus S_m$.
By the discussion above, the $J$-submolecules of $\Gamma_\lambda$ are
spanned by the $(\geqslant m)$-subclasses of~$\STD(\lambda)$.
\end{rema}
We shall need to use some properties of the well-known ``jeu-de-taquin''
operation on skew tableaux, which we now describe.

Fix a positive integer $n$ and a target set $\mathcal{T}=[m+1,m+n]$. It
is convenient to define a \emph{partial tableau} to be a bijection $t$
from a subset of $\{\,(i,j)\mid i,\,j\in\mathbb{Z}^+\,\}$ to~$\mathcal{T}$.
We shall also assume that the domain of $t$ is always of the form
$[\kappa/\xi]\setminus\{(i,j)\}$, where $\kappa/\xi$ is a skew partition
of~$n+1$ and $(i,j)\in[\kappa/\xi]$. If $(i,j)$ is $\xi$-addable
then $t$ is a $(\kappa/\pi)$-tableau, with
$[\pi]=[\xi]\cup\{(i,j)\}$, and if $(i,j)$ is $\kappa$-removable
then $t$ is a $(\lambda/\xi)$-tableau, with
$[\lambda]=[\kappa]\setminus\{(i,j)\}$.

Now suppose that $\lambda/\pi$ is a skew partition of $n$ and
$t\in\STD(\lambda/\pi)$, and suppose also that $c=(i,j)$
is a $\pi$-removable box.
Note that $t$ may be regarded as a partial tableau, since
$[\lambda/\pi]=[\kappa/\xi]\setminus\{(i,j)\}$, where
$[\kappa]=[\lambda]$ and $[\xi]=[\pi]\setminus\{(i,j)\}$.
The (forward) \emph{jeu-de-taquin slide on $t$ into $c$} is the
process $\jdt(c,t)$ given as follows.

Start by defining $t_0=t$ and $b_0 = (i,j)$. Proceeding recursively,
suppose that $k\geqslant 0$ and that $t_{k}$
and $b_{k}$ are defined, with $t_{k}$ a partial tableau whose
domain is $[\kappa/\xi]\setminus\{b_{k}\}$. If $b_{k}$ is
$\lambda$-removable then the process terminates, we define
$t'=t_{k}$ and put $m=k$. If $b_{k}=(g,h)$ is not
$\lambda$-removable we put $x=\min(t_k(g+1,h),t_k(g,h+1))$, define
$b_{k+1}=t_{k}^{-1}(x)$, and define $t_{k+1}$ to be the partial
tableau with domain $[\kappa/\xi]\setminus\{b_{k+1}\}$ given by
\[
t_{k+1}(b)=
\begin{cases}t_{k}(b)&\text{ whenever $b$ is in the domain of $t_{k}$
and $b\ne b_{k+1}$,}\\
x&\text{ if $b=b_{k}$.}\\
\end{cases}
\]
(We say that $x$ slides from $b_{k+1}$ into $b_{k}$.)
The tableau $t'$ obtained by the above process is denoted by $\jdt^{(c)}(t)$.
The sequence of boxes $[b_0,\,b_1,\,\ldots,\,b_m]$
is called the \emph{slide path} of $\jdt(c,t)$, and we say that
$\jdt(c,t)$ \emph{vacates} the box $d=b_m$ and \emph{produces}
the tableau $\jdt^{(c)}(t)$.

Backward jeu-de-taquin slides are defined by similar rules: the slide path of a backward
slide is the reverse of the slide path of a forward slide. The backward slide $\jdt(c,t)$
is defined whenever $c$ is a $\lambda$-addable box, its slide path terminates with
a $\pi$-addable box~$d$, and we write $\jdt_{(c)}(t)$ for the tableau produced.
We have the following well-known result.
\begin{prop}\label{forwardbackward}
Let $\lambda/\pi\vdash n$, let $c$ be a $\pi$-removable box and $d$ a
$\lambda$-removable box, and define $\pi'$ and $\lambda'$ by
$[\pi']=[\pi]\setminus\{c\}$ and $[\lambda']=[\lambda]\setminus\{d\}$.
Let $t\in\STD(\lambda/\pi)$ and $t'\in\STD(\lambda'/\pi')$. Then
$\jdt(c,t)$ vacates $d$ and produces $t'$ if and only if
$\jdt(d,t')$ vacates $c$ and produces $t$.
\end{prop}

\begin{exem}\label{slidepathexample}
Suppose that $t\in\STD((3,3,2)/(2,1))$ is given by
\[
t=\vcenter{\hbox{\ttab(\none\none2,\none34,15)}}\,,
\]
and note that the box $c=(1,2)$ is $(2,1)$-removable. The jeu-de-taquin slide on $t$
into $c$ is
\[
\vcenter{\hbox{\ttab(\none\none2,\none34,15)}}\to
\vcenter{\hbox{\ttab(\none2,\none34,15)}}\to
\vcenter{\hbox{\ttab(\none24,\none3,15)}},
\]
terminating here since the $(3,3,2)$-removable box $(2,3)$ has been vacated. Thus
\[\jdt^{(c)}(t)=\vcenter{\hbox{\ttab(\none24,\none3,15)}},\]
the slide path of $\jdt(c,t)$ is $[(1,2),(1,3),(2,3)]$, and the box vacated by
$\jdt(c,t)$ is $(2,3)$.
\end{exem}

The following observation follows immediately from the definition
of a slide path.

\begin{lemm}\label{slidepath}
Let $[(i_0,j_0),(i_1,j_1),\ldots,(i_m,j_m)]$ be the slide path of a jeu-de-taquin slide,
as described above. Then $i_0\leqslant i_1\leqslant\cdots\leqslant i_m$
and $j_0\leqslant j_1\leqslant\cdots\leqslant j_m$.
\end{lemm}

It is straightforward to check the following result.

\begin{lemm}\label{jofminandmaxt}
Let $\lambda=(\lambda_1,\lambda_2,\ldots,\lambda_l)\in P(n)$ and choose $m$
so that $\lambda_1=\lambda_2=\cdots=\lambda_m>\lambda_{m+1}$. Let $\pi=(1)\in P(1)$
and put $t=\tau_{\lambda/\pi}=\morethan{(\tau_{\lambda})}1-1$.
Then the slide $\jdt((1,1),t)$ vacates the box $(\lambda_1,m)$. Similarly, if
$u=\tau^{\lambda/\pi}=\morethan{(\tau^{\lambda})}1-1$ and $k$ satisfies
$\lambda^*_1=\lambda^*_2=\cdots=\lambda^*_k>\lambda^*_{k+1}$ then the slide
$\jdt((1,1),u)$ vacates the box $(k,\lambda_1^*)$.
\end{lemm}

A sequence of boxes $\beta=(b_1,\ldots,b_l)$ is called a \emph{slide sequence}
for a standard skew tableau $t$ if there exists a sequence of standard skew tableaux
$t=t_0,\,t_1,\,\ldots,\,t_l$ such that the jeu-de-taquin slide $\jdt(b_i,t_{i-1})$
is defined for each $i\in[1,l]$, and $t_i=\jdt^{(b_i)}(t_{i-1})$. We write
$t_l=\jdt_\beta(t)$. Note that the slide sequence $\beta=(b_1,\ldots,b_l)$ can be
extended to a longer slide sequence $b_1,\ldots,b_{l+1}$ if and only if the skew tableau
$t_l$ is not of normal shape. Indeed, if $\shape(t_l)=\lambda/\pi$ and $\pi$ is not
the empty partition then $b_{l+1}$ may be chosen to be any $\pi$-removable box. If
$t_l$ is of normal shape, so that the slide sequence $\beta$ is maximal, then
Theorem~\ref{jeudetaquintableau} below shows that the tableau $t_l$ depends only
on the initial skew tableau $t$ and not on the particular choice of maximal
slide sequence~$\beta$. Accordingly, we write $t_l=\jdt(t)$ whenever
$\beta$ is maximal.

\begin{exam*}
With $t$ as in Example~\ref{slidepathexample} above, there are two
maximal slide sequences for~$t$, namely $(1,2),(2,1),(1,1)$ and $(2,1),(1,2),(1,1)$.
They produce the following sequences of standard skew tableaux, ending at the
same standard tableau of normal shape:
\begin{gather*}
\vcenter{\hbox{\ttab(\none\none2,\none34,15)}} \longrightarrow
\vcenter{\hbox{\ttab(\none24,\none3,15)}} \longrightarrow
\vcenter{\hbox{\ttab(\none24,13,5)}} \longrightarrow
\vcenter{\hbox{\ttab(124,3,5)}},\\
\vcenter{\hbox{\ttab(\none\none2,\none34,15)}} \longrightarrow
\vcenter{\hbox{\ttab(\none\none2,134,5)}} \longrightarrow
\vcenter{\hbox{\ttab(\none24,13,5)}} \longrightarrow
\vcenter{\hbox{\ttab(124,3,5)}}.
\end{gather*}
\end{exam*}
In fact, the standard tableau $\jdt(t)$ produced by applying a maximal slide sequence
to a standard skew tableau~$t$ is the Robinson--Schensted insertion tableau of~$\word(t)$.

\begin{theo}[{\citex{Theorem~3.7.7}{sag:sym}}]\label{jeudetaquintableau}
Let $\lambda/\pi$ be a skew partition of $n$ and $t\in \STD(\lambda/\pi)$.
If $\beta$ is any maximal length slide sequence for~$t$ then
$\jdt_\beta(t)=\RSP(\word(t))$.
\end{theo}
Skew tableaux $u$ and $t$ are said to be \emph{dual equivalent} if
the skew tableaux $\jdt_\beta(u)$ and $\jdt_\beta(t)$ are of the same shape
whenever $\beta$ is a slide sequence for both $u$~and~$t$. Dual
equivalent skew tableaux are necessarily of the same shape, since
$\beta$ is allowed to be the empty sequence. It is easily shown
that if $u$ and $t$ are dual equivalent then every slide sequence
for $u$ is also a slide sequence for~$t$; so dual equivalence is indeed
an equivalence relation. Theorem~\ref{dualdualKnuthequivalence} below says
that this equivalence relation coincides with dual Knuth equivalence. It is easily
shown that $\D(t)=\D(\jdt(t))$ holds for all $t\in\STD(\lambda/\pi)$;
indeed, if $i\in[1,n-1]$ and $u$ is any partial tableau used
in the construction of~$\jdt(t)$, then $\col_t(i+1)\leqslant\col_t(i)$ if and only if
$\col_u(i+1)\leqslant\col_u(i)$.

\begin{theo}[{\citex{Theorem~3.8.8}{sag:sym}}]\label{dualdualKnuthequivalence}
Let $\lambda/\pi$ be a skew partition, and let $u$ and $t$ be standard
$(\lambda/\pi)$-tableaux. Then $u$ is dual equivalent to $t$
if and only if $u\approx t$.
\end{theo}
Note that Theorem~\ref{dualdualKnuthequivalence} generalizes the fact that the set
of standard tableaux of a given normal shape form a single dual Knuth equivalence class.

If $\lambda/\pi$ is a skew partition of $n$ then the
\emph{dual equivalence graph} $\mathcal{G}_{\lambda/\pi}$
has vertex set $\STD(\lambda/\pi)$ and edge set
$\big\{\{u,t\}\mid u,t\in\STD(\lambda/\pi)\text{ and
$u\to^{*1}t$ or $u\to^{*2}t$}\big\}$. By Theorem~\ref{dualdualKnuthequivalence},
the connected components of $\mathcal{G}_{\lambda/\pi}$ correspond to the dual
equivalence classes in $\STD(\lambda/\pi)$; if $C$ is a dual equivalence class
we write $\mathcal{G}_{\lambda/\pi}(C)$ for component with vertex set~$C$.
It follows from Proposition~\ref{oneclass} that if $\pi$ is the empty partition
then $\mathcal{G}_{\lambda}=\mathcal{G}_{\lambda/\pi}$ is connected.
We call $\mathcal{G}_{\lambda}$ the \emph{standard dual equivalence graph}
corresponding to $\lambda\in P(n)$.
\begin{prop}\label{glambdapiconnectedcomps}
Let $\lambda/\pi\vdash n$ and $C\subseteq\STD(\lambda/\pi)$ a dual equivalence class.
Then there exists a $\xi\in P(n)$ such that $u\mapsto\jdt(u)$ is a bijection
$C\to\STD(\xi)$ inducing a graph isomorphism from
$\mathcal{G}_{\lambda/\pi}(C)$ to $\mathcal{G}_{\xi}$.
Furthermore, this isomorphism preserves descent sets of vertices.
\end{prop}

\begin{proof}
We use induction on the cardinality of $[\pi]$, the result being trivial if $[\pi]=\emptyset$.
Now suppose that $[\pi]\ne\emptyset$, let $c\in[\pi]$ be $\pi$-removable,
and let $C'=\{\jdt^{(c)}(t)\mid t\in C\}$. It follows from the
definition of dual equivalence all elements of $C'$ have the same shape
$\lambda'/\pi'$, where $[\pi']=[\pi]\setminus\{c\}$ and
$[\lambda']=[\lambda]\setminus\{d\}$ for some $\lambda$-removable $d\in[\lambda]$.
Furthermore, since it is also clear from the definition that if $t_1,\,t_2\in C$
then $\jdt^{(c)}(t_1)$ and $\jdt^{(c)}(t_2)$ are dual equivalent, there is
a dual equivalence class containing all elements of~$C'$.
Now choose some $t\in C$ and let $t'=\jdt^{(c)}(t)$.
Let $u'\in\STD(\lambda'/\pi')$ and put $u=\jdt_{(d)}(u')$.
If $u'$ is dual equivalent to $t'$ then $u'\approx t'$
by~\citex{Theorem~3.8.8}{sag:sym}, that is,
$u'=u_0\to^{*i_{1}}u_1\to^{*i_{2}}\cdots \to^{*i_{l-1}} u_{l-1}\to^{*i_{l}}u_l=t'$, for
some integer $l$ and some $i_1,\ldots,i_l\in\{1,2\}$.
The proof of \citex{Theorem~3.8.8}{sag:sym} now shows that
$u=\jdt_{(d)}(u_0)\to^{*j_{1}}\jdt_{(d)}(u_1)\to^{*j_{2}}\cdots \to^{*j_{l-1}}\jdt(u_{l-1})\to^{*j_l} t$
for some $j_1,\ldots,j_l\in\{1,2\}$, and so
$u$ is dual equivalent to $t$ (by \citex{Theorem~3.8.8}{sag:sym}).
Hence $u\in C$, and so $u'=\jdt^{(c)}(u)\in C'$. Thus $C'$ is a dual equivalence
class in $\STD(\lambda'/\pi')$ and $\jdt^{(c)}\colon C\mapsto C'$ is surjective.
Since $\jdt_{(d)}(\jdt^{(c)}(t))=t$ for all $t\in C$,
the map~$\jdt^{(c)}$ is also injective. Moreover, the proof
of \text{\citex{Theorem~3.8.8}{sag:sym}} shows that $u$ and $t$ in $C$ are
related by a DKM if and only if $\jdt^{(c)}(u)$ and
$\jdt^{(c)}(t)$ in $C'$are related by a DKM.
It follows that $\jdt^{(c)}$ is an edge preserving bijection from $C$
to a dual equivalence class $C'$ in $\STD(\lambda'/\pi')$.
Since the jeu-de-taquin process also preserves descent sets, as we observed
above, the result now follows simply by induction.
\end{proof}

If $k\in[1,n-2]$ then each
$v\in\STD(\lambda/\pi)$ with $\D(v)\cap\{k,k+1\}=\{k\}$ is adjacent
in the dual equivalence graph to a unique $v'$ with
$\D(v')\cap\{k,k+1\}=\{k+1\}$, and each $v$ with
$\D(v)\cap\{k,k+1\}=\{k+1\}$
is adjacent to a unique $v'$ with $\D(v')\cap\{k,k+1\}=\{k\}$.
In fact, $v'=s_kv$ if
the box $v^{-1}(k+2)$ is between the boxes $v^{-1}(k)$ and $v^{-1}(k+1)$, in the sense that
$\col_v(k)<\col_v(k+2)\leq\col_v(k+1)$ or
$\col_v(k+1)<\col_v(k+2)\leq\col_v(k)$,
while $v'=s_{k+1}v$ if
$v^{-1}(k)$ is between $v^{-1}(k+1)$ and $v^{-1}(k+2)$ (meaning that
$\col_v(k+1)\leq\col_v(k)<\col_v(k+2)$ or
$\col_v(k+2)\leq\col_v(k)<\col_v(k+1)$).

\begin{defi}\label{k-neighbourdefinition}
We call the above tableau $v'$ the \emph{$k$-neighbour} of $v$, and write
$v'=k\neb(v)$.
\end{defi}
It follows from Remark~\ref{altDualKnuth} that the standard
dual equivalence graph $\mathcal{G}_{\lambda}$ is isomorphic to
the simple part of each Kazhdan--Lusztig left cell $\Gamma(C(t))$ for
$t\in\STD(\lambda)$.
Extending earlier work of Assaf~\cite{assaf:dualequigraphs}, Chmutov showed
in~\cite{ChuMic:typeAMol} that the simple part of an admissible $W_n$-molecule
is isomorphic to a standard dual equivalence graph. (Recall that
$W\!$-molecules need not be $W\!$-graphs, since the
Polygon Rule may not be satisfied.) The following result is
the main theorem of~\cite{ChuMic:typeAMol}.

\begin{theo}[{\citex{Theorem~2}{ChuMic:typeAMol}}]\label{molIsKL}
The simple part of an admissible molecule of type $A_{n-1}$ is
isomorphic to the simple part of a Kazhdan--Lusztig left cell.
\end{theo}

\begin{rema}\label{vertexsetofWngraph}
It follows that if $M=(V,\mu,\tau)$ is a molecule then there exists $\lambda\in P(n)$ and
a bijection $t\mapsto c_t$ from $\STD(\lambda)$ to $V$ such that
the simple edges of $M$ are the pairs $\{c_u,c_t\}$ such that $u,t\in\STD(\lambda)$ and there is a
DKM from $u$ to $t$ or from $t$ to $u$, and $\tau(c_t)=\{\,s_j\mid j\in\D(t)\,\}$.
The molecule~$M$ is said to be of type~$\lambda$.

Let $M=(V,\mu,\tau)$ be an arbitrary $S_n$-coloured molecular graph, and for
each $\lambda\in P(n)$ let $m_{\lambda}$
be the number of molecules of type $\lambda$ in $M$. For each $\lambda$ such
that $m_{\lambda}\neq 0$ let $\mathcal{I}_{\lambda}$ be some indexing set of
cardinality $m_{\lambda}$. Then we can write
\begin{equation}\label{vertexsetexp}
V = \bigsqcup_{\lambda\in\Lambda}\bigsqcup_{\alpha\in\mathcal{I}_{\lambda}}V_{\alpha,\lambda},
\end{equation}
where $\Lambda=\{\,\lambda\in P(n)\mid m_{\lambda}\neq 0\,\}$, each
$V_{\alpha,\lambda}=\{\,c_{\alpha,t}\mid t\in\STD(\lambda)\,\}$ is the vertex set
of a molecule of type~$\lambda$, the simple edges of $M$ are the pairs
$\{c_{\alpha,u},c_{\beta,t}\}$ such that $\alpha=\beta\in\mathcal{I}_{\lambda}$
for some $\lambda\in\Lambda$ and $u,t\in\STD(\lambda)$ are related by a
DKM, and $\tau(c_{\alpha,t})=\{\,s_j\mid j\in\D(t)\,\}$ for all $(\alpha,t)\in\mathcal{I}_\lambda\times\STD(\lambda)$.
We call $\Lambda$ the set of \emph{molecule types} for~$M$.

Note that if $\Gamma=(V,\mu,\tau)$ is an admissible $W_n$-graph then $\Gamma$
is an $S_n$-coloured molecular graph, by Remark~\ref{WgraphisWmoleculargraph}, and
hence Eq.~\eqref{vertexsetexp} can be used to describe the vertex set of $\Gamma$.
\end{rema}

\begin{rema}\label{Gamma-lambdaIsAdmissible}
We know from Remark~\ref{KLgraphisadmissibl} and Corollary~\ref{leftcelllambda} that, for each
$\lambda\in P(n)$, the $W_n$-graph $\Gamma_{\lambda}=(\STD(\lambda),\mu^{(\lambda)},\tau^{(\lambda)})$
is admissible. Since $\{u,t\}$ is a simple edge in $\Gamma_{\lambda}$
when $u,\, t\in\STD(\lambda)$ are related by a
DKM, and $\STD(\lambda)$ is a single dual Knuth equivalence class, we see that
$\Gamma_{\lambda}$ consists of a single molecule (of type $\lambda$).
\end{rema}

\begin{rema}\label{restrictionrem}
Let $\Gamma=(V,\mu,\tau)$ be an admissible $W_n$-graph, and continue
with the notation and terminology of Remark~\ref{vertexsetofWngraph}
above. Let $m\in [1,n]$, and let $K=S_m$ and $L=S_n\setminus S_m$.

Let $\lambda\in\Lambda$ and $\alpha\in\mathcal I_\lambda$, and let $\Theta$
be the molecule of $\Gamma$ whose vertex set is $V_{\alpha,\lambda}$.
By Remark~\ref{vertexsetofWngraph} applied to $\Theta_K$, the $W_K$-restriction
of~$\Theta$ (as defined in Section~\ref{sec:3} above), we may write
\[
V_{\alpha,\lambda}=\bigsqcup_{\kappa\in\Lambda_{K,\alpha,\lambda}}
\bigsqcup_{\,\,\,\,\beta\in\mathcal{I}_{K,\alpha,\lambda,\kappa}\!\!\!\!}\!\!V_{\alpha,\lambda,\beta,\kappa},
\]
where $\Lambda_{K,\alpha,\lambda}$ is the set of all $\kappa\in P(m)$ such that
$\Theta$ contains a $K$-submolecule of type~$\kappa$, and
$\mathcal I_{K,\alpha,\lambda,\kappa}$ is an indexing set whose size is
the number of such $K$-submolecules. Each $V_{\alpha,\lambda,\beta,\kappa}$
is the vertex set of a $K$-submolecule of $\Theta$ of type~$\kappa$. Writing
$V_{\alpha,\lambda,\beta,\kappa}=\{\,c'_{\beta,u}\mid u\in\STD(\kappa)\,\}$,
we see that each $c_{\alpha,t}\in V_{\alpha,\lambda}$ coincides with some
$c'_{\beta,v}$ with $\beta\in\mathcal I_{K,\alpha,\lambda,\kappa}$ and
$v\in\STD(\kappa)$.
It follows from Remark~\ref{Gamma-lambda-submolecules} above that the
$K$-submolecule of~$\Theta$ containing a given vertex $c_{\alpha,t}$ is
spanned by the
$(\leqslant m)$-subclass~$C_m(t)=\{\,u\in\STD(\lambda)\mid\morethan um=\morethan tm\,\}$.
Thus when we write $c_{\alpha,t}=c'_{\beta,v}$ as above, we can identify
$v$ with~$\Lessthan tk$.

Similarly, applying Remark~\ref{vertexsetofWngraph} to $\Theta_L$,
we may write
\[
V_{\alpha,\lambda}=\bigsqcup_{\theta\in\Lambda_{L,\alpha,\lambda}}
\bigsqcup_{\,\,\,\,\gamma\in\mathcal{I}_{L,\alpha,\lambda,\theta}\!\!\!\!}\!\!V_{\alpha,\lambda,\gamma,\theta},
\]
where $\Lambda_{L,\alpha,\lambda}$ is the set of all $\theta\in P(n-m+1)$ such
that $\Theta$ contains an $L$-submolecule of type~$\theta$, each
$V_{\alpha,\lambda,\gamma,\theta}$ is the vertex set of an $L$-submolecule
of type~$\theta$, and the set $\mathcal I_{L,\alpha,\lambda,\theta}$ indexes
these submolecules. Writing
$V_{\alpha,\lambda,\gamma,\theta}=\{\,c''_{\gamma,v}\mid v\in\STD_{m-1}(\theta)\,\}$
(where $\STD_{m-1}(\theta)$ is the set of standard $\theta$-tableaux with target
$[m,n]$), we see that each $c_{\alpha,t}\in V_{\alpha,\lambda}$ coincides
with some $c''_{\gamma,v}$ with $\gamma\in\mathcal I_{L,\alpha,\lambda,\theta}$
and $v\in\STD_{m-1}(\theta)$.
By Remark~\ref{Gamma-lambda-submolecules+} above we see that the
$L$-submolecule of~$\Theta$ containing a given vertex $c_{\alpha,t}$ is
spanned by the $(\geqslant m)$-subclass
$C^m(t) = \{\,u\in\STD(\lambda) \mid\lessthan um = \lessthan tm
\text{ and } \Morethan um\approx \Morethan tm\,\}$.
Since the condition $\Morethan um\approx\Morethan tm$ is equivalent to
$1-m + (\Morethan um)\approx 1-m+ (\Morethan tm)$, it follows from
Proposition~\ref{glambdapiconnectedcomps} that
when we write $c_{\alpha,t}=c''_{\beta,v}$
as above we can identify $v$ with~$\jdt(\Morethan tm)=m-1+\jdt(1-m+(\Morethan tm))$.
\end{rema}

\section{Extended Bruhat order on \texorpdfstring{$\STD(n)$}{Std(n)} and paired dual Knuth equivalence relation}
\label{sec:8}

Let $n\geqslant 1$ and let $(W_n,S_n)$ be the Coxeter group of type $A_{n-1}$
Recall from Proposition~\ref{equidombruhat}
that if $\lambda\in P(n)$ and $u,\,t\in\STD(\lambda)$ then $u\leqslant t$
if and only if $\shape(\Lessthan um) \leqslant \shape(\Lessthan tm)$ for all $m\in [1,n]$.
Hence it is natural to make the following definition.

\begin{defi}\label{ExtDominance Order Tableaux}
Let $\lambda,\,\pi\in P(n)$, and let $u\in\STD(\lambda)$ and $t\in\STD(\pi)$.
We write $u \leqslant t$ if $\shape(\Lessthan um) \leqslant
\shape(\Lessthan tm)$ for all $m\in [1,n]$.
\end{defi}

It is obvious that this is a partial order on
$\STD(n)=\bigcup_{\lambda\in P(n)}\STD(\lambda)$ extending the Bruhat order on each
$\STD(\lambda)$.
\begin{exam*}
For $n=2$ and $n=3$ we obtain
\[
\ttab(1,2)<\ttab(12) \qquad\text{and}\qquad
\ttab(1,2,3) < \ttab(13,2) < \ttab(12,3) < \ttab(123).
\]
\end{exam*}

Observe that $u \leqslant t$ if and only if $\shape(u) \leqslant \shape(t)$ and
$\Lessthan u{(n-1)} \leqslant \Lessthan t{(n-1)}$.

We remark that in \cite{brini:superalg} this order was used in the context of
the representation theory of symmetric groups, while in~\cite{CasNik:domPerm}
it was used in the context of combinatorics of permutations.

\begin{lemm}\label{nsbeforent}
Let $\pi,\,\lambda\in P(n)$, let $u\in\STD(\pi)$ and
$t\in\STD(\lambda)$, and let $\sigma=\shape(\lessthan un)$ and
$\theta=\shape(\lessthan tn)$. Suppose that $\sigma\leqslant\theta$
and $\col_u(n)\leqslant\col_t(n)$. Then $\pi\leqslant\lambda$.
\end{lemm}

\begin{proof}
Let $\col_u(n)=p$ and $\col_t(n)=r$, and assume that $p\leqslant r$.
Recall that $\sigma\leqslant\theta$ is equivalent to $\sigma\trianglerighteq\theta$,
and so
$\sum_{m=1}^l \sigma_m \geqslant \sum_{m=1}^l \theta_m$ for all~$l\in[1,r-1]$.
Hence for all $l\in[1,p-1]$ we have
\begin{equation*}
\sum_{m=1}^l \pi_m =\sum_{m=1}^l \sigma_m \geqslant \sum_{m=1}^l \theta_m =\sum_{m=1}^l \lambda_m,
\end{equation*}
while for all $l\in[p,r-1]$ we have
\begin{equation*}
\sum_{m=1}^l \pi_m = (\sigma_p+1) + \sum_{\substack{m=1\\m\neq p}}^l \sigma_m
> \sum_{m=1}^l \theta_m=\sum_{m=1}^l \lambda_m,
\end{equation*}
and for all $l>r$ we have
\begin{equation*}
\sum_{m=1}^l \pi_m =(\sigma_p+1) + \sum_{\substack{m=1\\m\neq p}}^l \sigma_m
\geqslant
(\theta_r+1) + \sum_{\substack{m=1\\m\neq r}}^l \theta_m =\sum_{m=1}^l \lambda_m.
\end{equation*}
Hence $\pi\leqslant\lambda$.
\end{proof}

\begin{lemm}\label{shapetheta}
Let $\lambda\in P(n)$ and $t\in\STD(\lambda)$. Suppose that $i\in\SD(t)$,
and let $p=\col_t(i)$ and $j=\col_t(i+1)$. For all
$h\in[1,n-1]$ let $\lambda^{(h)}=\shape(\Lessthan th)$ and $\theta^{(h)}=\shape(\Lessthan{s_it}h)$.
Then
\begin{equation}\label{sumthetatoh}
\sum_{m=1}^l \theta^{(i)}_m =
\begin{cases}
\sum_{m=1}^l \lambda^{(i)}_m = \sum_{m=1}^l \lambda^{(i+1)}_m = \sum_{m=1}^l \lambda^{(i-1)}_m & \text{if $l<j$}\\
\sum_{m=1}^l \lambda^{(i)}_m + 1 = \sum_{m=1}^l \lambda^{(i+1)}_m = \sum_{m=1}^l \lambda^{(i-1)}_m + 1 & \text{if $j\leqslant l<p$}\\
\sum_{m=1}^l \lambda^{(i)}_m = \sum_{m=1}^l \lambda^{(i+1)}_m - 1= \sum_{m=1}^l \lambda^{(i-1)}_m + 1& \text{if $p<l$}
\end{cases}
\end{equation}
and
\begin{equation}\label{sumnutoh}
\sum_{m=1}^l \lambda^{(i)}_m =
\begin{cases}
\sum_{m=1}^l \theta^{(i)}_m = \sum_{m=1}^l \theta^{(i+1)}_m = \sum_{m=1}^l \theta^{(i-1)}_m & \text{if $l<j$}\\
\sum_{m=1}^l \theta^{(i)}_m - 1 = \sum_{m=1}^l \theta^{(i+1)}_m - 1= \sum_{m=1}^l \theta^{(i-1)}_m & \text{if $j\leqslant l<p$}\\
\sum_{m=1}^l \theta^{(i)}_m = \sum_{m=1}^l \theta^{(i+1)}_m - 1= \sum_{m=1}^l \theta^{(i-1)}_m + 1& \text{if $p<l$}.
\end{cases}
\end{equation}
\end{lemm}
\begin{proof}
The results given by Eq.~(\ref{sumthetatoh}) and Eq.~(\ref{sumnutoh}) are readily obtained from the following formulae
\begin{equation}\label{thetam}
\theta^{(i)}_m =
\begin{cases}
\lambda^{(i)}_m + 1 = \lambda^{(i+1)}_{m} = \lambda^{(i-1)}_m + 1 & \text{if $m=j$}\\
\lambda^{(i)}_m - 1 = \lambda^{(i+1)}_{m} - 1 = \lambda^{(i-1)}_m & \text{if $m=p$}\\
\lambda^{(i)}_m = \lambda^{(i+1)}_{m} = \lambda^{(i-1)}_m & \text{if $m\neq j,p$}
\end{cases}
\end{equation}
and
\begin{equation}\label{num}
\lambda^{(i)}_m =
\begin{cases}
\theta^{(i)}_m - 1 = \theta^{(i+1)}_{m} - 1 = \theta^{(i-1)}_m & \text{if $m=j$}\\
\theta^{(i)}_m + 1 = \theta^{(i+1)}_{m} = \theta^{(i-1)}_m + 1 & \text{if $m=p$}\\
\theta^{(i)}_m = \theta^{(i+1)}_{m} = \theta^{(i-1)}_m & \text{if $m\neq j,p$},
\end{cases}
\end{equation}
respectively.
\end{proof}

\begin{lemm}\label{characterisationofdominanceorderontab}
Let $\pi,\,\lambda\in P(n)$, let $u\in\STD(\pi)$ and
$t\in\STD(\lambda)$. Suppose that $i\in\SD(u)\cap\SD(t)$. Then
$u \leqslant t$ if and only if $s_{i}u \leqslant s_{i}t$.
\end{lemm}

\begin{proof}
Let $j=\col_t(i+1)$ and let $k=\col_u(i+1)$.
For all $h\in[1,n]$ let $\lambda^{(h)} = \shape(\Lessthan th)$, let
$\theta^{(h)} = \shape(\Lessthan{s_{i}t}h)$,
let $\pi^{(h)} = \shape(\Lessthan uh)$ and let
 $\sigma^{(h)} = \shape(\Lessthan{s_{i}u}h)$.

Suppose that $u\leqslant t$. Since $s_iu$ and $s_it$ differ from
$u$ and $t$ respectively only in the positions of $i$ and $i+1$, we have
$\pi^{(h)}=\sigma^{(h)}$ and $\lambda^{(h)}=\theta^{(h)}$ for all $h\neq i$.
But since $\pi^{(h)}\leqslant\lambda^{(h)}$ for all $h$ by our assumption,
it follows that $\sigma^{(h)}\leqslant\theta^{(h)}$ for all $h\neq i$.
Hence to show that $s_iu\leqslant s_it$ it suffices to show that
$\sigma^{(i)}\leqslant\theta^{(i)}$.
Let $l\in\mathbb{Z}^+$ be arbitrary.

\begin{Case}{1.}
Suppose that $l\geqslant k$. By Lemma~\ref{shapetheta} applied to $u$,
we have $\sum_{m=1}^l \sigma^{(i)}_m = \sum_{m=1}^l \pi^{(i-1)}_m + 1$,
by the last two formulae of Eq.(\ref{sumthetatoh}).
Since $\pi^{(i-1)}\leqslant\lambda^{(i-1)}$ gives
$\sum_{m=1}^l \pi^{(i-1)}_m \geqslant \sum_{m=1}^l \lambda^{(i-1)}_m$,
it follows that
$\sum_{m=1}^l \sigma^{(i)}_m \geqslant \sum_{m=1}^l \lambda^{(i-1)}_m + 1$.
But by Lemma~\ref{shapetheta} applied to~$t$, in each case
in Eq.(\ref{sumthetatoh}) we have
$\sum_{m=1}^l \lambda^{(i-1)}_m + 1 \geqslant \sum_{m=1}^l \theta^{(i)}_m$.
Hence $\sum_{m=1}^l \sigma^{(i)}_m \geqslant \sum_{m=1}^l \theta^{(i)}_m$.
\end{Case}

\begin{Case}{2.}
Suppose that $l<k$. By Lemma~\ref{shapetheta} applied to $u$, we have
$\sum_{m=1}^l \sigma^{(i)}_m = \sum_{m=1}^l \pi^{(i)}_m
= \sum_{m=1}^l \pi^{(i+1)}_m\!$, by the first formula of Eq.~(\ref{sumthetatoh}).
Since $\pi^{(i)}\leqslant\lambda^{(i)}$
and $\pi^{(i+1)}\leqslant\lambda^{(i+1)}$, for each $h\in\{i,i+1\}$
we obtain
$\sum_{m=1}^l \pi^{(h)}_m \geqslant \sum_{m=1}^l \lambda^{(h)}_m$,
and hence $\sum_{m=1}^l \sigma^{(i)}_m\geqslant \sum_{m=1}^l \lambda^{(h)}_m$.
By Lemma~\ref{shapetheta} applied to $t$,
in each case in~Eq.(\ref{sumthetatoh}) there exists $h\in\{i,i+1\}$
such that $\sum_{m=1}^l \lambda^{(h)}_m = \sum_{m=1}^l \theta^{(i)}_m$.
Hence $\sum_{m=1}^l \sigma^{(i)}_m \geqslant \sum_{m=1}^l \theta^{(i)}_m$.
\end{Case}

Conversely, suppose that $s_iu\leqslant s_it$. As above, it suffices to show that
$\pi^{(i)}\leqslant\lambda^{(i)}$. Let $l\in\mathbb{Z}^+$ be arbitrary.

\begin{Case}{1.}
Suppose that $l\geqslant j$. By Lemma~\ref{shapetheta} applied to $t$,
we have $\sum_{m=1}^l \lambda^{(i)}_m = \sum_{m=1}^l \theta^{(i+1)}_m - 1$,
by the last two formulae of Eq.(\ref{sumnutoh}).
Since $\sigma^{(i+1)}\leqslant\theta^{(i+1)}$ gives
$\sum_{m=1}^l \sigma^{(i+1)}_m \geqslant \sum_{m=1}^l \theta^{(i+1)}_m$,
it follows that
$\sum_{m=1}^l \sigma^{(i+1)}_m - 1\geqslant \sum_{m=1}^l \lambda^{(i)}_m$.
But by Lemma~\ref{shapetheta} applied to~$u$, in each case
in Eq.(\ref{sumnutoh}) we have
$\sum_{m=1}^l \pi^{(i)}_m\geqslant \sum_{m=1}^l \sigma^{(i+1)}-1$.
Hence $\sum_{m=1}^l \pi^{(i)}_m \geqslant \sum_{m=1}^l \lambda^{(i)}_m$.
\end{Case}

\begin{Case}{2.}
Suppose that $l<j$. By Lemma~\ref{shapetheta} applied to $t$, we have
$\sum_{m=1}^l \lambda^{(i)}_m = \sum_{m=1}^l \theta^{(i-1)}_m=\sum_{m=1}^l \theta^{(i)}_m$,
by the first formula of Eq.~(\ref{sumnutoh}). Since $\theta^{(i-1)}\geqslant\sigma^{(i-1)}$ and
$\theta^{(i)}\geqslant\sigma^{(i)}$, for each $h\in\{i-1,i\}$ we obtain
$\sum_{m=1}^l \theta^{(h)}_m\leqslant\sum_{m=1}^l \sigma^{(h)}_m$, and hence
$\sum_{m=1}^l \lambda^{(i)}_m\leqslant\sum_{m=1}^l \sigma^{(h)}_m$.
By Lemma~\ref{shapetheta} applied to $u$,
in each case in Eq.~\ref{sumnutoh}) there exists $h\in\{i-1,i\}$
such that $\sum_{m=1}^l \sigma^{(h)}_m = \sum_{m=1}^l \pi^{(i)}_m$.
Hence $\sum_{m=1}^l \lambda^{(i)}_m\leqslant\sum_{m=1}^l \pi^{(i)}_m $.\qedhere
\end{Case}
\end{proof}

\begin{defi}\label{paireddualKmove}
Let $\lambda,\,\pi\in P(n)$ and let $1\leqslant m\leqslant n$.
Let $u,\,v \in \STD(\pi)$ and $t,\,x\in\STD(\lambda)$, and let
$i\in\{1,2\}$. We say that there is a \emph{paired
$(\leqslant m)$-DKM
of the $i$-th kind from $(u,t)$ to $(v,x)$} if there
exists $k\leqslant m-1$ such that $u\to^{*i}v$ and $t\to^{*i}x$ are
$(\leqslant m)$-DKMs
of index~$k$.
When this holds we write $(u,t)\to^{*i}(v,x)$, and call $k$ the index
of the paired move.
\end{defi}
\begin{defi}
Let $\lambda,\,\pi\in P(n)$. The \emph{paired $(\leqslant m)$ dual
Knuth equivalence relation} is the equivalence relation $\approx_m$ on
$\STD(\pi)\times\STD(\lambda)$ generated by paired
$(\leqslant m)$-DKMs.
When $m=n$ we write $\approx$ for $\approx_n$, and call it
the paired dual Knuth equivalence relation.
\end{defi}
We denote by $C_m(u,t)$ the $\approx_m$ equivalence class that contains $(u,t)$.

\begin{rema}\label{mimpliesn}
It is clear that $(u,t)\approx_m (v,x)$ implies $(u,t)\approx_{m'}(v,x)$ whenever $m\leqslant m'$.
In particular, $(u,t)\approx_m (v,x)$ implies $(u,t)\approx (v,x)$.
\end{rema}

\begin{rema}\label{C_m(u,t)-rmk}
By Remark~\ref{approxm-rmk}, if $(v,x)\in C_m(u,t)$ then $(v,x)=(wu,wt)$ for some $w\in W_m$.
Thus $\morethan vm = \morethan um$ and $\morethan xm = \morethan tm$. Now suppose that
$\Lessthan um=\Lessthan tm$, and write $\xi=\shape(\Lessthan um)$. It is clear
that every $(v,x)\in C_m(u,t)$ satifies $\Lessthan vm=\Lessthan xm$. Furthermore, since
$\STD(\xi)$ is a single dual Knuth equivalence class, for every $y\in\STD(\xi)$ there is a
sequence of $(\leqslant m)$-DKMs taking $\Lessthan um$ to $y$. This same sequence of DKMs takes
$(u,t)$ to $(v,x)$, where $v$ satisfies $\Lessthan vm = y$ and $\morethan vm= \morethan um$
and $x$ satisfies $\Lessthan xm = y$ and $\morethan xm= \morethan tm$. So it follows that if
$\Lessthan um=\Lessthan tm$ then
$C_m(u,t)=\{\,(v,x)\in\STD(\pi)\times\STD(\lambda)\mid (v,x)=(wu,wt)\text{ for some $w\in W_m$}\,\}$.
\end{rema}

\begin{exam*}
Suppose that $\pi=(3,1)$ and $\lambda=(2,1,1)$. Then the set
$\STD(\pi)\times\STD(\lambda)$ has $9$ elements.
It is easily verified that there are seven paired dual Knuth equivalence classes, of
which two classes have $2$ elements and five classes have $1$ element only.
The two non-trivial classes are
\begin{align*}
&\left\{\left(\,\vcenter{\hbox{\ttab(14,2,3),\ttab(124,3)}}\right), \left(\,\vcenter{\hbox{\ttab(13,2,4),\ttab(123,4)}}\right)\right\}\\ 
\intertext{and}
&\left\{\left(\,\vcenter{\hbox{\ttab(13,2,4),\ttab(134,2)}}\right),
\left(\,\vcenter{\hbox{\ttab(12,3,4),\ttab(124,3)}}\right)\right\}.
\end{align*}
\end{exam*}

Let $\pi,\,\lambda\in P(n)$ and $(u,t),\,(v,x)\in \STD(\pi)\times\STD(\lambda)$,
and suppose that $(v,x) = (s_iu,s_it)$ for some $i\in [1,n-1]$. If
$i\in \SD(u)\cap SD(t)$ then $u\leqslant t$ if and only if $v\leqslant x$,
by Lemma~\ref{characterisationofdominanceorderontab}, and it follows by
interchanging the roles of $(u,t)$ and $(v,x)$ that the same is true
if $i\in \SA(u)\cap SA(t)$. In particular, if there is a paired
DKM from $(u,t)$ to $(v,x)$ or from $(v,x)$ to $(u,t)$ then
$u\leqslant t$ if and only if $v\leqslant x$. An obvious induction
now yields the following result.

\begin{prop}\label{bruhatorderpreserved}
Let $\pi,\,\lambda\in P(n)$. Let $(u,t),\,(v,x)\in \STD(\pi)\times\STD(\lambda)$ and
suppose that $(u,t)\approx (v,x)$. Then $u\leqslant t$ if and only if
$v\leqslant x$.
\end{prop}
Let $\pi,\,\lambda\in P(n)$ and $(u,t),\,(v,x)\in \STD(\pi)\times\STD(\lambda)$.
and suppose that $(v,x) = (s_iu,s_it)$ for some $i\in [1,n-1]$. If
$i\in \SD(u)\cap SD(t)$ then $l(v)-l(x)=(l(u)-1)-(l(t)-1) = l(u)-l(t)$, and
if $i\in \SA(u)\cap\SA(t)$ then $l(v)-l(x)=(l(u)+1)-(l(t)+1) = l(u)-l(t)$.
In particular, $l(v)-l(x)=l(u)-l(t)$ if there is a paired DKM
from $(u,t)$ to $(v,x)$ or from $(v,x)$ to $(u,t)$. It follows
that $l(x)-l(v)$ is constant for all $(v,x)\in C(u,t)$. Hence we
obtain the following result.

\begin{prop}\label{weakorderpreserved}
Let $\pi,\,\lambda\in P(n)$. Let $(u,t),\,(v,x)\in \STD(\pi)\times\STD(\lambda)$ and
suppose that $(u,t)\approx (v,x)$. Then $u\leqslant\lside v$ if and only if
$t\leqslant\lside x$.
\end{prop}
\begin{proof}
Since $(u,t)\approx (v,x)$ there exists $w\in W_m$ such that $v=wu$
and $x=wt$. Now
$u\leqslant\lside v$ if and only if $l(v)-l(u)=l(w)$, and
$t\leqslant\lside x$ if and only if $l(x)-l(t)=l(w)$,
by the definition of the left weak order. Since
$(u,t)\approx (v,x)$ implies that $l(v)-l(u)=l(x)-l(t)$, the result
follows.
\end{proof}

\begin{defi}\label{k-restrictable-ed}
Let $\pi,\,\lambda\in P(n)$ and $(u,t)\in\STD(\pi)\times\STD(\lambda)$.
If $j\in [1,n]$ and $\Lessthan uj=\Lessthan tj$ then we say that the
pair $(u,t)$ is $j$-\emph{restrictable}.
\end{defi}

\begin{rema}\label{k-restrictable-rem}
It is clear that the set
$R(u,t)=\{\,j\in[1,n]\mid (u,t)\text{ is }j\text{-restrictable}\,\}$
is always nonempty, since $1\in R(u,t)$. Moreover, $R(u,t)=[1,k]$ for
some $k\in [1,n]$.
\end{rema}

\begin{defi}\label{restrictedNumber}
Let $\pi,\,\lambda\in P(n)$ and $(u,t)\in\STD(\pi)\times\STD(\lambda)$.
We shall call the number $k$ satisfying $R(u,t)=[1,k]$
the \emph{restriction number} of the pair $(u,t)$. If $k$ is the
restriction number of $(u,t)$ then we say that $(u,t)$ is
$k$-\emph{restricted}.
\end{defi}

\begin{rema}\label{k-restrricted-rem}
With $(u,t)$ as above, the restriction number of $(u,t)$ is at least $1$ and at
most~$n$. If $k\in[1,n]$ then $(u,t)$ is $k$-restricted if and only
if it is $k$-restrictable and not $(k+1)$-restrictable. If $(u,t)$ is
$k$-restricted then $k=n$ if and only if $u=t$, and if $k<n$
then $\col_u(k+1)\ne\col_t(k+1)$ and $\row_u(k+1)\ne\row_t(k+1)$.
\end{rema}

\begin{lemm}\label{knuthweakordercase23}
Let $\pi,\,\lambda\in P(n)$, and let $u\in\STD(\pi)$ and
$t\in\STD(\lambda)$. If $n<4$ then $\D(u)=\D(t)$ implies $u=t$.
\end{lemm}
\begin{proof}
This is trivially proved by listing all the standard tableaux.
\end{proof}

\begin{defi}\label{k-reduced}
Let $\pi,\,\lambda\in P(n)$ and $(u,t)\in\STD(\pi)\times\STD(\lambda)$.
We say that the pair $(u,t)$ is \emph{favourable} if the restriction number
of~$(u,t)$ lies in $\D(u)\oplus\D(t)$, the symmetric difference of the descent
sets of $u$~and~$t$.
\end{defi}

Note that $k\in\D(u)\oplus\D(t)$ if in one of the two tableaux the column number
of $k+1$ is less than or equal to the column number of~$k$, and in the other
the row number of $k+1$ is less than or equal to the row number of~$k$.

\begin{exem}\label{resfavexample} The pair
\[(u,t) = \left(\,\vcenter{\hbox{\ttab(12,34,5)}},\,\vcenter{\hbox{\ttab(124,35)}}\right)\]
is $3$-restricted, since 1, 2 and 3 occupy the same boxes in $u$ and $t$ but
4 does not, and is not favourable since $3$ is not a descent of either $u$~or~$t$.
\end{exem}

\begin{rema}\label{k-reducedrem}
Let $\pi,\,\lambda\in P(n)$, and suppose that $(u,t)\in\STD(\pi)\times\STD(\lambda)$
is $k$-restricted. Then no element of $[1,k-1]$ can belong to $\D(u)\oplus\D(t)$,
since the fact that $\Lessthan uk=\Lessthan tk$ means that $\D(u)\cap[1,k-1]=\D(t)\cap[1,k-1]$.
So if $(u,t)$ is favourable then $k =\min(\D(u)\oplus\D(t))$, and if
$(u,t)$ is not favourable and $\D(u)\oplus\D(t)$ is nonempty then $k<\min(\D(u)\oplus\D(t))$.
\end{rema}

Let $\pi,\,\lambda\in P(n)$, and let $(u,t)\in\STD(\pi)\times\STD(\lambda)$.
Let $i$ be the restriction number of~$(u,t)$, and suppose that $i\ne n$.
Let $w=\Lessthan ui=\Lessthan ti\in\STD(\xi)$, where
$\xi=\shape(w)$, and let also $(g,p)=u^{-1}(i+1)$ and $(h,r)=t^{-1}(i+1)$,
the boxes of $u$ and $t$ that contain~$i+1$. Thus $(g,p)$ and $(h,r)$
are $\xi$-addable, and $(g,p)\ne (h,r)$ since $(u,t)$ is not
$(i+1)$-restrictable. Clearly there is at least one $\xi$-removable
box $(d,m)$ that lies between $(g,p)$ and $(h,r)$ (in the sense that
either $g>d\geqslant h$ and $p\leqslant m<r$, or $h>d\geqslant g$ and
$r\leqslant m<p$), and note that $i\in \D(u)\oplus\D(t)$ if and only if
the $\xi$-removable box $w^{-1}(i)$ is such a box.

With $(d,m)$ as above, suppose that $w'\in\STD(\xi)$ satisfies $w'(d,m)=i$.
Then there exist unique $v\in\STD(\pi)$ and $x\in\STD(\lambda)$ such that
$v$ satisfies $\Lessthan vi = w'$ and $\morethan vi= \morethan ui$ and $x$
satisfies $\Lessthan xi = w'$ and $\morethan xi= \morethan ti$. We see
that $(v,x)$ is $i$-restricted and favourable. Furthermore, it follows from
Remark~\ref{C_m(u,t)-rmk} that $(v,x)\approx_i (u,t)$.

We denote by $F(u,t)$ the set of all $(v,x)$ obtained as above as $(d,m)$
and $w'$ vary. Thus
\[
F(u,t) = \{(v,x)\in C_i(u,t) \mid v^{-1}(i)=x^{-1}(i) \text{ lies between }
u^{-1}(i+1) \text{ and } t^{-1}(i+1)\}
\]
where $i$ is the restriction number of~$(u,t)$.
Note that $(u,t)\in F(u,t)$ if and only if $(u,t)$ is favourable.

\begin{exam*}
Let $u$ and $t$ be as in Example~\ref{resfavexample} above,
so that $i=3$ and $w=\Lessthan u3=\Lessthan t3=\vcenter{\hbox{\ttab(12,3)}}$.

The shape of $w$ is $\xi=(2,1)$, and the boxes of $u$ and $t$ that contain~4
are $(g,p)=(2,2)$ and $(h,r)=(1,3)$. Observe that these are indeed both
$\xi$-addable. There are two $\xi$-removable boxes, namely $(2,1)$ and $(1,2)$;
the fact that $(u,t)$ is not favourable corresponds to the fact that $w^{-1}(3)=(2,1)$
does not lie between $(g,p)$ and $(h,r)$. Putting $(d,m)=(1,2)$, which (necessarily)
does lie between $(g,p)$ and $(h,r)$, we find that (in this small example) there
is only one standard $\xi$-tableau $w'$ satisfying $w'(d,m)=3$, namely
$w'=w^*$. Hence the set $F(u,t)$ consists of a single pair $(v,x)$:
\[
F(u,t)=\{(v,x)\}=\left\{\left(\,\vcenter{\hbox{\ttab(13,24,5)}}\,,\,
\vcenter{\hbox{\ttab(134,25)}}\right)\right\}.
\]
Since $w$ and $w'$ are the only standard $\xi$-tableaux, we see
that $C_3(u,t)=\{(u,t),(v,x)\}$.
\end{exam*}

\begin{lemm}\label{knuthmoveweakorderrelationLem}
Let $\pi,\,\lambda\in P(n)$ and let $(u,t)\in\STD(\pi)\times\STD(\lambda)$
with $u\ne t$. Let $i$ be the restriction
number of~$(u,t)$, and assume that $i\notin \D(u)\oplus\D(t)$.
Let $(v,x)\in F(u,t)$.
Then either
$\D(x)\setminus\D(v)=\D(t)\setminus\D(u)$
and
$\D(v)\setminus\D(x)=\{i\}\cup(\D(u)\setminus\D(t))$,
this alternative occurring in the case that
$\col_u(i+1)<\col_t(i+1)$, or else
$\D(x)\setminus\D(v)=\{i\}\cup(\D(t)\setminus\D(u))$
and
$\D(v)\setminus\D(x)=\D(u)\setminus\D(t)$,
this occurring in the case that
$\col_t(i+1)<\col_u(i+1)$.
\end{lemm}

\begin{proof}
The construction of $(v,x)$ is given in the discussion above. Since
$(v,x)$ and $(u,t)$ are both $i$-restricted,
$\D(v)\cap[1,i-1]=\D(x)\cap[1,i-1]$ and $\D(u)\cap[1,i-1]=\D(t)\cap[1,i-1]$.
That is,
$(\D(v)\oplus \D(x))\cap[1,i-1] = (\D(u)\oplus \D(t))\cap[1,i-1] =\emptyset$.
Furthermore, since $\morethan vi= \morethan ui$ and
$\morethan xi= \morethan ti$ it follows that
$(\D(v)\setminus\D(x))\cap[i+1,n-1]=(\D(u)\setminus\D(t))\cap[i+1,n-1]$ and
$(\D(x)\setminus\D(v))\cap[i+1,n-1]=(\D(t)\setminus\D(u))\cap[i+1,n-1]$.
It remains to observe that if
$p=\col_v(i+1)\leqslant m=\col_v(i)=\col_x(i)<r=\col_x(i+1)$ then
$i\in\D(v)\setminus\D(x)$, while if
$r\leqslant m<p$ then $i\in\D(x)\setminus\D(v)$.
\end{proof}

\begin{lemm}\label{knuthmoveweakorderrelationLemb}
Let $\pi,\,\lambda\in P(n)$ and let $(u,t)\in\STD(\pi)\times\STD(\lambda)$.
Assume that the restriction number of~$(u,t)$ lies in
$\D(u)\oplus\D(t)$, and let $(v,x)\in F(u,t)$.
Then $\D(v)\setminus\D(x)=\D(u)\setminus\D(t)$ and
$\D(x)\setminus\D(v)=\D(t)\setminus\D(u)$.
\end{lemm}

\begin{proof}
The proof is the same as the proof of Lemma~\ref{knuthmoveweakorderrelationLem}, except
that it can be seen now that
$i\in\D(v)\setminus\D(x)$ if $i\in\D(u)\setminus\D(t)$ and
$i\in\D(x)\setminus\D(v)$ if $i\in\D(t)\setminus\D(u)$.
\end{proof}

\begin{lemm}\label{noname2}
Let $\pi,\,\lambda\in P(n)$ and $(u,t)\in\STD(\pi)\times\STD(\lambda)$,
and $i$ the restriction number of~$(u,t)$.
Suppose that $\D(t)\subsetneqq\D(u)$ and $i<j$, where
$j=\min (\D(u)\setminus\D(t))$. Let $(v,x)\in F(u,t)$.
If $\col_u(i+1)<\col_t(i+1)$ then
$\D(v)\setminus\D(x)=\{i\}\cup(\D(u)\setminus\D(t))$ and
$\D(x)\setminus\D(v)=\emptyset$, while if $\col_t(i+1)<\col_u(i+1)$ then
$\D(v)\setminus\D(x)=\D(u)\setminus\D(t)$ and
$\D(x)\setminus\D(v)=\{i\}$. In the former case $\D(v)\cap\{i,j\}=\{i,j\}$ and
$\D(x)\cap\{i,j\}=\emptyset$, while in the latter case $\D(v)\cap\{i,j\}=\{j\}$ and
$\D(x)\cap\{i,j\}=\{i\}$.
\end{lemm}

\begin{proof}
Since $\D(t)\subsetneqq\D(u)$ we have $\D(u)\oplus\D(t)=\D(u)\setminus\D(t)\neq\emptyset$.
So $j=\min(\D(u)\oplus\D(t))$, and since $j>i$ we have $i\notin\D(u)\oplus\D(t)$. Hence
$(v,x)\in F(u,t)$ satisfies
the extra properties specified in Lemma~\ref{knuthmoveweakorderrelationLem}.

If $\col_u(i+1)<\col_t(i+1)$ then Lemma~\ref{knuthmoveweakorderrelationLem} gives
$\D(v)\setminus\D(x)=\{i\}\cup(\D(u)\setminus\D(t))$ and $\D(x)\setminus\D(v)=\emptyset$,
since $\D(t)\setminus\D(u)=\emptyset$ by hypothesis. In particular, since $j\in\D(u)\setminus\D(t)$,
we see that $\D(v)\setminus\D(x)$ contains both $i$~and~$j$.

If $\col_t(i+1)<\col_u(i+1)$ then Lemma~\ref{knuthmoveweakorderrelationLem} combined together
with $\D(t)\setminus\D(u)=\emptyset$ gives
$\D(v)\setminus\D(x)=\D(u)\setminus\D(t)$ and $\D(x)\setminus\D(v)=\{i\}$.
In particular it follows that $j\in\D(v)\setminus\D(x)$ and $i\in\D(x)\setminus\D(v)$.
\end{proof}

Let $\Gamma=\Gamma(C,\mu,\tau)$ be a $W_n$-molecular graph, and let $\Lambda$
be the set of molecule types for~$\Gamma$. For each $\lambda\in\Lambda$ let
$m_{\lambda}$ be the number of molecules of type $\lambda$ in $\Gamma$, and
$\mathcal{I}_{\lambda}$ some indexing set of cardinality $m_{\lambda}$.
As in Remark~\ref{vertexsetofWngraph}, the
vertex set of $\Gamma$ can be expressed in the form
\begin{equation*}
C=\bigsqcup_{\lambda\in\Lambda}\bigsqcup_{\alpha\in\mathcal{I}_{\lambda}}C_{\alpha,\lambda},
\end{equation*}
where $C_{\alpha,\lambda}=\{c_{\alpha,t}\mid t\in\STD(\lambda)\}$ for each
$\alpha\in\mathcal{I}_{\lambda}$,
and the simple edges of $\Gamma$ are the pairs
$\{c_{\beta,u},c_{\alpha,t}\}$ such that $\alpha=\beta\in\mathcal{I}_{\lambda}$ for some $\lambda\in\Lambda$
and $u,\,t\in\STD(\lambda)$ are related by a DKM.
Furthermore, $\tau(c_{\alpha,t})=\{\,s_j\in S_n\mid j\in\D(t)\,\}$, whenever $\lambda\in\Lambda$ and
$(\alpha,t)\in\mathcal{I}_\lambda\times\STD(\lambda)$.

Now let $\lambda,\,\pi\in\Lambda$, and let $(\alpha,t)\in\mathcal{I}_\lambda\times\STD(\lambda)$ and
$(\beta,u)\in\mathcal{I}_\pi\times\STD(\pi)$, so that $c_{\alpha,t}$ and
$c_{\beta,u}$ are vertices of~$\Gamma$. Suppose that
$\D(u)\setminus\D(t)\neq\emptyset$, and let $j\in\D(u)\setminus\D(t)$.

Suppose that there exist $i<j$ and
$(v,x)\in\STD(\pi)\times\STD(\lambda)$ such that $(u,t)$ and $(v,x)$
are related by a paired $(\leqslant i)$-DKM.
Then $j\in\D(\morethan ui)\setminus\D(\morethan ti)$, since $j\in\D(u)\setminus\D(t)$
and $j>i$. Thus $j\in\D(\morethan vi)\setminus\D(\morethan xi)$, since
$(v,x)\approx_i(u,t)$ gives $\morethan vi=\morethan ui$ and $\morethan xi=\morethan ti$.
Hence $j\in\D(v)\setminus\D(x)$. Moreover, since $(u,t)$ and $(v,x)$ are related by a
paired $(\leqslant i)$-DKM,
there are $k,l\leqslant i-1$ with $|k-l|=1$ such that
\begin{align*}
\qquad&&\D(x)\cap \{k,l,j\} &= \{k\},&\D(v)\cap \{k,l,j\} &= \{k,j\},&&\qquad\\
\qquad&&\D(t)\cap \{k,l,j\} &= \{l\},&\D(u)\cap \{k,l,j\} &= \{l,j\},&&\qquad
\end{align*}
and it follows from Proposition~\ref{arctransport}
that $\mu(c_{\beta,v},c_{\alpha,x})=\mu(c_{\beta,u},c_{\alpha,t})$.

More generally, suppose that $i<j$ and $(v,x)\in\STD(\pi)\times\STD(\lambda)$ satisfy
$(v,x)\approx_i (u,t)$, so that for some $m\in\mathbb{N}$ there exist
$(u_0,t_0)$, $(u_1,t_1)$, \dots, $(u_m,t_m)$ in $\STD(\pi)\times\STD(\lambda)$,
with $(u_{h-1},t_{h-1})$ and $(u_h,t_h)$ related by a paired
$(\leqslant i)$-DKM
for each $h\in[1,m]$, and $(u_0,t_0)=(u,t)$
and $(u_m,t_m)=(v,x)$. Applying the argument in the preceding paragraph and a trivial induction,
we deduce that $j\in\D(u_h)\setminus\D(t_h)$ and
$\mu(c_{\beta,u_h},c_{\alpha,t_h})=\mu(c_{\beta,u},c_{\alpha,t})$ for all $h\in[0,m]$.
Thus we obtain the following result.

\begin{lemm}\label{noname0}
Let $\Gamma$ be a $W_n$-molecular graph. Using the notation as above, let $\lambda,\,\pi\in\Lambda$,
and let $(\alpha,t)\in\mathcal{I}_\lambda\times\STD(\lambda)$ and
$(\beta,u)\in\mathcal{I}_\pi\times\STD(\pi)$. Suppose that $\D(u)\setminus\D(t)\neq\emptyset$, and
let $j\in\D(u)\setminus\D(t)$. Then for all $i<j$ and all $(v,x)\in C_i(u,t)$ we have
$j\in\D(v)\setminus\D(x)$ and $\mu(c_{\beta,v},c_{\alpha,x})=\mu(c_{\beta,u},c_{\alpha,t})$.
\end{lemm}

\begin{coro}\label{arweightone}
Continuing with the same notation,
let $\lambda\in\Lambda$ and $u,\,t\in\STD(\lambda)$, and suppose that
$u=s_jt>t$ for some $j\in[1,n-1]$. Then $\mu(c_{\alpha,u},c_{\alpha,t}) = 1$
for all $\alpha\in\mathcal{I}_{\lambda}$.
\end{coro}

\begin{proof}
Since $t<s_jt=u$, it follows from Remark~\ref{altDualKnuth} that if $\D(t)\nsubseteq\D(u)$
then there is a DKM
from $t$ to $u$, and $\{c_{\alpha,u},c_{\alpha,t}\}$
is a simple edge. Thus $\mu(c_{\alpha,u},c_{\alpha,t}) = 1$ in this case, and so we may
assume that $\D(t)\subsetneqq\D(u)$.

Since $u$ is obtained from $t$ by interchanging $j$ and $j+1$, it is clear
that $j-1$ is the restriction number of $(u,t)$, and since $\D(u)\setminus\D(t)=\{j\}$,
by Lemma~\ref{dessjtminust}, we see that the hypotheses of Lemma~\ref{noname2} are
satisfied with $i=j-1$. Since $\col_t(i+1)<\col_u(i+1)$, it
follows that $(u,t)\approx_i(v,x)$, for some $(v,x)\in F(u,t)$ satisfying
$\D(v)\cap\{i,j\}=\{j\}$ and $\D(x)\cap\{i,j\}=\{i\}$. Since
$(u,t)\approx_i(v,x)$ there exists $w\in W_i$ with $v=wu$ and $x=wt$, and
$s_jw=ws_j$ since$j>i$. Thus $s_jx=s_jwt=ws_jt=wu=v$. Furthermore $s_jx>x$, since
$j\notin\D(x)$, and $\D(x)\nsubseteq\D(v)$ since $i\in\D(x)\setminus\D(v)$.
So there is a DKM
indexed by $j$ from $x$ to $v$, and so
$\{c_{\alpha,v},c_{\alpha,x}\}$ is a simple edge. Thus
$\mu(c_{\alpha,v},c_{\alpha,x})=1$, and so $\mu(c_{\alpha,u},c_{\alpha,t})=1$ by
Lemma~\ref{noname0}.
\end{proof}

\begin{lemm}\label{noname1muequallambda}
Continue with the notation used in Lemma~\ref{noname0} above. Let
$\pi,\,\lambda\in\Lambda$ and let $(u,t)\in\STD(\pi)\times\STD(\lambda)$.
Suppose that $\D(u) = \{n-1\} \cup \D(t)$ and that $\mu(c_{\beta,u},c_{\alpha,t}) \neq 0$
for some $\beta\in\mathcal{I}_{\pi}$ and $\alpha\in\mathcal{I}_{\lambda}$.
Suppose also that the restriction number of~$(u,t)$
is $i< n-2$. Then $\col_u(i+1)<\col_t(i+1)$, and $(u,t)\approx_i (v,x)$ for some
$(v,x)\in\STD(\pi)\times\STD(\lambda)$ such that
$\D(v)=\D(x)\cup\{i,n-1\}$ and $\mu(c_{\beta,v},c_{\alpha,x})=\mu(c_{\beta,u},c_{\alpha,t}) \neq 0$.
\end{lemm}

\begin{proof}
Since $u\ne t$, the set $F(u,t)$ is defined and nonempty.
Let $(v,x)\in F(u,t)$. Then it follows by Lemmas~\ref{noname2} and
\ref{noname0} that $(u,t)\approx_i (v,x)$
and $\mu(c_{\beta,v},c_{\alpha,x})=\mu(c_{\beta,u},c_{\alpha,t}) \neq 0$.
Moreover, if $\col_u(i+1)<\col_t(i+1)$ then
Lemma~\ref{noname2} gives $\D(v)=\D(x)\cup\{i,n-1\}$. Thus it
remains to show that $\col_u(i+1)<\col_t(i+1)$. Suppose otherwise.
Then Lemma~\ref{noname2} shows that $n-1\in\D(v)\setminus\D(x)$
and $i\in\D(x)\setminus\D(v)$, and now the
Compatibility Rule
says that
$i$ and $n-1$ must be joined by a bond in the Coxeter diagram of $W_n$.
But this contradicts $i<n-2$.
\end{proof}

\begin{lemm}\label{nminus1order0}
Suppose that $u,\,t\in\STD(n)$ are such that the restriction number of~$(u,t)$
is $n-1$ and $\D(u) = \{n-1\} \cup \D(t)$. Then
$\col_u(n)<\col_t(n)$. Thus $\shape(u)<\shape(t)$ and $u<t$.
\end{lemm}

\begin{proof}
Clearly $n\geqslant 2$. Since $\Lessthan u{(n-1)}=\Lessthan t{(n-1)}$
we have $\shape(\lessthan un)=\shape(\lessthan tn)$, and
since $n-1\in\D(u)\setminus\D(t)$ we have
$\col_u(n)\leqslant\col_u(n-1)=\col_t(n-1)<\col_t(n)$.
Hence $\shape(u)<\shape(t)$ by Lemma~\ref{nsbeforent}, and $u<t$
by Definition~\ref{ExtDominance Order Tableaux}.
\end{proof}

\begin{lemm}\label{nminus1order}
Continue with the notation used in Lemma~\ref{noname0} above. Let
$\pi,\,\lambda\in\Lambda$, let $u\in\STD(\pi)$, and let $t\in\STD(\lambda)$.
Suppose that $\D(u) = \{n-1\} \cup \D(t)$ and that $\mu(c_{\beta,u},c_{\alpha,t}) \neq 0$
for some $\beta\in\mathcal{I}_{\pi}$ and $\alpha\in\mathcal{I}_{\lambda}$,
and suppose that the restriction number of $(u,t)$ is $n-2$.
Then $(u,t)\approx_{n-2} (v,x)$ for some
$(v,x)\in\STD(\pi)\times\STD(\lambda)$ with
$\mu(c_{\beta,v},c_{\alpha,x})=\mu(c_{\beta,u},c_{\alpha,t})$, and
either $u<t$ and $\D(v) = \{n-2,n-1\} \cup \D(x)$
(in the case $\col_u(n-1)<\col_t(n-1)$), or else $(\lambda,\alpha)=(\pi,\beta)$
and $u=s_{n-1}t>t$, and $\mu(c_{\beta,u},c_{\alpha,t}) = 1$ (in the
case $\col_t(n-1)<\col_u(n-1)$).
\end{lemm}

\begin{proof}
Clearly $n\geqslant 3$. We have $\Lessthan u{(n-2)} = \Lessthan t{(n-2)}$
and $\col_u(n-1) \neq \col_t(n-1)$, since $(u,t)$ is $(n-2)$-restricted.
Let $(v,x)\in F(u,t)$ be arbitrary, and note that the hypotheses
of Lemma~\ref{noname2} hold with $i=n-2$ and $j=n-1$. Moreover, since
$(u,t)\approx_{n-2}(v,x)$, we may apply Lemma~\ref{noname0} with $i=n-2$ and
$j=n-1$, and deduce that
$\mu(c_{\beta,v},c_{\alpha,x})=\mu(c_{\beta,u},c_{\alpha,t}) \neq 0$.

\begin{Case}{1.}Suppose that $\col_u(n-1)<\col_t(n-1)$. Since
$\shape(\Lessthan u{(n-2)})=\shape(\Lessthan t{(n-2)})$
it follows from Lemma~\ref{nsbeforent} that
$\shape(\Lessthan u{(n-1)})<\shape(\Lessthan t{(n-1)})$. Moreover, since
$n-1\in\D(u)\setminus\D(t)$, it follows that
$\col_u(n)\leqslant\col_u(n-1)<\col_t(n-1)<\col_t(n)$. Hence $\pi<\lambda$
by Lemma~\ref{nsbeforent}, and $u<t$ by Definition~\ref{ExtDominance Order Tableaux}.
Moreover, since $\col_u(n-1)<\col_t(n-1)$ and $\D(u)\setminus\D(t)=\{n-1\}$, it
follows from Lemma~\ref{knuthmoveweakorderrelationLem} that
$\D(v)=\D(x)\cup\{n-2,n-1\}$.
\end{Case}
\begin{Case}{2.}Suppose that $\col_t(n-1)<\col_u(n-1)$. Lemma~\ref{knuthmoveweakorderrelationLem}
gives $\D(x)\cap\{n-2,n-1\}=\{n-2\}$ and $\D(v)\cap\{n-2,n-1\}=\{n-1\}$, and since
$\mu(c_{\beta,v},c_{\alpha,x})\ne0$ it follows from
Simplicity Rule
that $\mu(c_{\beta,v},c_{\alpha,x})=1$.
Hence $\mu(c_{\beta,u},c_{\alpha,t})=1$. Moreover, since $\{c_{\beta,v},c_{\alpha,x}\}$ is a
simple edge, it follows that from Theorem~\ref{molIsKL} and Remark~\ref{vertexsetofWngraph}
that $\lambda=\pi$ and $\alpha=\beta$. Hence $u=s_{n-1}t$, since
$\Lessthan u{(n-2)} = \Lessthan t{(n-2)}$, and $u>t$ since $\col_t(n-1)<\col_u(n-1)$.
\qedhere
\end{Case}
\end{proof}

\begin{rema}\label{nminus1iplusonest}
Continue with the notation used in Lemma~\ref{noname0} above. Let
$\pi,\,\lambda\in\Lambda$, let $u\in\STD(\pi)$, and let $t\in\STD(\lambda)$.
Suppose that $\D(u) = \{n-1\} \cup \D(t)$ and that $\mu(c_{\beta,u},c_{\alpha,t}) \neq 0$
for some $\beta\in\mathcal{I}_{\pi}$ and $\alpha\in\mathcal{I}_{\lambda}$. Let
$i$ be the restriction number of $(u,t)$, and note that $i\leqslant n-1$. If $i<n-2$ then
$\col_u(i+1)<\col_t(i+1)$ by Lemma~\ref{noname1muequallambda}, and if $i=n-1$ then
$\col_u(n)<\col_t(n)$ by Lemma~\ref{nminus1order0}. In the remaining case, namely
$i=n-2$, if $\col_u(n-1)>\col_t(n-1)$ then $u=s_{n-1}t>t$ by
Lemma~\ref{nminus1order}. Thus $\col_u(i+1)<\col_t(i+1)$
unless $i=n-2$ and $u=s_{n-1}t>t$.
\end{rema}

\begin{rema}\label{iplusonest}
Continue with the notation used in Lemma~\ref{noname0} above.
Suppose that $\pi,\,\lambda\in\Lambda$, and
let $(\beta,u)\in\mathcal{I}_{\pi}\times\STD(\pi)$ and
$(\alpha,t)\in\mathcal{I}_{\lambda}\times\STD(\lambda)$ satisfy
$\mu(c_{\beta,u},c_{\alpha,t}) \neq 0$ and $\D(t) \subsetneqq \D(u)$.
Let $j=\min(\D(u)\setminus\D(t))$, and~$i$ the restriction number of~$(u,t)$.
Note that $i\leqslant j$.

Let $K=\{s_1,s_2,\ldots,s_j\}\subseteq S_n$, and
$\Gamma_K$ the $W_K$-restriction of~$\Gamma$.
As in Remark~\ref{restrictionrem}, for each $\lambda\in\Lambda$
and $\alpha\in\mathcal{I}_\lambda$ we define $\Lambda_{K,\alpha,\lambda}$ to be the
set of all $\kappa\in P(j+1)$ such that the molecule of $\Gamma$ with the vertex set
$C_{\alpha,\lambda}$ contains a $K$-submolecule of type~$\kappa$,
and let $\mathcal{I}_{K,\alpha,\lambda,\kappa}$ index these
submolecules. Let $\Lambda_K=\bigcup_{\alpha,\lambda}\Lambda_{K,\alpha,\lambda}$, the
set of molecule types for $\Gamma_K$, and for each $\kappa\in\Lambda_K$ let
$\mathcal{I}_{K,\kappa}=\bigsqcup_{\{(\alpha,\lambda)|\kappa\in\Lambda_{K,\alpha,\lambda}\}}
\mathcal{I}_{K,\alpha,\lambda,\kappa}$. For each $\beta\in\mathcal{I}_{K,\kappa}$ we
write $\{c'_{\beta,u}\mid u\in\STD(\kappa)\}$ for the vertex set of the corresponding
$K$-submolecule of~$\Gamma$.

Let $v=\Lessthan u{(j+1)}$ and $x=\Lessthan t{(j+1)}$, and write $\sigma=\shape(v)$
and $\theta=\shape(x)$. Then by Remark~\ref{restrictionrem}, we can identify the vertex
$c_{\beta,u}$ of $\Gamma_K$ with $c'_{\delta,v}$ for some
$\delta\in\mathcal{I}_{K,\beta,u,\sigma}$, and
the vertex $c_{\alpha,t}$ of $\Gamma_K$ with $c'_{\gamma,x}$ for some
$\gamma\in\mathcal{I}_{K,\alpha,\lambda,\theta}$. It is clear that $\D(v)=\D(x)\cup\{j\}$,
and it follows that $\mu(c'_{\delta,v},c'_{\gamma,x})=\mu_K(c_{\beta,u},c_{\alpha,t})\neq 0$.
Moreover, since $i\leqslant j$, the restriction number of $(v,x)$ is also $i$. We can
now apply Remark~\ref{nminus1iplusonest} with $j+1$ in place of~$n$ and $\Gamma_K$ in
place of~$\Gamma$, and with $(v,x)$ in place of the $(u,t)$ of
Remark~\ref{nminus1iplusonest}. The conclusion is that $\col_v(i+1)<\col_x(i+1)$
unless $v=s_jx>x$, in which case $j=i+1$. So $\col_u(i+1)<\col_t(i+1)$ unless $u=s_{i+1}t>t$.
\end{rema}
We end this section with two technical lemmas that will be used throughout
Sections \ref{sec:8ex}~and~\ref{sec:9}. They will be useful when the Polygon Rule
is to be applied.

Recall that if $t\in\STD(n)$ and $i\in[1,n-1]$ then $s_it\in\STD(n)$ if and only if
$i\in\SA(t)\cup\SD(t)$.

\begin{lemm}\label{variousaltpaths}
Let $t\in\STD(n)$ and let $i\in\A(t)$ and $j\in\SD(t)$. Put $v=s_jt$.
\begin{enumerate}[label=(\roman*),topsep=1 pt]
\item\label{lemma7.29_i} Suppose that $i<j-1$. Then $i\notin\D(v)$ and $j\notin\D(v)$.

\noindent
Additionally, if $i\in\SA(v)$ then $i\in\D(s_{i}v)$ and $j\notin\D(s_{i}v)$.
\item\label{lemma7.29_ii} Suppose that $i=j-1$ and $\col_t(j+1)>\col_t(j-1)$.
Then $j-1\notin\D(v)$ and $j\notin\D(v)$.

\noindent
Additionally, if $j-1\in\SA(v)$ then
$j-1\in\D(s_{j-1}v)$ and $j\notin\D(s_{j-1}v)$.
\item\label{lemma7.29_iii} Suppose that $i=j-1$ and $\col_t(j+1)<\col_t(j-1)$. Then
$j-1\in\SD(v)$. Writing $w=s_{j-1}v$, we have
$j-1\in\D(v)$ and $j\notin\D(v)$, and $j-1\notin\D(w)$ and
$j\notin\D(w)$.

\noindent
Additionally, if $j\in\SA(w)$, then $j-1\in\SA(s_jw)$,
and we have $j\in\D(s_jw)$ and $j-1\notin\D(s_{j}w)$, and
$j-1\in\D(s_{j-1}s_jw)$ and $j\notin\D(s_{j-1}s_jw)$.
\end{enumerate}
\end{lemm}

\begin{proof} \ \\*[-1.2em]

\ref{lemma7.29_i}~Since $v=s_jt$ and $j\in\SD(t)$, it follows that $j\in\SA(v)$,
whence $j\notin\D(v)$.
Since $v$ is obtained from $t$ by switching the positions of~$j$ and $j+1$,
and since $i+1<j$, it follows that $i$ and $i+1$ have the same row and column
index in $v$ as they have in $t$. Since $i\notin\D(t)$, this shows that
$i\notin\D(v)$.

If $i\in\SA(v)$ then $s_iv$ is standard and $i\in\D(s_iv)$.
Since $s_iv$ is obtained from $v$ by switching $i$~and~$i+1$,
and since $j>i+1$, it follows that $j$ and $j+1$ have the same row and column
index in $s_iv$ as in $v$. Since $j\notin\D(v)$ it follows
that $j\notin\D(s_iv)$.

\ref{lemma7.29_ii}~Since $v=s_jt$ and $j\in\SD(t)$, it follows that $j\in\SA(v)$,
whence $j\notin\D(v)$. Now since $\col_v(j-1)=\col_t(j-1)$ and $\col_v(j)=\col_t(j+1)$,
and $\col_t(j-1)<\col_t(j+1)$ by assumption, it follows that $\col_v(j-1)<\col_v(j)$.
That is, $j-1\notin\D(v)$.

If $j-1\in\SA(v)$ then $s_{j-1}v$ is standard and $j-1\in\D(s_{j-1}v)$.
Since $j-1$ and $j$ are both ascents of $v$, we have $\col_v(j-1)<\col_v(j)<\col_v(j+1)$,
and since $s_{j-1}v$ is obtained from $v$ by switching $j-1$ and $j$, we have
$\col_{s_{j-1}v}(j)=\col_v(j-1)$ and $\col_{s_{j-1}v}(j+1)=\col_v(j+1)$, and
it follows that $\col_{s_{j-1}v}(j)<\col_{s_{j-1}v}(j+1)$. Thus $j\notin\D(s_{j-1}v)$.

\ref{lemma7.29_iii}~As in~\ref{lemma7.29_i} and~\ref{lemma7.29_ii} we have $j\notin\D(v)$. The assumption $\col_t(j+1)<\col_t(j-1)$
gives $\col_v(j)<\col_v(j-1)$, and so $j-1\in\SD(v)$. Hence $w=s_{j-1}v$ is standard,
and $j-1\in\SA(w)$. Since $\col_w(j+1)=\col_v(j+1)=\col_t(j)$ and
$\col_w(j)=\col_v(j-1)=\col_t(j-1)$, and since $j-1\in\A(t)$ by assumption, it follows
that $j\in \A(w)$. Thus $j-1\in\D(v)$ and $j\notin\D(v)$, and $j-1\notin\D(w)$ and
$j\notin\D(w)$, as required.

If $j\in\SA(w)$ then $s_jw\in\STD(\lambda)$. Since $j-1$ and $j$ are both strong ascents of
$w$, we have $\row_w(j-1)>\row_w(j)>\row_w(j+1)$, and since $s_jw$ is obtained from $w$
by switching $j$ and $j+1$, we have $\row_{s_jw}(j-1)=\row_w(j-1)$ and
$\row_{s_jw}(j)=\row_w(j+1)$, and it follows that $\row_{s_jw}(j-1)>\row_{s_jw}(j)$.
Thus $j-1\in\SA(s_jw)$.

Now $j-1\in\SA(s_jw)$ gives $j-1\notin\D(s_jw)$, and gives
$j-1\in\D(s_{j-1}s_jw)$. Similarly, $j\in\SA(w)$ gives
$j\in\D(s_jw)$. Finally, the assumption $\col_t(j+1)<\col_t(j-1)$ gives
$\col_{s_{j-1}s_jw}(j)=\col_{s_jw}(j-1)=\col_t(j+1)<\col_t(j-1)=\col_{s_jw}(j+1)
=\col_{s_{j-1}s_jw}(j+1)$, and $j\notin\D(s_{j-1}s_jw)$.
\end{proof}

Recall from Remark~\ref{LexTableauEquiv} that if $\lambda\in P(n)$ and
$u,\,t\in\STD(\lambda)$ then $t>_{\lex}u$ if and only if there exists $l\in[1,n]$
such that $\col_t(l)<\col_u(l)$ and $\morethan tl=\morethan ul$.
\begin{lemm}\label{xlexlessthant}
Let $\lambda\in P(n)$ and $0\leqslant i\leqslant n-1$. Let $t,\,t'\in\STD(\lambda)$ satisfy
$\morethan ti=\morethan {(t')}i$. Let $j\in \SD(t)$ and put $v=s_jt$, and suppose that
$i\in\A(t)$ and $i<j$. Then $v<_{\lex}t'$, and the following all hold.
\begin{enumerate}[label=(\roman*)]
\item\label{lemma7.30_i} If $i\in\SA(v)$ then $s_iv\in\STD(\lambda)$ and $s_iv<_{\lex}t'$.
\item\label{lemma7.30_ii} If $y\in\STD(\lambda)$ and $y<v$ then $y<_{\lex}t'$.
\item\label{lemma7.30_iii} Suppose that $i=j-1$ and that $\col_t(j+1)<\col_t(j-1)$, and
let $w=s_{j-1}v$. Then $w\in\STD(\lambda)$ and $w<_{\lex}t'$.
If $j\in\SA(w)$ then $s_{j-1}s_jw\in\STD(\lambda)$ and $s_{j-1}s_jw<_{\lex}t'$.
\item\label{lemma7.30_iv} Suppose that $i=j-1$ and that $\col_t(j+1)<\col_t(j-1)$, and
let $w=s_{j-1}v$. Let $x\in\STD(\lambda)$ be such that $x<w$ and $\D(x)$
contains exactly one of $j-1$ or $j$, and let $y$ be the $(j-1)$-neighbour
of $x$ (see Definition~\ref{k-neighbourdefinition}). Then $y<_{\lex}t'$.
\end{enumerate}
\end{lemm}

\begin{proof}
Since $j\in\SD(t)$ we have $t>s_jt=v$, and hence
$t>_{\lex}v$ by Corollary~\ref{domimplieslextableau}. Indeed,
$\col_t(j+1)<\col_t(j)=\col_v(j+1)$ and $\morethan t{(j+1)}=\morethan v{(j+1)}$.
Since $\morethan ti=\morethan {(t')}i$ and $j+1>i$ it follows that
$\col_{t'}(j+1)<\col_v(j+1)$ and
$\morethan {(t')}{(j+1)}=\morethan v{(j+1)}$,
giving $t'>_{\lex}v$.

\ref{lemma7.30_i}\quad The assumption $i\in\SA(v)$ gives $s_iv\in\STD(\lambda)$,
and since $j+1>i+1$ it follows that $\col_{t'}(j+1)<\col_v(j+1)=\col_{s_iv}(j+1)$
and $\morethan{(t')}{(j+1)}=\morethan{(s_iv)}{(j+1)}$. So $t'>_{\lex}s_iv$.

\ref{lemma7.30_ii}\quad If $y<v$ then $y<_{\lex} v$, by Corollary~\ref{domimplieslextableau}, and
since $v<_{\lex} t'$ this gives $y<_{\lex}t'$.

\ref{lemma7.30_iii}\quad Since $\col_v(j)=\col_t(j+1)<\col_t(j-1)=\col_v(j-1)$, we have
$j-1\in\SD(v)$, and since this gives $s_{j-1}v\in\STD(\lambda)$, an argument similar
to that for~\ref{lemma7.30_i} yields $w<_{\lex}t'$.

If $j\in\SA(w)$ then $s_jw\in\STD(\lambda)$. Since $j-1\in\SA(s_jw)$
by Lemma~\ref{variousaltpaths}~\ref{lemma7.29_iii}, we also see that $s_{j-1}s_jw\in\STD(\lambda)$.
Since $\col_t(j+1)<\col_t(j-1)=\col_{s_{j-1}s_jw}(j+1)$, and
since $j+1>i+1$, it follows that $\col_{t'}(j+1)<\col_{s_{j-1}s_jw}(j+1)$ and
$\morethan{(t')}{(j+1)}=\morethan{(s_{j-1}s_jw)}{(j+1)}$. Hence $t'>_{\lex}s_{j-1}s_jw$.

\ref{lemma7.30_iv}\quad There are two cases to consider.

\begin{Case}{1.} Suppose that $\D(x)\cap\{j-1,j\}=\{j-1\}$ and $\D(y)\cap\{j-1,j\}=\{j\}$.
Then either $y=s_jx>x$ or $y=s_{j-1}x<x$.

Suppose first that $y=s_{j}x>x$. Since $x<w$ and $w=s_{j-1}v<v$ by the proof of~\ref{lemma7.30_iii}, it
follows that $x<v$. Since $v<s_jv=t$ and $x<s_jx=y$, it follows by
 Lemma~\ref{characterisationofdominanceorderontab} that $y<t$. Hence
$y<_{\lex}t$ by Corollary~\ref{domimplieslextableau}. That is, there exists
$l\in[1,n]$ such that $\col_{t}(l)<\col_y(l)$ and
$\morethan tl=\morethan yl$. Suppose, for a contradiction, that $l\leqslant j-1$.
Then $\morethan t{(j-1)}=\morethan y{(j-1)}$, giving
$\morethan {(s_jt)}{(j-1)}=\morethan{(s_jy)}{(j-1)}$, that is,
$\morethan v{(j-1)}=\morethan x{(j-1)}$. Therefore $\morethan vj=\morethan xj$, which
gives $\morethan wj=\morethan{(s_{j-1}v)}j=\morethan vj=\morethan xj$.
Now since $\col_x(j)=\col_v(j)<\col_v(j-1)=\col_w(j)$ (using again
$s_{j-1}v<v$), it follows that $w<_{\lex}x$,
contradicting the assumption that $x<w$. Thus $l\geqslant j$. Since
$\morethan t{(j-1)}=\morethan {t'}{(j-1)}$, it follows that $\col_{t'}(l)=\col_t(l)<\col_y(l)$
and $\morethan {t'}l=\morethan yl$, and hence $y<_{\lex}t'$, as required.

Suppose now that $y=s_{j-1}x<x$. Since $x<w$, we have $y<w$,
and by Corollary~\ref{domimplieslextableau} this gives $y<_{\lex}w$.
But since $w<_{\lex}t'$ by~\ref{lemma7.30_iii}, this yields $y<_{\lex}t'$.
\end{Case}

\begin{Case}{2.} Suppose that $\D(x)\cap\{j-1,j\}=\{j\}$ and $\D(y)\cap\{j-1,j\}=\{j-1\}$.
Then either $y=s_jx<x$ or $y=s_{j-1}x>x$.

Suppose first that $y=s_{j-1}x>x$. Since $x<s_{j-1}x=y$ and $w<s_{j-1}w=v$, the assumption
$x<w$ gives $y<v$ by Lemma~\ref{characterisationofdominanceorderontab}. Thus
$y<_{\lex}t'$ by~\ref{lemma7.30_ii}.

Suppose now that $y=s_jx<x$. Since $x<w$, we have $y<w$, and by
Corollary~\ref{domimplieslextableau} this gives $y<_{\lex}w$. But since
$w<_{\lex}t'$ by~\ref{lemma7.30_iii}, this yields $y<_{\lex}t'$.\qedhere
\end{Case}
\end{proof}

\section{Ordered admissible \texorpdfstring{$W$}{W}-graphs in type \texorpdfstring{$A$}{A}}
\label{sec:8ex}

Let $\Gamma=\Gamma(C,\mu,\tau)$ be an admissible $W_n$-graph, and let
$\Lambda\subseteq P(n)$ be the set of molecule types for $\Gamma$. As in
Remark~\ref{vertexsetofWngraph} we write
\begin{equation*}
C=\bigsqcup_{\lambda\in\Lambda}\bigsqcup_{\alpha\in\mathcal{I}_{\lambda}}C_{\alpha,\lambda},
\end{equation*}
where for each $\lambda\in\Lambda$ the set $\mathcal{I}_{\lambda}$ indexes
the molecules of $\Gamma$ of type~$\lambda$, and for each $\lambda\in\Lambda$
and $\alpha\in\mathcal{I}_{\lambda}$ the set
$C_{\alpha,\lambda}=\{c_{\alpha,t}\mid t\in\STD(\lambda)\}$ is the vertex
set of a molecule of type~$\lambda$.

Recall that, by Theorem~\ref{combinatorialCharacterisation}, $\Gamma$ satisfies the
Compatibility Rule, the Simplicity Rule,
the Bonding Rule and the Polygon Rule.
In particular,
in view of Definition~\ref{simplicity}, it follows that whenever vertices $c_{\alpha,t}$
and $c_{\beta,u}$ belong to different molecules and
$\mu(c_{\beta,u},c_{\alpha,t})\neq 0$, we must have $\D(t)\subsetneqq\D(u)$ and
$\mu(c_{\alpha,t},c_{\beta,u})=0$.

We make the following definition.
\begin{defi}\label{orderedWgraphdef}
Let $\Gamma=\Gamma(C,\mu,\tau)$ be an admissible $W_n$-graph, and let
\begin{equation*}
C=\bigsqcup_{\lambda\in\Lambda}\bigsqcup_{\alpha\in\mathcal{I}_{\lambda}}C_{\alpha,\lambda},
\end{equation*}
as above. Then $\Gamma$ is said to be \emph{ordered}
if for all vertices $c_{\alpha,t}$ and $c_{\beta,u}$ with
$\mu(c_{\beta,u},c_{\alpha,t})\neq 0$, either $u<t$ (in the extended
Bruhat order) or else $\alpha=\beta$ and $u=st>t$ for some $s\in S_n$.
\end{defi}

Note that $\mu(c_{\beta,u},c_{\alpha,t})\neq 0$
implies that $\D(u)\nsubseteq\D(t)$. In particular, since $S_1=\emptyset$,
the condition $\mu(c_{\beta,u},c_{\alpha,t})\neq 0$ can never be satisfied
in the case $n=1$. Thus it is vacuously true that every $W_1$-graph is ordered.

Our main objective in this section is to prove the next theorem, which is
one of the main results of this paper.

\begin{theo}\label{inductivestep}
All admissible $W_n$-graphs are ordered.
\end{theo}

The proof of Theorem~\ref{inductivestep} will proceed by induction on~$n$.
Accordingly, we assume now that $n$ is a positive integer and that all admissible
$W_m$-graphs are ordered for $1\leqslant m<n$.
We let $\Gamma=\Gamma(C,\mu,\tau)$ be an admissible $W_n$-graph, and use the
notation introduced in the preamble to this section: $\Lambda$ is the set of molecule
types of $\Gamma$, and for each $\lambda\in\Lambda$ the set $\mathcal{I}_\lambda$
indexes the molecules of type~$\lambda$. We fix $K=S_n\setminus\{s_{n-1}\}$ and
$L=S_n\setminus\{s_1\}$, and we let $\Gamma_K$
and $\Gamma_L$ be the $W_K$-graph and $W_L$-graph obtained by restricting
$\Gamma$ to $W_K$ and $W_L$. Since $|K|=|L|=n-1$, the
inductive hypothesis tells us that $\Gamma_K$ and $\Gamma_L$ are ordered.

For $\lambda\in\Lambda$ and $\alpha\in\mathcal{I}_\lambda$ let $\Theta_{\alpha,\lambda}$
be the molecule of $\Gamma$ with vertex set~$C_{\alpha,\lambda}$, and consider
the $W_K$-restriction of~$\Theta_{\alpha,\lambda}$, as in Remark~\ref{restrictionrem}.
Write $\Lambda_{K,\alpha,\lambda}$ for the set of all $\kappa\in P(n-1)$
such that $\Theta_{\alpha,\lambda}$ contains a $K$-submolecule of type $\kappa$,
and for each $\kappa\in\Lambda_{K,\alpha,\lambda}$ let $\mathcal{I}_{K,\alpha,\lambda,\kappa}$
be a set that indexes the $K$-submolecules of~$\Theta_{\alpha,\lambda}$ of type $\kappa$.
For each $\gamma\in\mathcal{I}_{K,\alpha,\lambda,\kappa}$ let
$\{\,c'_{\gamma,x}\mid x\in\STD(\kappa)\,\}$ be the vertex set of the corresponding
$K$-submolecule of~$\Theta_{\alpha,\lambda}$. Now
$\Lambda_K=\bigcup_{\alpha,\lambda}\Lambda_{K,\alpha,\lambda}$ is the set of molecule
types for $\Gamma_K$, for each $\kappa\in\Lambda_K$ the set
$\mathcal{I}_{K,\kappa}=\bigsqcup_{\{\alpha,\lambda\mid \kappa\in\Lambda_{K,\alpha,\lambda}\}}\mathcal{I}_{K,\alpha,\lambda,\kappa}$
indexes the $K$-submolecules of~$\Gamma$ of type~$\kappa$, and since the vertex set of
$\Gamma_K$ is~$C$ we deduce that
\begin{equation}\label{verticesofGammaK}
C = \bigsqcup_{\kappa\in\Lambda_K}\{c'_{\gamma,x}\mid (\gamma,x)\in\mathcal{I}_{K,\kappa}\times\STD(\kappa)\}.
\end{equation}

Similarly, by Remark~\ref{restrictionrem}, the set of molecule types for $\Gamma_L$ is
$\Lambda_L=\bigcup_{\alpha,\lambda}\Lambda_{L,\alpha,\lambda}$,
where $\Lambda_{L,\alpha,\lambda}$ is the set of all $\theta\in P(n-1)$
such that $\Theta_{\alpha,\lambda}$ has an $L$-submolecule of type $\theta$, and
for each $\theta\in\Lambda_L$ the $L$-submolecules of $\Gamma$ of type $\theta$
are indexed by
$\mathcal{I}_{L,\theta}=\bigsqcup_{\{\alpha,\lambda\mid \theta\in\Lambda_{L,\alpha,\lambda}\}}
\mathcal{I}_{L,\alpha,\lambda,\theta}$, where $\mathcal{I}_{L,\alpha,\lambda,\theta}$
indexes the $L$-submolecules of type $\theta$ in $\Theta_{\alpha,\lambda}$. The vertex
set of $\Gamma_L$ is
\begin{equation}\label{verticesofGammaL}
C = \bigsqcup_{\theta\in\Lambda_L}\{c''_{\epsilon,y}\mid (\epsilon,y)\in\mathcal{I}_{L,\theta}\times\STD(\theta)\}.
\end{equation}

\begin{lemm}\label{kldominancesamenolecule0}
Let $\pi,\,\lambda\in\Lambda$ with $\pi\leqslant\lambda$, and
let $(\beta,u)\in\mathcal{I}_{\pi}\times\STD(\pi)$ and
$(\alpha,t)\in\mathcal{I}_{\lambda}\times\STD(\lambda)$ satisfy the condition
$\mu(c_{\beta,u},c_{\alpha,t}) \neq 0$ and $\D(t) \subsetneqq \D(u)$.
Let $j=\min(\D(u)\setminus\D(t))$ and assume that $j<n-1$. Then $u < t$ unless
$\alpha = \beta$ and $u=s_j t>t$.
\end{lemm}

\begin{proof}
Since $j$ is at least $1$, the requirement that $n-1>j$ implies that $n\geqslant 3$.
Let $v=\Lessthan u{(n-1)}$ and $x=\Lessthan t{(n-1)}$, and write
$\sigma=\shape(v)$ and $\theta=\shape(x)$. By Eq.~\eqref{verticesofGammaK} we can
identify $c_{\beta,u}\in C$ with
$c'_{\delta,v}$ for some $\delta\in\mathcal{I}_{K,\beta,\pi,\sigma}$, and
$c_{\alpha,t}\in C$ with $c'_{\gamma,x}$ for some
$\gamma\in\mathcal{I}_{K,\alpha,\lambda,\theta}$.

Since $\D(v)\setminus\D(x)=\D(u)\cap[1,n-2])\setminus(\D(t)\cap[1,n-2])$, and since we are
given that $j\in\D(u)\setminus\D(t)$ and $j<n-1$,
it follows that
$\mu(c'_{\delta,v},c'_{\gamma,x})=\mu_K(c_{\beta,u},c_{\alpha,t})\neq 0$.
Since $\Gamma_K$ is ordered, we have either $v<x$ or $\gamma=\delta$ and
$v=s_ix>x$ for some $i\in[1,n-2]$. In the former case, since
$\shape(u)=\pi\leqslant\lambda=\shape(t)$ by hypothesis, and
since $\Lessthan u{(n-1)}=v<x=\Lessthan t{(n-1)}$,
we have $u<t$ by the remark following Definition~\ref{ExtDominance Order Tableaux}
In the latter case, the fact that $\gamma=\delta$ implies that
$\mathcal{I}_{K,\alpha,\lambda,\theta}=\mathcal{I}_{K,\beta,\pi,\sigma}$
and hence that $\alpha=\beta$. Moreover, since $v=s_ix>x$ it follows
from Lemma~\ref{dessjtminust} that $i$ is the unique element of
$\D(v)\setminus\D(x)$. So $i=j$, and $u=s_jt>t$.
\end{proof}

\begin{prop}\label{kldominancesamenolecule}
Let $\pi,\,\lambda\in\Lambda$ with $\pi\leqslant\lambda$, and
suppose that $(\beta,u)\in\mathcal{I}_{\pi}\times\STD(\pi)$ and
$(\alpha,t)\in\mathcal{I}_{\lambda}\times\STD(\lambda)$ satisfy
$\mu(c_{\beta,u},c_{\alpha,t}) \neq 0$. Then $u < t$ unless
$\alpha = \beta$ and $u=s_i t>t$ for some $i\in[1,n-1]$.
\end{prop}

\begin{proof}
Since $\mu(c_{\beta,u},c_{\alpha,t}) \neq 0$, it follows that
$\D(u)\nsubseteq\D(t)$. If $\D(t)\nsubseteq\D(u)$ also holds then the
Simplicity Rule
shows that
$\{c_{\beta,u},c_{\alpha,t}\}$ is a simple edge, so that $\alpha=\beta$ and
$u=s_i t$ for some $i\in[1,n-1]$. Thus we may assume that $\D(t)\subsetneqq\D(u)$.
If $\min(\D(u)\setminus\D(t))<n-1$ then
the result is given by Lemma~\ref{kldominancesamenolecule0}.

It remains to consider the case $\D(u)=\D(t)\cup\{n-1\}$. Let $i$ be the restriction
number of the pair $(u,t)$ and note that $i<n$ by Remark~\ref{k-restrricted-rem}.
If $i=n-1$ or $i=n-2$ then the results are given
by Lemma~\ref{nminus1order0} and Lemma~\ref{nminus1order}, respectively. So we may assume that
$i<n-2$, and it follows by Lemma~\ref{noname1muequallambda} that
$(u,t)\approx_i (v,x)$ for some $(v,x)\in\STD(\pi)\times\STD(\lambda)$ satisfying
the condition $\mu(c_{\beta,u},c_{\alpha,t}) = \mu(c_{\beta,v},c_{\alpha,x}) \neq 0$
and $\D(x)\subsetneqq\D(v)=\D(x)\cup\{i,n-1\}$. We can now apply
Lemma~\ref{kldominancesamenolecule0} with $c_{\beta,v}$ and $c_{\alpha,x}$ in place
of $c_{\beta,u}$ and $c_{\alpha,t}$. Since $\D(v)=\D(x)\cup\{i,n-1\}$ it follows
from Lemma~\ref{dessjtminust} that $v\ne s_ix$, and so we must have $v<x$.
So $u<t$, by Proposition~\ref{bruhatorderpreserved}.
\end{proof}

\begin{nota}\label{inidef}
Given $\lambda\in\Lambda$, let $C'_\lambda=
C\setminus\bigl(\bigsqcup_{\alpha\in\mathcal{I}_{\lambda}}C_{\alpha,\lambda}\bigr)$,
the set of vertices of $\Gamma$ belonging to molecules of type different
from~$\lambda$. We define $\Ini_{\lambda}(\Gamma)$ to be the set of
$(\alpha,t)\in\mathcal{I}_{\lambda}\times\STD(\lambda)$ such that there exists
an arc from $c_{\alpha,t}$ to some vertex in $C'_\lambda$. That is,
\[
\Ini_{\lambda}(\mk \Gamma \mk )\Mk = \Mk 
\left\{\Mk \Mk (\mk \alpha,\mk t\mk )\Mk \in\Mk \mathcal{I}_{\lambda}\Mk \times\Mk \STD(\lambda)\bigm|
\mu(c_{\beta,u},c_{\alpha,t})\Mk \neq\Mk 0 \text{ for some } (\beta,u)\Mk \in\Mk \!\!\!\!
\bigsqcup_{\pi\in\Lambda\setminus\{\lambda\}}\Mk \!\!\!
(\mathcal{I}_{\pi}\Mk \times\Mk \STD(\pi))\Mk \Mk \right\}\Mk .
\]
For each $\alpha\in\mathcal{I}_\lambda$ we define
$\Ini_\Gamma(\alpha,\lambda)=\{\,t\in\STD(\lambda)\mid(\alpha,t)\in\Ini_{\lambda}(\Gamma)\,\}$,
and we also define
$\Ini_\Gamma(\lambda)=\bigcup_{\alpha\in\mathcal{I}_\lambda}\Ini_\Gamma(\alpha,\lambda)$.
If $\Ini_{\lambda}(\Gamma)\neq\emptyset$, so that
$\Ini_\Gamma(\lambda)$ is a nonempty subset of~$\STD(\lambda)$, we define $t_\lambda$ to
be the element of $\Ini_\Gamma(\lambda)$ that is minimal in the lexicographic order
on $\STD(\lambda)$.
\end{nota}

It is clear that whenever $\lambda\in\Lambda$ and $\Ini_{\lambda}(\Gamma)\neq\emptyset$
there must be at least one $\alpha\in\mathcal{I}_\lambda$ having the property that
$t_\lambda\in\Ini_\Gamma(\alpha,\lambda)$, but for an arbitrary
$\alpha\in\mathcal{I}_\lambda$ it may or may not be the case that
$t_\lambda\in\Ini_\Gamma(\alpha,\lambda)$.

The following definitions will be used in results dealing with~$t_\lambda$.

\begin{defi}\label{k-minimalappr}
Let $\pi,\lambda\in P(n)$. Let $(u,t)\in\STD(\pi)\times\STD(\lambda)$, and
let $k$ be the restriction number of $(u,t)$.
The pair $(u,t)$ is said to be $k$-\emph{minimal}, and
$t$ is said to be $k$-\emph{minimal
with respect to $u$}, if $\D(t)\subsetneqq\D(u)$ and $\Morethan tk$
is $k$-critical, and $\lessthan tk$ is the minimal tableau of its shape.
\end{defi}

\begin{exam*} Consider the following three pairs $(u_i,t_i)$:
\begin{align*}
(u_1,t_1)&=\left(\,\vcenter{\hbox{\ttab(123,4,5)}},\,\vcenter{\hbox{\ttab(123,45)}}\right),\\
(u_2,t_2)&=\left(\,\vcenter{\hbox{\ttab(125,3,4)}},\,\vcenter{\hbox{\ttab(123,45)}}\right),\\
(u_3,t_3)&=\left(\,\vcenter{\hbox{\ttab(125,3,4)}},\,\vcenter{\hbox{\ttab(125,34)}}\right).
\end{align*}
The pair $(u_1,t_1)$ is $4$-restricted,
$\lessthan{(t_1)}4=\lessthan{(u_1)}4$ is the minimal tableau of shape~$(1,1,1)$,
and from Definition~\ref{m-critical} it is readily checked that $\Morethan{(t_1)}4$ is
the $4$-critical tableau of shape $(2,2,1)/(1,1,1)$.
Since $\D(t_1)=\{3\}\subsetneqq\{3,4\}=\D(u_1)$, it follows that
$(u_1,t_1)$ is $4$-minimal.

The pair $(u_2,t_2)$ is $2$-restricted, $\lessthan{(t_2)}2$ is minimal, and
$\D(t_2)=\{3\}\subsetneqq\{2,3\}=\D(u_2)$. But
$\Morethan{(t_2)}2$ is not $2$-critical, since $2$ is not in its first
nonempty column. So $(u_2,t_2)$ is not $2$-minimal.

The pair $(u_3,t_3)$ is $3$-restricted, $\lessthan{(t_3)}3$ is minimal,
$\Morethan{(t_3)}3$ is the $3$-critical tableau of shape $(2,2,1)/(1,1,1)$,
and $\D(t_3)=\{2\}\subsetneqq\{2,3\}=\D(u_3)$.
So $(u_3,t_3)$ is $3$-minimal.
\end{exam*}

Let $\pi,\,\lambda\in P(n)$, and let $(u,t)\in\STD(\pi)\times\STD(\lambda)$.
Let $k$ be the restriction number of~$(u,t)$, and assume that
$k\in[1,n-1]$ (or, equivalently, $u\neq t$).
Recall that
\begin{equation*}
F(u,t) = \{(v,x)\in C_k(u,t) \mid v^{-1}(k)=x^{-1}(k) \text{ lies between } u^{-1}(k+1) \text{ and } t^{-1}(k+1)\}.
\end{equation*}

\begin{defi}\label{k-minimalapproximate}
Let $\pi,\lambda\in P(n)$ and $(u,t)\in\STD(\pi)\times\STD(\lambda)$
with $u\neq t$, and let $k$ be the restriction number of~$(u,t)$.
We define $A(u,t) = \{(v,x)\in F(u,t) \mid \col_x(k)=\col_t(k+1)-1\}$ and call any
element of $A(u,t)$ an \emph{approximate} of~$(u,t)$.
\end{defi}
Since every pair $(v,x)\in F(u,t)$ must satisfy either $\col_u(k+1)\leqslant\col_x(k)<\col_t(k+1)$
or $\col_t(k+1)\leqslant\col_x(k)<\col_u(k+1)$, it is clear that $A(u,t)\ne\emptyset$ only
if $\col_u(k+1)<\col_t(k+1)$. Conversely, suppose that $\col_u(k+1)<\col_t(k+1)$.
Let $r=\col_t(k+1)$ and $\xi=\shape(\Lessthan tk)$, and note that $\xi_{r-1}>\xi_r$,
since $t^{-1}(k+1)$ is $\xi$-addable and in column~$r$. Hence $(\xi_{r-1},r-1)$ is
$\xi$-removable, and we can choose a tableau $w'\in\STD(\xi)$ with $w'(\xi_{r-1},r-1)=k$.
The unique pair $(v,x)\in C_k(u,t)$ with $\Lessthan vk=\Lessthan xk=w'$ is then an
element of $A(u,t)$, and so $A(u,t)\ne\emptyset$.

\begin{exem}\label{approxexample}
Let
\[
(u,t)=\left(\,\vcenter{\hbox{\ttab(123,4,5,6,7)}},\,\vcenter{\hbox{\ttab(1236,47,5)}}\right).
\]
We see that $(u,t)$ is $5$-restricted and $\col_u(6)=1<4=\col_t(6)$. By
Definition~\ref{k-minimalapproximate}, the set $A(u,t)$ consists of all
$(v,x)\in F(u,t)$ with $\col_x(5)=\col_t(6)-1=3$. So as well as being
5-restricted and satisfying the conditions $\morethan v5=\morethan u5$ and
$\morethan x5=\morethan t5$, each $(v,x)\in F(u,t)$ satisfies
$v^{-1}(5)=x^{-1}(5)=(1,3)$. The tableau $\Lessthan v4=\Lessthan x4$ can be
any element of~$\STD(3,1)$. So
\[
A(u,t)\Mk =\Mk\left\{\Mk\Mk\left(\,\vcenter{\hbox{\ttab(145,2,3,6,7)}},\vcenter{\hbox{\ttab(1456,27,3)}}\Mk\right)\Mk,
\left(\,\vcenter{\hbox{\ttab(135,2,4,6,7)}}\Mk, \vcenter{\hbox{\ttab(1356,27,4)}}\Mk\right)\Mk,
\left(\,\vcenter{\hbox{\ttab(125,3,4,6,7)}}\Mk,\vcenter{\hbox{\ttab(1256,37,4)}}\Mk\right)\Mk
\Mk\right\}\Mk\Mk.
\]
\end{exem}

\begin{rema}\label{k-1subclassofk}
Let $u,\,t$ be as in Definition~\ref{k-minimalapproximate}, and assume that $A(u,t)\ne\emptyset$.
It is immediate from the definition that every approximate $(v,x)$ of~$(u,t)$ is
$k$-restricted and satisfies $(v,x)\approx_k(u,t)$. Furthermore, if $r=\col_t(k+1)$
then $v^{-1}(k)=x^{-1}(k)=(d,m)$, where $m=r-1$ and $d$ is the $(r-1)$-th part
of $\xi=\shape(\Lessthan tk)$.
So $A(u,t)$ is a nonempty $(k-1)$-subclass of $C_k(u,t)$. Furthermore, if we define
$\kappa\in P(k-1)$ by $[\kappa]=[\xi]\setminus\{(d,m)\}$, then
there is a bijection from $\STD(\kappa)$ to $A(u,t)$ such that $w\mapsto (v,x)$
if and only if $\Lessthan {v}{(k-1)}=\Lessthan {x}{(k-1)}=w$.
We use this bijection to transfer the partial order
$\leqslant$ from $\STD(\kappa)$ to $A(u,t)$, and observe that $A(u,t)$
has a unique minimal element, given by $w=\tau_{\kappa}$, and a unique maximal element,
given by $w=\tau^{\kappa}$. We call these elements of $A(u,t)$
the \emph{minimal approximate} of~$(u,t)$ and the \emph{maximal approximate} of~$(u,t)$,
respectively.
\end{rema}

In Example~\ref{approxexample}, the first and the last are the minimal and maximal approximates of~$(u,t)$.

\begin{rema}\label{nonemptyAkst}
Let $\pi,\,\lambda\in\Lambda$, and
let $(\beta,u)\in\mathcal{I}_{\pi}\times\STD(\pi)$ and
$(\alpha,t)\in\mathcal{I}_{\lambda}\times\STD(\lambda)$ satisfy the condition
$\mu(c_{\beta,u},c_{\alpha,t}) \neq 0$ and $\D(t) \subsetneqq \D(u)$.
Let $k\in[1,n-1]$ be the restriction number of the pair
$(u,t)$. Let $l=\min(\D(u)\setminus\D(t))$, and let
$H=\{s_1,\ldots,s_l\}$. Remark~\ref{iplusonest} applied to $\Gamma_H$,
the $W_H$-restriction of $\Gamma$, shows that $\col_u(k+1)<\col_t(k+1)$
unless $u=s_{k+1}t>t$. It follows that $A(u,t)\ne\emptyset$
unless $u=s_{k+1}t>t$.

Since $u=s_{k+1}t$ forces $\pi=\lambda$, Proposition~\ref{kldominancesamenolecule}
shows that $\mu(c_{\beta,u},c_{\alpha,t}) \neq 0$ and $u=s_{k+1}t>t$ occurs only
if $\alpha=\beta$. Moreover, if $\alpha\neq\beta$ then
$c_{\beta,u}$ and $c_{\alpha,t}$ are in different molecules, and
$\mu(c_{\beta,u},c_{\alpha,t})\neq 0$ implies that $\D(t) \subsetneqq \D(u)$.
So when $\alpha\neq\beta$ the condition $\mu(c_{\beta,u},c_{\alpha,t})\neq 0$
suffices to ensure that $A(u,t)\neq\emptyset$.
\end{rema}

\begin{lemm}\label{k-minimallemma1}
Let $\pi,\lambda\in P(n)$ and $(u,t)\in\STD(\pi)\times\STD(\lambda)$
with $u\neq t$, and let $k$ be the restriction number of~$(u,t)$.
Assume that $A(u,t)\neq\emptyset$, and let
$(v,x)\in A(u,t)$. Then $(v,x)$ is $k$-restricted and satisfies
$(v,x)\approx_k(u,t)$.
Moreover, if $D(t)\subsetneqq\D(u)$ then we have
$\D(x)\subsetneqq\D(v)$ and $k=\min(\D(v)\setminus\D(x))$.
\end{lemm}

\begin{proof}
It follows from Remark~\ref{k-1subclassofk} that $(v,x)$ is $k$-restricted and
satisfies $(v,x)\approx_k(u,t)$. So it remains to show that if $\D(t)\subsetneqq\D(u)$
then $\D(x)\subsetneqq\D(v)$ and $k=\min(\D(v)\setminus\D(x))$. So assume that
$\D(t)\subsetneqq\D(u)$.

Since $\col_u(k+1)<\col_t(k+1)$
(since $A(u,t)\neq\emptyset$), Lemma~\ref{knuthmoveweakorderrelationLem}
and Lemma~\ref{knuthmoveweakorderrelationLemb} show that
$\D(x)\setminus\D(v)=\D(t)\setminus\D(u)$ and
$\D(v)\setminus\D(x)\supseteq\D(u)\setminus\D(t)$, and since
$\D(t)\subsetneqq\D(u)$ it follows that $\D(x)\subsetneqq\D(v)$.
Since $(v,x)$ is favourable, we have $k=\min(\D(v)\oplus\D(x)$
by Remark~\ref{k-reducedrem}, and it follows that $k=\min(\D(v)\setminus\D(x))$.
\end{proof}

\begin{lemm}\label{k-minimallemma}
Let $\pi,\lambda\in\Lambda$ with $\pi\neq\lambda$, and let
$(\beta,u)\in\mathcal{I}_{\pi}\times\STD(\pi)$ and
$(\alpha,t)\in\mathcal{I}_{\lambda}\times\STD(\lambda)$ satisfy
$\mu(c_{\beta,u},c_{\alpha,t}) \neq 0$. Let $k$ be the
restriction number of $(u,t)$. Then $A(u,t)\neq\emptyset$, and for
all $(v,x)\in A(u,t)$ the following three conditions hold\textup{:}
\begin{enumerate}[label=(\roman*),topsep=1 pt]
\item
$(v,x)\approx (u,t)$,
\item
$\D(x)\subsetneqq\D(v)$ and $k=\min(\D(v)\setminus\D(x))$,
\item
$\mu(c_{\beta,v},c_{\alpha,x})=\mu(c_{\beta,u},c_{\alpha,t})$.
\end{enumerate}
\end{lemm}

\begin{proof}
Remark~\ref{nonemptyAkst} gives $A(u,t)\neq\emptyset$. Let
$(v,x)\in A(u,t)$ be arbitrary. By Lemma~\ref{k-minimallemma1}, we have
$(u,t)\approx_k (v,x)$, whence $(u,t)\approx (v,x)$. Since $c_{\beta,u}$
and $c_{\alpha,t}$ are in distinct molecules (since $\pi\neq\lambda$),
and since $\mu(c_{\beta,u},c_{\alpha,t}) \neq 0$, we have $\D(t)\subsetneqq\D(u)$,
and it follows by Lemma~\ref{k-minimallemma1}
that $\D(x)\subsetneqq\D(v)$ and $k=\min(\D(v)\setminus\D(x))$.
It remains to show that
$\mu(c_{\beta,v},c_{\alpha,x})=\mu(c_{\beta,u},c_{\alpha,t})$.
Let $l=\min(\D(u)\setminus\D(t))$, and note that $k\leqslant l$
since $(u,t)$ is $k$-restricted.

Suppose first that $k<l$. Since $(u,t)\approx_k (v,x)$ and
$k<l\in\D(u)\setminus\D(t)$, the result follows from Lemma~\ref{noname0}.

Suppose now that $k=l=\min(\D(u)\setminus\D(t))$. In particular,
this means that $k\in\D(u)\setminus\D(t)$. Let $w=\Lessthan tk=\Lessthan uk$
and $\xi=\shape(w)$, and let $(h,r)=t^{-1}(k+1)$ and $(g,p)=t^{-1}(k)$,
the boxes of $t$ that contain $k+1$ and $k$. Note that $\xi_{r-1}>\xi_r$,
since $(h,r)$ is $\xi$-addable. Since $k\notin\D(t)$, it follows that
$g\geqslant h$ and $p<r$. If $p=r-1$ then $(u,t)\in A(u,t)$, and then
$(v,x)\in C_{k-1}(u,t)$, since Remark~\ref{k-1subclassofk} tells us that $A(u,t)$
is a single $(k-1)$-subclass. The desired conclusion
$\mu(c_{\beta,v},c_{\alpha,x})=\mu(c_{\beta,u},c_{\alpha,t})$
then follows from Lemma~\ref{noname0}, since $k-1<k\in\D(u)\setminus\D(t)$.
Thus we can assume that $p<r-1$.

Let $(d,m)=(\xi_{r-1},r-1)$, noting that $(g,p)$ and $(d,m)$ are distinct
$\xi$-removable boxes. Let $\rho=\shape(\lessthan wk)\in P(k-1)$, so that
$[\rho]=[\xi]\setminus \{(g,p)\}$, and let $\rho'\in P(k-2)$ satisfy
$[\rho']=[\xi]\setminus \{(g,p),(d,m)\}$. Let $(i,j)$ be a
$\rho'$-removable box that lies between $(g,p)$ and $(d,m)$
(in the sense that $g>i\geqslant d$ and $p\leqslant j<m$), and observe
that we can choose $w'\in\STD(\rho)$ satisfying $w'(i,j)=k-2$ and $w'(d,m)=k-1$.
Since $\shape(w)=\shape(\lessthan tk)=\shape(\lessthan uk)$, we can
define $(u_1,t_1)\in C_{k-1}(u,t)$ by the conditions that
$\lessthan {(u_1)}k = w'$ and $\Morethan {(u_1)}k = \Morethan {u}k$,
and $\lessthan {(t_1)}k = w'$ and $\Morethan {(t_1)}k = \Morethan tk$.
Since $(u_1,t_1)\in C_{k-1}(u,t)$ and $k-1<k\in\D(u)\setminus\D(t)$, it follows
from Lemma~\ref{noname0} that
$\mu(c_{\beta,u_1},c_{\alpha,t_1})=\mu(c_{\beta,u},c_{\alpha,t})$ and
$k\in\D(u_1)\setminus\D(t_1)$.

Note that $u_1^{-1}(k)=t_1^{-1}(k)=(g,p)$ and $u_1^{-1}(k-1)=t_1^{-1}(k-1)=(d,m)$.
Since $p<m$ we have that $k-1\in\SD(u_1)\cap\SD(t_1)$, and therefore we can define
$(u_2,t_2)\in\STD(\pi)\times\STD(\lambda)$ by $(u_2,t_2)=(s_{k-1}u_1,s_{k-1}t_1)$.
Since $\Lessthan{(u_1)}k=\Lessthan{(t_1)}k$, clearly also
$\Lessthan{(u_2)}k=\Lessthan{(t_2)}k$. Furthermore,
$\morethan{(u_2)}k=\morethan{(u_1)}k=\morethan uk$ and
$\morethan{(t_2)}k=\morethan{(t_1)}k=\morethan tk$. Thus it follows that
$(u_2,t_2)\in C_k(u,t)$. We check that in fact $(u_2,t_2)\in A(u,t)$. One of the
requirements is that $\col_{t_2}(k)=\col_t(k+1)-1$, which is satisfied
since
\[
t_2^{-1}(k)=t_1^{-1}(k-1)=(w')^{-1}(k-1)=(d,m)=(\xi_{r-1},r-1),
\]
while $t^{-1}(k+1)=(h,r)$. The other requirement is that $t_2^{-1}(k)$ lies
between $u^{-1}(k+1)$ and $t^{-1}(k+1)$, and this holds since
\[
\col_t(k+1)=r>\col_{t_2}(k)=r-1\geqslant p=\col_t(k)=\col_u(k)\geqslant\col_u(k+1),
\]
the last inequality because $k\in\D(u)$. Since we also have $(v,x)\in A(u,t)$,
it follows from Remark~\ref{k-1subclassofk} that $(u_2,t_2)\in C_{k-1}(v,x)$,
and since $k-1<k\in\D(v)\setminus\D(x)$ it follows from Lemma~\ref{noname0}
that $\mu(c_{\beta,u_2},c_{\alpha,t_2})=\mu(c_{\beta,v},c_{\alpha,x})$
and $k\in\D(u_2)\setminus\D(t_2)$.

Our task is now reduced to proving that
$\mu(c_{\beta,u_1},c_{\alpha,t_1})=\mu(c_{\beta,u_2},c_{\alpha,t_2})$,
and for this we apply Proposition~\ref{arctransport} with
$\{c_{\beta,u_1},c_{\beta,u_2}\}$ in place of $\{v,v'\}$ and
$\{c_{\alpha,t_1},c_{\alpha,t_2}\}$ in place of $\{u,u'\}$, and with
$\{s_{k-2},s_{k-1}\}$ in place of~$\{s,t\}$ and $s_k$ in place of~$r$.

We must first check that $\{c_{\beta,u_1},c_{\beta,u_2}\}$ and
$\{c_{\alpha,t_1},c_{\alpha,t_2}\}$ are simple edges of~$\Gamma$. For this
it suffices to show that there are DKMs taking $u_2$ to $u_1$ and
taking $t_2$ to~$t_1$. Observe first that
$\col_{u_1}(k-2)=\col_{t_1}(k-2)=j$, which gives
$\col_{u_1}(k)\leqslant\col_{u_1}(k-2)<\col_{u_1}(k-1)$ and
$\col_{t_1}(k)\leqslant\col_{t_1}(k-2)<\col_{t_1}(k-1)$,
since $p\leqslant j<m$. Applying $s_{k-1}$ interchanges $k-1$ and~$k$,
and so
$\col_{u_2}(k-1)\leqslant\col_{u_2}(k-2)<\col_{u_1}(k)$ and
$\col_{t_2}(k-1)\leqslant\col_{t_2}(k-2)<\col_{t_2}(k)$.
Now since $\col_{u_2}(k-1)\leqslant\col_{u_2}(k-2)$ and $\col_{u_2}(k-1)<\col_{u_2}(k)$,
it follows from Definition~\ref{ascentsdescentsdef} that $k-2\in\D(u_2)$
and $k-1\notin\D(u_2)$. So $\D(u_2)\cap \{k-2,k-1\}=\{k-2\}$. In a similar
way, $\D(u_1)\cap \{k-2,k-1\}=\{k-1\}$. Hence, since $u_1=s_{k-1}u_2$, it follows
from the definition that $u_2\to u_1$ is a DKM (of the first kind) of index $k-1$.
Completely analogous reasoning shows that $t_2\to t_1$ is also a DKM (of the
first kind) of index $k-1$. Hence $\{c_{\beta,u_1},c_{\beta,u_2}\}$ and
$\{c_{\alpha,t_1},c_{\alpha,t_2}\}$ are simple edges, as required.

To show that the hypotheses of Proposition~\ref{arctransport} are satisfied, it
remains to check that
\begin{align*}
\D(t_1)\cap \{k-2,k-1,k\} &= \{k-1\},&\D(u_1)\cap \{k-2,k-1,k\} &= \{k-1,k\},\\
\D(t_2)\cap \{k-2,k-1,k\} &= \{k-2\},&\D(u_2)\cap \{k-2,k-1,k\} &= \{k-2,k\}.
\end{align*}
\looseness-1
But these are just the DKM conditions established in the previous paragraph together
with $k\in \D(u_1)\setminus\D(t_1)$ and $k\in \D(u_2)\setminus\D(t_2)$, both
of which have been established above. Hence Proposition~\ref{arctransport} applies,
and $\mu(c_{\beta,u_2},c_{\alpha,t_2})=\mu(c_{\beta,u_1},c_{\alpha,t_1})$, as required.
\end{proof}

\begin{prop}\label{cellorder}
Let $\lambda,\,\pi\in\Lambda$ with $\Ini_\lambda(\Gamma)\ne\emptyset$ and
$\pi\in\Lambda\setminus\{\lambda\}$, and let $\alpha\in\mathcal I_\lambda$
and $(\beta,u')\in\mathcal I_\pi\times\STD(\pi)$ satisfy
$\mu(c_{\beta,u'},c_{\alpha,t_\lambda})\neq 0$.
Let $k$ be the restriction number of~$(u',t_\lambda)$, and let $(u,t)\in A(u',t_\lambda)$.
Then $\Morethan tk$ is $k$-critical. In particular, if $(u,t)$ is
the minimal approximate of $(u',t_\lambda)$ then $t$ is $k$-minimal
with respect to~$u$.
\end{prop}

\begin{proof}
Lemma~\ref{k-minimallemma} tells us that $(u,t)\approx (u',t_\lambda)$, that
$\D(t)\subsetneqq\D(u)$ and $k=\min(\D(u)\setminus\D(t))$,
and that $\mu(c_{\beta,u},c_{\alpha,t})=\mu(c_{\beta,u'},c_{\alpha,t_\lambda})\neq 0$.
Note that $\col_t(k+1)=\col_t(k)+1$, since $(u,t)$ is an approximate of
$(u',t_\lambda)$ (see Definition~\ref{k-minimalapproximate}). Thus, by Lemma~\ref{critical-tableau},
to show that $\Morethan tk$
is $k$-critical it will suffice to show that every $j\in\D(t)$ with
$j>k+1$ is in $\WD(t)$, and that either $\col_t(k+2)=\col_t(k)$
or $k+1\notin\SD(t)$. We do both parts of this by contradiction.

For the first part, suppose that $j>k+1$ and $j\in\SD(t)$. Let $v=s_jt$,
which is standard since $j\in\SD(t)$. Observe that the conditions of
Lemma~\ref{xlexlessthant} are satisfied with $k$ in place of~$i$
and $t_\lambda$ in place of~$t'$. Since $k+1<j$, the conditions
$0\leqslant k\leqslant n-1$ and $k<j$ are certainly satisfied, and
we also have $j\in\SD(t)$ and $v=s_jt$. The condition
$\morethan {\morethan tk=(t_\lambda)}k$ holds since $(u,t)\approx_k(u',t_\lambda)$
(see Remark~\ref{C_m(u,t)-rmk}), and $k\in\A(t)$ holds since
$k\in\D(u)\setminus\D(t)$ (and $\A(t)$ is the complement of $\D(t)$).
So Lemma~\ref{xlexlessthant} applies, and in particular shows
that $v<_{\lex}t_\lambda$.

Note that $(c_{\alpha,t},c_{\beta,u})$ is an arc in~$\Gamma$, since
$\mu(c_{\beta,u},c_{\alpha,t})\neq 0$, and note that $(c_{\alpha,v},c_{\alpha,t})$
is also an arc, since $\mu(c_{\alpha,t},c_{\alpha,v})=1$ by Corollary~\ref{arweightone}.
We claim that the directed path $(c_{\alpha,v},c_{\alpha,t},c_{\beta,u})$
is of alternating type~$(s_j,s_k)$. That is, we claim that $k,j\notin\D(v)$,
that $j\in\D(t)$ and $k\notin\D(t)$, and that $k,j\in\D(u)$.
The first of these follows from Lemma~\ref{variousaltpaths}~\ref{lemma7.29_i}, since
$v=s_jt$ and $k<j-1$, and we have $k\in\A(t)$ and $j\in\SD(t)$.
The others follow from facts established above, namely that $k\in\D(u)\setminus\D(t)$
and $j\in\SD(t)\subseteq\D(t)\subsetneqq\D(u)$.

\looseness-1
To simplify the notation slightly, if $c,\,\mk d \Mk \in \Mk C$ we write
$N^h_{j,k}(c,d)$ for $N^h_{s_j,s_k}(\Gamma;c,d)$ (de\-fined in Eq.~\eqref{altsums}
in Section~\ref{sec:5} above). Since $\Gamma$ is admissible the quantity~$N^{2}_{j,k}(c_{\alpha,v},c_{\beta,u})$ is a sum of positive terms, and there
is at least one term since $(c_{\alpha,v},c_{\alpha,t},c_{\beta,u})$ is an
alternating path of type~$(j,k)$. Hence
$N^{2}_{j,k}(c_{\alpha,v},c_{\beta,u})\ne0$, and since $\Gamma$ satisfies
the Polygon Rule it follows from Definition~\ref{polygon} that
$N^{2}_{k,j}(c_{\alpha,v},c_{\beta,u})\ne0$. So there must exist at least
one $\nu\in\Lambda$ and one
$(\gamma,y)\in\mathcal{I}_{\nu}\times\STD(\nu)$
such that $(c_{\alpha,v},c_{\gamma,y},c_{\beta,u})$ is an alternating directed
path of type $(k,j)$. If $\nu\ne\lambda$ this implies that
$(\alpha,v)\in\Ini_\lambda(\Gamma)$, and
$v\in\Ini_\Gamma(\alpha,\lambda)\subseteq\Ini_\Gamma(\lambda)$.
But this is impossible since $t_\lambda$ is lexicographically minimal
in $\Ini_\Gamma(\lambda)$, and we have shown above that $v<_{\lex}t_\lambda$.
So $\nu=\lambda$. Hence, by Proposition~\ref{kldominancesamenolecule}, we must
have either $y<v$, or $\gamma=\alpha$ and $y=s_kv>v$.

If $y=s_kv>v$ then Lemma~\ref{xlexlessthant}~\ref{lemma7.30_i} applies (since the condition
$k\in\SA(v)$ is equivalent to $v<s_kv\in\STD(\lambda)$),
giving $y=s_kv <_{\lex}t_\lambda$. On the other hand, if $y<v$ then
Lemma~\ref{xlexlessthant}~\ref{lemma7.30_ii} applies, again giving $y<_{\lex}t_\lambda$.
But since $(c_{\gamma,y},c_{\beta,u})$ is an arc from a molecule of type $\nu=\lambda$
to a molecule of type $\pi\ne\lambda$ it follows that
$y\in\Ini_\Gamma(\lambda)$, and $y<_{\lex}t_\lambda$ contradicts the
minimality of~$t_\lambda$. This completes
the proof of the first part.

For the second part, suppose that $k+1\in\SD(t)$ and
$\col_t(k+2)\neq\col_t(k)$.

Since $(u,t)\in A(u',t_\lambda)$, it follows from Definition~\ref{k-minimalapproximate}
that $\col_t(k)=\col_{t_\lambda}(k+1)-1$, and therefore
$\col_t(k)=\col_t(k+1)-1$, since $\morethan tk=\morethan{(t_\lambda)}k$.
But $\col_t(k+1)>\col_t(k+2)$ since $k+1\in\SD(t)$, and so
$\col_t(k)\geqslant\col_t(k+2)$. Thus the assumption $\col_t(k+2)\neq\col_t(k)$
in fact means that $\col_t(k+2)<\col_t(k)$.

Let $v=s_{k+1}t$, noting that $v\in\STD(\lambda)$ since $k+1\in\SD(t)$.
We now apply Lemma~\ref{variousaltpaths}~\ref{lemma7.29_iii} with $k$ and $k+1$ in place
of $i$~and~$j$. The hypotheses of Lemma~\ref{variousaltpaths} are that
$k\in\A(t)$, which holds since $k\in\D(u)\setminus\D(t)$, and
$k+1\in\SD(t)$, which is given here, and the additional hypothesis
for \ref{variousaltpaths}~\ref{lemma7.29_iii} is $\col_t(k+2)<\col_t(k)$, which is
also given. The conclusions are that $k\in\SD(v)$ and $k+1\notin\D(v)$,
and also that $k\notin\D(w)$ and $k+1\notin\D(w)$, where $w=s_kv$.

We can now check that $(c_{\alpha,w},c_{\alpha,v},c_{\alpha,t},c_{\beta,u})$ is
an alternating directed path of type $(k,k+1)$. It is certainly a path, since
$\mu(c_{\alpha,v},c_{\alpha,w})=\mu(c_{\alpha,t},c_{\alpha,v})=1$
by Corollary~\ref{arweightone}, and we have already seen that
$\mu(c_{\beta,u},c_{\alpha,t})\neq 0$. To show that the path is
alternating of type $(k,k+1)$ we must show that $k,\,k+1\notin \D(w)$,
that $k\in\D(v)$ and $k+1\notin\D(v)$, that $k\notin\D(t)$ and
$k+1\in\D(t)$, and that $k,\,k+1\in\D(u)$. These all appear explicitly
in the last paragraph above (given that $\SD(t)\subseteq\D(t)$ and
$\SD(v)\subseteq\D(v)$), apart from the statement that $k+1\in\D(u)$.
But this also holds, since $k+1\in\D(t)$ and $\D(t)\subsetneqq\D(u)$.

Admissibility of~$\Gamma$ tells us that all arc weights are positive,
and so it follows that $N^{3}_{k,k+1}(c_{\alpha,w},c_{\beta,u})>0$.
By the Polygon Rule, $N^{3}_{k+1,k}(c_{\alpha,w},c_{\beta,u})$ is also
nonzero. So there must exist $\xi\in\Lambda$ and
$(\delta,x)\in\mathcal{I}_{\xi}\times\STD(\xi)$, and
$\nu\in\Lambda$ and $(\gamma,y)\in\mathcal{I}_{\nu}\times\STD(\nu)$,
such that $(c_{\alpha,w}, c_{\delta,x}, c_{\gamma,y}, c_{\beta,u})$
is an alternating directed path of type $(k+1,k)$.

Lemma~\ref{xlexlessthant} applies, as in the first part above, with $k$
in place of~$i$ and $t_\lambda$ in place of~$t'$, but now with $j=k+1$.
The conditions $0\leqslant k\leqslant n-1$ and $k<j$ are still satisfied, and
we have $j\in\SD(t)$ and $v=s_jt$. The conditions
$\morethan {\morethan tk=(t_\lambda)}k$ and $k\in\A(t)$ still hold, as
before, and the extra condition needed for \ref{xlexlessthant}~\ref{lemma7.30_iii} is
$\col_t(k+2)<\col_t(k)$, which we also have. Hence $w<_{\lex}t_\lambda$.
Now by the definition of $t_\lambda$ it follows that
$w\notin\Ini_\Gamma(\lambda)$,
and since $(c_{\alpha,w}, c_{\delta,x})$ is an arc we conclude that
the molecules containing $c_{\alpha,w}$ and $c_{\delta,x}$ are of
the same type. Thus $\xi=\lambda$, and $x\in\STD(\lambda)$. Furthermore, by
Proposition~\ref{kldominancesamenolecule}, either $x<w$ or else
$\delta=\alpha$ and $x=s_lw>w$ for some $l\in[1,n-1]$. In the latter
case Lemma~\ref{dessjtminust} tells us that $\D(x)\setminus\D(w)=\{l\}$,
but the fact that $(c_{\alpha,w}, c_{\delta,x}, c_{\gamma,y}, c_{\beta,u})$ is
alternating of type $(k+1,k)$ means that $k+1\in\D(x)\setminus\D(w)$,
and it follows that $l=k+1$. So either $x<w$ or $x=s_{k+1}w>w$.

Since $\D(x)\cap\{k,k+1\}=\{k+1\}$ and $\D(y)\cap\{k,k+1\}=\{k\}$,
and since $\mu(c_{\gamma,y},c_{\delta,x})\neq 0$, it follows that
$\{c_{\delta,x},c_{\gamma,y}\}$ is a simple edge of~$\Gamma$, by
the Simplicity Rule. Thus $\nu=\lambda$ and $\gamma=\delta$, and
$y$ and $x$ are related by a dual Knuth move. Note that $y=k\neb(x)$
(see Definition~\ref{k-neighbourdefinition}). Note also that
$y\in\Ini_\Gamma(\lambda)$,
since $(c_{\gamma,y}, c_{\beta,u})$ is an arc whose head $c_{\gamma,y}$
is in a molecule of type~$\nu=\lambda$ and whose tail $c_{\beta,u}$
is in a molecule of type~$\pi\ne\lambda$.

If $x<w$ then it follows from Lemma~\ref{xlexlessthant}~\ref{lemma7.30_iv} that
$y<_{\lex}t_\lambda$, contradicting the fact that $t_\lambda$
is lexicographically minimal in $\Ini_\Gamma(\lambda)$.
So we must have $x=s_{k+1}w>w$. Since this says that $k+1\in\SA(w)$,
it follows from Lemma~\ref{variousaltpaths}~\ref{lemma7.29_iii} that
$s_ks_{k+1}w=k\neb(s_{k+1}w)$. That is, $s_kx=y$. But it
follows from Lemma~\ref{xlexlessthant}~\ref{lemma7.30_iii} that
$s_ks_{k+1}w<_{\lex}t_\lambda$. So $y<_{\lex}t_\lambda$,
contradicting the minimality of~$t_\lambda$, and completing the
proof that $\Morethan tk$ is $k$-critical.

The minimal approximate of $(u',t_\lambda)$ is by definition the
element $(u,t)\in A(u',t_\lambda)$ such that
$\lessthan uk=\lessthan tk=\tau_\kappa$, where
$\kappa=\shape(\lessthan uk)$. By the proof above, it
must also have the property that $\D(t)\subsetneqq\D(u)$ and
$\Morethan tk$ is $k$-critical, and by
Definition~\ref{k-minimalappr} this means $t$ is $k$-minimal
with respect to $u$.
\end{proof}

\begin{coro}\label{talphanystructure}
Let $\lambda,\,\pi\in\Lambda$ with $\Ini_\lambda(\Gamma)\ne\emptyset$ and
$\pi\in\Lambda\setminus\{\lambda\}$, and let $\alpha\in\mathcal I_\lambda$
and $(\beta,u')\in\mathcal I_\pi\times\STD(\pi)$ satisfy
$\mu(c_{\beta,u'},c_{\alpha,t_\lambda})\neq 0$.
Let $k$ be the restriction number of~$(u',t_\lambda)$.
Then $\morethan {(t_{\lambda})}{(k+1)}$ is minimal, and if $k+1\in\SD(t_{\lambda})$
then $\col_{t_{\lambda}}(k+1)=1+\col_{t_{\lambda}}(k+2)$.
\end{coro}

\begin{proof}
Let $(u,t)\in A(u',t_\lambda)$.
Then $\Morethan tk$ is $k$-critical, by Proposition~\ref{cellorder}. So
$\morethan t{(k+1)}$ is minimal, and if $k+1\in\SD(t)$ then
$\col_t(k+1)=1+\col_t(k+2)$. Since
$\morethan{(t_{\lambda})}k=\morethan tk$, the result follows.
\end{proof}

\begin{lemm}\label{1-minimalcomp}
Let $n\geq 2$, and let $\pi,\,\lambda\in P(n)$.
Let $t\in\STD(\lambda)$ and $u\in\STD(\pi)$ and suppose that $t$ is $1$-minimal with
respect to $u$. Then $\pi<\lambda$.
\end{lemm}

\begin{proof}
By Definition~\ref{k-minimalappr} the tableau $t$ is $1$-critical and the pair
$(u,t)$ is not $2$-restrictable. It follows that
$t(1,1)=u(1,1)=1$ and $t(1,2)=u(2,1)=2$. Thus $1\in\D(u)\setminus\D(t)$.
Furthermore, if $n=2$ then
$\pi=(2)$ and $\lambda=(1,1)$, giving $\pi<\lambda$ by
Definitions \ref{Dominance Order}~and~\ref{BruhatforC(n)}.

We proceed inductively on $n\geq 3$. If $\lambda=(1,1,\ldots,1)$ then it
follows readily from Definitions \ref{Dominance Order}~and~\ref{BruhatforC(n)}
that $\nu<\lambda$ holds for all $\nu\in P(n)\setminus\{\lambda\}$, and so
$\pi<\lambda$. So we may assume that $\lambda_1\geqslant 2$, and since $t$ is
$1$-critical we have $t(i,1)=i+1$ for all $i\in[2,\lambda_1]$. In particular,
$\col_t(3)=1<2=\col_t(2)$, and $\col_t(i+1)=\col_t(i)=1$ for all $i\in[3,\lambda_1]$.
It follows that $i\in\D(t)$ for all $i\in[2,\lambda_1]$. But Definition~\ref{k-minimalappr}
requires that $\D(t)\subsetneqq\D(u)$, and so for all $i\in[2,\lambda_1]$
we must have $\col_u(i+1)\leqslant\col_u(i)$. So all the numbers from $1$ to
$\lambda_1+1$ are in the first column of~$u$, and so $\pi_1\geqslant\lambda_1+1$.
In particular, $\pi\ne\lambda$.

Put $(u',t')=(\Lessthan u{(n-1},\Lessthan t{(n-1)})$, and let
$\sigma=\shape(u')$ and $\theta=\shape(t')$. Since $1$ is the restriction
number of~$(u,t)$ it is clear that $1$ is also the restriction number of~$(u',t')$,
and since $t$ is $1$-critical and $n\geqslant 3$ it is clear that $t'$
is $1$-critical. Since $\D(t)$ is a subset of $\D(u)$, and since
$\D(u')=\D(u)\setminus\{n-1\}$ and $\D(t')=\D(t)\setminus\{n-1\}$, it follows
that $\D(t')$ is a subset of~$\D(u')$. Hence $\D(t')\subsetneqq\D(u')$,
since $1\in\D(u')\setminus\D(t')$. So $t'$ is $1$-minimal with respect to
$u'$, whence $\sigma<\theta$ by the inductive hypothesis. In view of
Lemma~\ref{nsbeforent} it now suffices to show that
$\col_u(n)\leqslant\col_t(n)$. Since this obviously holds if
$\col_u(n)=1$, we may assume that $\col_u(n)>1$.

\looseness-1
Let $r=\max\{\,i\mid\theta_i>0\,\}$, the number of parts of~$\theta$, and note
that $\sum_{i=1}^r\theta_i=n-1$. Since $\sigma<\theta$ it follows that
$\sum_{i=1}^r\sigma_i\geqslant n-1$, and since $\sigma$ is a partition of $n-1$
this means that $\sigma$ has at most $r$ nonzero parts. In other words, all the numbers
from $1$~to~$n-1$ are in the first~$r$ columns of~$u$. Thus $\col_u(n)\leqslant r+1$.
We may now assume that $\col_t(n)\leqslant r$, since otherwise the
desired conclusion $\col_u(n)\leqslant\col_t(n)$ obviously holds.

Since $t$ and $t'$ are $1$-critical, $\morethan t2$ and $\morethan{(t')}2$ are
minimal of their respective shapes, by Definition~\ref{m-critical}. Thus
$\col_t(n-1)\leqslant\col_t(n)$, and $n-1$ is in the last column of $\morethan{(t')}2$.
Note that unless $\morethan{(t')}2$ has only one column, it has the same number of
columns as~$t'$, namely~$r$ columns.

Consider first the case that $\morethan{(t')}2$ has $r$ columns. Then
$r=\col_t(n-1)\leqslant\col_t(n)\leqslant r$, so that $\col_t(n)=\col_t(n-1)=r$,
and $n-1$ is a (weak) descent of~$t$. Since $\D(t)\subsetneqq\D(u)$ it follows that
$n-1\in\D(u)$, and so $\col_u(n)\leqslant\col_u(n-1)\leqslant r=\col_t(n)$,
as desired.

It remains to consider the case that $\morethan{(t')}2$ has only one column.
Then $\theta=(n-2,1)$, and $n-1=\theta_1+1\leqslant\lambda_1+1\leqslant\pi_1$.
Thus $\pi=(n-1,1)$, since we have assumed that $\col_u(n)\ne 1$.
So $\col_u(n-1)=1$ and $\col_u(n)=2$, which means that $n-1\notin\D(u)$.
So $n-1\notin\D(t)$, and so $\col_t(n)\geqslant\col_t(n-1)+1=2=\col_u(n)$,
as desired.
\end{proof}

\begin{lemm}\label{k-minimalonerow}
Let $\lambda,\,\pi\in\Lambda$ with $\Ini_\lambda(\Gamma)\ne\emptyset$ and
$\pi\in\Lambda\setminus\{\lambda\}$, and suppose that $\alpha\in\mathcal I_\lambda$
and $(\beta,u')\in\mathcal I_\pi\times\STD(\pi)$ satisfy
$\mu(c_{\beta,u'},c_{\alpha,t_\lambda})\neq 0$.
Let $k$ be the restriction number of~$(u',t_\lambda)$, and assume that
$k\geqslant 3$ and $(2,k-1)\notin[\pi]$. Then there is no $(u,t)\in A(u',t_\lambda)$
having the properties that $u(1,k)=n$ and $\col_t(n)=k-1$, and satisfying
\[
\Lessthan tk=\Lessthan uk\ =\ \vcenter{\offinterlineskip\small
\hbox{\vrule height 0.4 pt depth 0 pt width 108.4pt}
\hbox{\vrule height 9 pt depth 4 pt\hbox to 13pt{\hfil$1$\hfil}\vrule
\hbox to 13pt{\hfil$2$\hfil}\vrule\hbox to 30pt{\hfil$\cdots$\hfil}\vrule
\hbox to 25pt{\hfil$k-2$\hfil}\vrule
\hbox to 25pt{\hfil$k-1$\hfil}\vrule}
\hrule
\hbox{\vrule height 9 pt depth 4 pt\hbox to 13pt{\hfil$k$\hfil}\vrule}
\hrule width 14.2 pt}\ .
\]
\end{lemm}

\begin{proof}
Assume to the contrary that $(u,t)\in A(u',t_\lambda)$ has the stated properties.
By Remark~\ref{k-1subclassofk} the pair $(u,t)$ has the same restriction number
as $(u',t_\lambda)$, namely~$k$, and $A(u',t_\lambda)=C_{k-1}(u,t)$ is in bijective
correspondence with the set of standard tableaux of the same shape as $\lessthan tk$.
Since there is only one standard tableau of shape $(1^{k-1})$, it follows that
$A(u',t_\lambda)=\{(u,t)\}$.

It follows from Lemma~\ref{k-minimallemma} that $\D(t)\subsetneqq\D(u)$ and
$k=\min(\D(u)\setminus\D(t))$, and also that
$\mu(c_{\beta,u},c_{\alpha,t})=\mu(c_{\beta,u'},c_{\alpha,t_\lambda})\neq 0$. Furthermore,
$\Morethan tk$ is $k$-critical, by Lemma~\ref{cellorder}.
Hence $\col_t(k+1)=\col_t(k)+1$, and so $t(2,2)=k+1$. Since
$k\in\D(u)$ we have $\col_u(k+1)\leqslant\col_u(k)$. Hence $\col_u(k+1)=1$,
and it follows that $u(3,1)=k+1$. Thus $k+1\ne n$, since it is given that
$u(1,k)=n$.

\begin{Case}{1.} Suppose that $(u,t)=(u',t_\lambda)$.

Since $k\geqslant 3$, we have
$\col_u(k)=1<k-1=\col_u(k-1)$, and so
$k-1\in\SD(u)\subseteq\D(u)$. Let $v=s_{k-1}u$. Since
$k-1\in\SD(u)$, it follows that
$v\in\STD(\pi)$ and $k-1\notin\D(v)$.
Since $\col_u(k-2)=k-2<k-1=\col_u(k-1)$, it follows that
$k-2\notin\D(u)$. Moreover,
since $v$ is obtained from $u$ by switching the positions of $k-1$ and $k$,
and since $k\geqslant 3$, we have
$\col_v(k-1)=\col_u(k)=1\leqslant k-2=\col_u(k-2)=\col_v(k-2)$,
and so $k-2\in\D(v)$. So $v\to u$ is a DKM
(of the first kind) of index~$k-1$, and so
$\{c_{\beta,u},c_{\beta,v}\}$ is a simple edge in~$\Gamma$.

Since $k\geqslant 3$, we have $1\leqslant k-2=\col_t(k-2)<k-1=\col_t(k-1)$,
and it follows that $k-2\notin\D(t)$. Since $k\in\D(u)\setminus\D(t)$, we
also have $k\notin\D(t)$. Similarly, since $\Lessthan uk = \Lessthan tk$,
we have $k-2\notin\D(u)$, but since
$k\in\D(u)\setminus\D(t)$, we have $k\in\D(u)$.
We showed above that $k-2\in\D(v)$. Since
$\col_v(k+1)=\col_u(k+1)=1<k-1=\col_u(k-1)=\col_v(k)$, as
$v$ is obtained from $u$ by switching the positions of $k-1$ and $k$,
we also have $k\in\D(v)$.
Moreover, since $\mu(c_{\beta,u},c_{\alpha,t})\neq 0$ and
$\mu(c_{\beta,v},c_{\beta,u})=1$ (as
$\{c_{\beta,u},c_{\beta,v}\}$ is a simple edge),
it follows that $(c_{\alpha,t},c_{\beta,u},c_{\beta,v})$ is
an alternating directed path of type~$(k,k-2)$.
So $N^{2}_{k,k-2}(c_{\alpha,t},c_{\beta,v})>0$.

Since $\Gamma$ satisfies the
Polygon Rule,
$N^{2}_{k-2,k}(c_{\alpha,t},c_{\beta,v})=N^{2}_{k,k-2}(c_{\alpha,t},c_{\beta,v})\ne 0$,
and it follows that there are $\xi\in\Lambda$ and $(\gamma,x)\in\mathcal{I}_{\xi}\times\STD(\xi)$
such that $(c_{\alpha,t},c_{\gamma,x},c_{\beta,v})$
is an alternating directed path of type $(k-2,k)$.

Suppose that $k-1\notin\D(x)$. Since $\col_t(k-1)=k-1>1=\col_t(k)$ we see that
$k-1\in\D(t)$, and so $\D(t)\nsubseteq\D(x)$. Since $\mu(c_{\gamma,x},c_{\alpha,t})\ne0$,
it follows from the Simplicity Rule that $\{c_{\alpha,t},c_{\gamma,x}\}$ is a simple edge.
So $\xi=\lambda$ and $\gamma=\alpha$, and $x$ and $t$ are related by a~DKM.
By Lemma~\ref{dessjtminust} the index of this DKM is either the unique element
of $\D(x)\setminus\D(t)$ or the unique element of~$\D(t)\setminus\D(x)$. Now
$k-2\in\D(x)\setminus\D(t)$, since $(c_{\alpha,t},c_{\gamma,x},c_{\beta,v})$ is
alternating of type $(k-2,k)$, but it is not the case that $x=s_{k-2}t$ since
$s_{k-2}t$ is not standard. So the index of the DKM is the unique element of
$\D(t)\setminus\D(x)$, which must be~$k-1$ since $k-1\in\D(t)\setminus\D(x)$.
So $x=s_{k-1}t$,
and $\col_x(k+1)=\col_t(k+1)=2\leq k-1=\col_t(k-1)=\col_x(k)$. But this means that
$k\in\D(x)$, which is not allowed since $(c_{\alpha,t},c_{\gamma,x},c_{\beta,v})$ is
alternating of type $(k-2,k)$. So we must have $k-1\in\D(x)$.

Since $\mu(c_{\beta,v},c_{\gamma,x})\neq 0$, and since
$\D(x)\cap\{k-1,k\}=\{k-1\}$ and $\D(v)\cap\{k-1,k\}=\{k\}$, it follows from
the Simplicity Rule that $\{c_{\beta,v},c_{\gamma,x}\}$ is a
simple edge. Thus $\xi=\pi$ and $\gamma=\beta$, and $x$ and $v$
are related by a DKM. Furthermore, $x=(k-1)\neb(v)$
(see Definition~\ref{k-neighbourdefinition}). But since by definition
$v=s_{k-1}u$, and $\col_u(k+1)=\col_u(k)=1<k-1=\col_u(k-1)$, we see that
$\col_v(k+1)=\col_v(k-1)<\col_v(k)$, and so $(k-1)\neb(v)=s_kv$.
Therefore, $x=s_kv$.

Since $t=t_\lambda$ we have $\mu(c_{\beta,x},c_{\alpha,t_\lambda})\ne0$, and we
can apply Corollary~\ref{talphanystructure} with $x$ in place of~$u'$.
Since $x=s_ks_{k-1}u$ and $(u,t)$ is $k$-restricted, we see that $(x,t)$
is $(k-2)$-restricted, and so $k-2$ will replace the $k$ of~\ref{talphanystructure}.
Since $k-1\in\SD(t)$, it follows from Corollary~\ref{talphanystructure} that
$\col_t(k-1)=\col_t(k)+1$. But $\col_t(k-1)=k-1$ and $\col_t(k)=1$; hence
$k=3$. We were given that $\col_t(n)=k-1$; so $\col_t(n)=2$.
Recall also that $t(2,1)=k=3$ and $t(2,2)=k+1=4$. Since Corollary~\ref{talphanystructure}
also tells us that $\morethan t2$ is the minimal tableau of its shape,
it follows that $(3,1)\notin[\lambda]$ (since otherwise minimality of
$\morethan t2$ would require that $t(3,1)=t(2,1)+1=4$). So $\lambda_1=2$,
and so $\lambda_2\leqslant 2$.
Hence $(2,2)=t^{-1}(4)$ is the last box in the second column of~$[\lambda]$,
and since $\col_t(n)=2$ this forces $n=4$. This contradicts $n>k+1$,
which was proved above.
\end{Case}

\begin{Case}{2. } Suppose that $(u,t)\neq(u',t_\lambda)$.

Since $(u,t)\approx_k(u',t_\lambda)$ by Lemma~\ref{k-minimallemma1}, it follows that
$\Lessthan {(t_\lambda)}{k}=\Lessthan {(u')}k$ is a standard tableau of the same shape as
$\Lessthan {t}{k}=\Lessthan {u}k$. Hence there is an $i\in[1,k-2]$
such that $u'$ and $t_\lambda$ satisfy
\[
w'=\Lessthan {(t_\lambda)}{k}=\Lessthan {(u')}k\ =\ \vcenter{\offinterlineskip\small
\hbox{\vrule height 0.4 pt depth0 pt width 152.4pt}
\hbox{\vrule height 9 pt depth 4 pt\hbox to 25pt{\hfil$1$\hfil}\vrule
\hbox to 30pt{\hfil$\cdots$\hfil}\vrule\hbox to 13pt{\hfil$i$\hfil}\vrule
\hbox to 25pt{\hfil$i+2$\hfil}\vrule\hbox to 30pt{\hfil$\cdots$\hfil}\vrule\
\hbox to 25pt{\hfil$k$\hfil}\vrule}
\hrule
\hbox{\vrule height 9 pt depth 4 pt\hbox to 25pt{\hfil$i+1$\hfil}\vrule}
\hrule width 25 pt}\ .
\]
Furthermore, $\morethan {(t_\lambda)}k=\morethan tk$ and $\morethan {(u')}k=\morethan uk$
must also hold.

Recall that $t(2,2)=k+1$ and $u(3,1)=k+1$ have been proved above. Recall also that
$(2,k-1)\notin[\pi]$ is given, as are $u(1,k)=n$ and $\col_t(n)=k-1$.

Suppose first that $k=3$. Since $u^{-1}(n)$ is necessarily the last box
in its row and the last box in its column, and $u^{-1}(n)=(1,k)=(1,3)$,
it follows that $\pi_3=1$ and that $\pi$ has only three parts. And $\pi_2=1$, since
$(2,2)=(2,k-1)\notin[\pi]$. So $\pi=(n-2,1,1)$, the first row of $u$ is
$\vcenter{\offinterlineskip
\hrule
\hbox{\vrule height 8 pt depth 2 pt\hbox to 15pt{\hfil$1$\hfil}\vrule
\hbox to 15pt{\hfil$2$\hfil}\vrule
\hbox to 15pt{\hfil$n$\hfil}\vrule}
\hrule}\,$,
and $u(j,1)=j+1$ for all $j\in[2, n-2]$. Observe that $n-1\notin\D(u)$, since
$\col_u(n-1)=1<3=\col_u(n)$, and since $\D(t)\subsetneqq\D(u)$ it follows
that $n-1\notin\D(t)$. That is, $\col_t(n-1)<\col_t(n)=k-1=2$. Since $\morethan t4$ is
minimal of its shape, by Corollary~\ref{talphanystructure},
\[
1\leqslant \col_t(5)\leqslant\col_t(6)\leqslant\cdots\leqslant\col_t(n-1)<\col_t(n)=2,
\]
and so $\col_t(j)=1$ for all $j\in[5,n-1]$ and $\col_t(n)=2$. Since we also
have $t(2,2)=4$, it follows that $\lambda=(n-3,3)$, the first three rows of $t$ are
$\vcenter{\offinterlineskip
\hrule
\hbox{\vrule height 8 pt depth 2 pt\hbox to 15pt{\hfil$1$\hfil}\vrule
\hbox to 15pt{\hfil$2$\hfil}\vrule}
\hrule}\,$,
$\vcenter{\offinterlineskip
\hrule
\hbox{\vrule height 8 pt depth 2 pt\hbox to 15pt{\hfil$3$\hfil}\vrule
\hbox to 15pt{\hfil$4$\hfil}\vrule}
\hrule}\,$
and
$\vcenter{\offinterlineskip
\hrule
\hbox{\vrule height 8 pt depth 2 pt\hbox to 15pt{\hfil$5$\hfil}\vrule
\hbox to 15pt{\hfil$n$\hfil}\vrule}
\hrule}\,$, and $t(j,1)=j+2$ for $j\in [4, n-3]$. Now since there are only
two standard tableaux of shape~$(2,1)$ we see that $C_k(u,t)=\{(u,t),(s_2u,s_2t)\}$,
and hence $(u',t_\lambda)=(s_2u,s_2t)$. But it is easily checked that
$\D(s_2u)=\D(s_2t)=\{1\}\cup[3,n-2]$, and it follows that
$\mu(c_{\beta,s_2u},c_{\alpha,s_2t})=0$. Since this contradicts the hypothesis
that $\mu(c_{\beta,u'},c_{\alpha,t_\lambda})\neq 0$, it follows that $k\geqslant 4$.

Since $\Lessthan {(t_\lambda)}{k}=\Lessthan {(u')}k$, we have
$\D(t_\lambda)\cap[1,k-1]=\D(u')\cap[1,k-1]$. Since $k\geqslant 4$ and
$\morethan{(t_\lambda)}k=\morethan tk$, we have
$\col_{t_\lambda}(k+1)=\col_t(k+1)=2<k-1=\col_t(k-1)=\col_{t_\lambda}(k)$, and similarly
we have $\col_{u'}(k+1)=\col_u(k+1)=1<k-1=\col_u(k-1)=\col_{u'}(k)$.
Thus $k\in\D(t_\lambda)\cap\D(u')$, and so
$\D(t_\lambda)\cap[1,k]=\D(u')\cap[1,k]$. So every element $l\in\D(u')\setminus\D(t_\lambda)$
satisfies $l>k$. But $\D(t_\lambda)\subsetneqq\D(u')$, since
$\mu(c_{\beta,u'},c_{\alpha,t_\lambda})\neq 0$ and $\alpha\ne\beta$, and so
it follows that
$\D(\morethan{(t_\lambda)}k)\subsetneqq\D(\morethan{(u')}k)$.

Suppose that $i>1$. Then
$\col_{w'}(i+1)=1\leqslant i-1=\col_{w'}(i-1)<\col_{w'}(i-1)+1=\col_{w'}(i)$,
and so $i\in\SD(w')\subseteq\D(w')$ and $i-1\notin\D(w')$. Thus
$\D(w')\cap\{i-1,i\}=\{i\}$. It follows from $i\in\SD(w')$, or
by inspection, that $s_iw'$ is standard and $i\notin\D(s_iw')$.
Furthermore, since $\col_{s_iw'}(i-1)=\col_{w'}(i)>\col_{w'}(i+1)=\col_{s_iw'}(i)$, it
follows that $i-1\in\D(s_iw')$, and hence $\D(s_iw')\cap\{i-1,i\}=\{i-1\}$.
Thus $s_iw'\to^{*1} w'$, a DKM of index $i\leqslant k-1$.
Since the same DKM
takes $(s_iu',s_it_\lambda)$ to $(u',t_\lambda)$, we have
$(s_iu',s_it_\lambda)\approx_k(u',t_\lambda)$. It follows by Lemma~\ref{noname0} that
$\mu(c_{\beta,s_iu'},c_{\alpha,s_it_\lambda})=\mu(c_{\beta,u'},c_{\alpha,t_\lambda})\neq 0$.
Since $\pi\ne\lambda$ the vertices $c_{\alpha,s_it_\lambda}$ and $c_{\beta,s_iu'}$
lie in molecules of different types, and so it follows that
$(\alpha,s_it_\lambda)\in\Ini_{\lambda}(\Gamma)$, and
$s_it_\lambda\in\Ini_\Gamma(\lambda)$. But it follows from
Remark~\ref{equalsubsetsS} that $s_it_\lambda<t_\lambda$, and hence
$s_it_\lambda<_{\lex} t_\lambda$ by Corollary~\ref{domimplieslextableau}.
Since this contradicts the minimality of~$t_\lambda$, we conclude that $i=1$.

We make use of $\Gamma_L$, the $W_L$-restriction of $\Gamma$,
using the notation of Remark~\ref{restrictionrem} and Eq.~\eqref{verticesofGammaL}.
Let $v=\jdt(\morethan{(u')}{1})$ and $x=\jdt(\morethan{(t_\lambda)}{1})$, and write
$\rho=\shape(v)$ and $\xi=\shape(x)$.

By Remark~\ref{restrictionrem}, we can identify the vertex $c_{\beta,u'}$ of
$\Gamma_L$ with $c''_{\delta,v}$ for some $\delta\in\mathcal{I}_{L,\beta,\pi,\rho}$, and
the vertex $c_{\alpha,t_\lambda}$ of $\Gamma_L$ with $c''_{\gamma,x}$ for some
$\gamma\in\mathcal{I}_{L,\alpha,\lambda,\xi}$. Note that
$\mathcal{I}_{L,\beta,\pi,\xi}\cap\mathcal{I}_{L,\alpha,\lambda,\rho}=\emptyset$, since
$\pi\neq\lambda$, and so it follows that $\delta\neq\gamma$.
Now $\D(v)\nsubseteq\D(x)$, since $\D(\morethan{(u')}{1})\nsubseteq\D(\morethan{(t_\lambda)}{1})$,
and so it follows that $\mu_L(c''_{\delta,v},c''_{\gamma,x})=\mu(c_{\beta,u'},c_{\alpha,t_\lambda})\neq 0$.
Since $\Gamma_L$ is ordered, and since
$\delta\neq\gamma$, we deduce that $v<x$. Thus $\rho\leqslant\xi$.

Observe that $\col_{t_\lambda}(j)=\col_{w'}(j)\leqslant k-1$ for all $j\in[1,k]$. Since
$\morethan {(t_\lambda)}k=\morethan tk$, we have $\col_{t_\lambda}(k+1)=2$ and
$\col_{t_\lambda}(n)=k-1$. Furthermore, $\morethan{t_\lambda}{(k+1)}$ is the minimal
tableau of its shape, by Corollary~\ref{talphanystructure}, and so
$\col_{t_\lambda}(k+2)\leqslant\col_{t_\lambda}(k+3)\leqslant\cdots\leqslant\col_{t_\lambda}(n)$.
It follows that $\col_{t_\lambda}(j)\leqslant\col_{t_\lambda}(n)=k-1$ for all $j\in[1,n]$,
showing that the partition~$\pi$ has exactly $k-1$ parts.
Since $u(1,k)=n$ the partition $\pi$ has exactly $k$ parts, and since also $(2,k-1)\notin\pi$
it follows that $\pi_{k-1}=\pi_k=1$. Now let $(g,p)$ and $(h,r)$ be the boxes vacated in
$\jdt(\morethan{(u')}{1})$ and $\jdt(\morethan{(t_\lambda)}{1})$ respectively. Recalling that
$u'(2,1)=t_\lambda(2,1)=2$, since $i=1$, we see that the box $(2,1)$ is in both slide paths,
and so $g\geqslant 2$ and $h\geqslant 2$. In particular, $(g,p)\ne(1,k)$, and so
$\rho_k=\pi_k=1$. On the other hand, $\xi$ has only $k-1$ parts, since $\lambda$ has
only $k-1$ parts. Hence we find that
$\sum_{m=1}^{k-1}\xi_m=n-1=\sum_{m=1}^{k}\rho_m=1+\sum_{m=1}^{k-1}\rho_m>\sum_{m=1}^{k-1}\rho_m$,
contradicting $\rho\leqslant\xi$.\qedhere
\end{Case}
\end{proof}

\begin{lemm}\label{kminimimalcompared}
Let $\lambda,\,\pi\in\Lambda$ with $\Ini_\lambda(\Gamma)\ne\emptyset$ and
$\pi\in\Lambda\setminus\{\lambda\}$, and suppose that $\alpha\in\mathcal I_\lambda$
and $(\beta,u')\in\mathcal I_\pi\times\STD(\pi)$ satisfy
$\mu(c_{\beta,u'},c_{\alpha,t_\lambda})\neq 0$.
Then $\pi<\lambda$.
\end{lemm}

\begin{proof}
Clearly $n$ is at least $2$. Since the vertices $c_{\beta,u'}$ and $c_{\alpha,t_\lambda}$
belong to different molecules we must have $\D(t_\lambda)\subsetneqq\D(u')$
and $\mu(c_{\alpha,t_\lambda},c_{\beta,u'})=0$. Let $k$ be the restriction number of
$(u',t_\lambda)$, noting that $1\leqslant k\leqslant n-1$ since $\pi\ne\lambda$.
We write $\nu=\shape(\Lessthan{(u')}k) = \shape(\Lessthan{(t_\lambda)}k)$, and
let $(u,t)$ be an arbitrary element of $A(u',t_\lambda)$ (which is nonempty,
by Lemma~\ref{k-minimallemma}).

Lemma~\ref{k-minimallemma1} shows that $(u,t)$ is $k$-restricted
and $(u,t)\approx_k (u',t_\lambda)$, and Proposition~\ref{cellorder}
shows that $\Morethan{t}{k}$ is $k$-critical.
Furthermore, we have $\D(t)\subsetneqq\D(u)$ and $k=\min(\D(u)\setminus\D(t))$, and
$\mu(c_{\beta,u},c_{\alpha,t})=\mu(c_{\beta,u'},c_{\alpha,t_\lambda})\neq 0$,
by Lemma~\ref{k-minimallemma}.
Since $\Morethan tk$ is~$k$-critical, it follows from Definition~\ref{m-critical}
that the following all hold:
\begin{gather}
\col_t(k)=\min\{\col_t(i)\mid i\in[k,n]\,\},\label{kfirst}\\
\col_t(k+1)=\col_t(k)+1,\label{colk+1}\\
\col_t(k+2)\leqslant\col_t(k+3)\leqslant\cdots\leqslant\col_t(n-1)\leqslant\col_t(n).
\label{incseq}
\end{gather}
Note that $\col_t(i)=\col_{t_\lambda}(i)$ for all $i\in[k+1,n]$, since
$\morethan tk=\morethan{(t_\lambda)}k$. Let $l=\col_t(n)=\col_{t_\lambda}(n)$.

If $k=1$ then since $t$ is $1$-minimal with respect to $u$, it follows
by Lemma~\ref{1-minimalcomp} that $\pi<\lambda$, while if $k=n-1$ then we must have
$\D(u)=\{n-1\}\cup\D(t)$, and $\pi<\lambda$ follows by Lemma~\ref{nminus1order0}.
So may assume that $1<k<n-1$.

We make use of $\Gamma_K$ and $\Gamma_L$, the restrictions of
$\Gamma$ to $W_K$~and~$W_L$, using the notation from Remark~\ref{restrictionrem}
and Eqs~\eqref{verticesofGammaK}~and~\eqref{verticesofGammaL} above.
Looking at $\Gamma_K$ first, let $v=\Lessthan u{(n-1)}$ and
$x=\Lessthan t{(n-1)}$, and let $\sigma=\shape(v)$ and $\theta=\shape(x)$.
Note that $\shape(\morethan vk)=\shape(\morethan{(\Lessthan u{n-1})}k)=\sigma/\nu$,
and similarly $\shape(\morethan xk)=\theta/\nu$.
So $\sigma/\nu\vdash (n-1-k)$ and $\theta/\nu\vdash (n-1-k)$.

By Eq.~\eqref{verticesofGammaK} we have $c_{\beta,u}=c'_{\delta,v}$ and
$c_{\alpha,t}=c'_{\gamma,x}$ for some $\delta\in\mathcal{I}_{K,\sigma}$ and
$\gamma\in\mathcal{I}_{K,\theta}$. Since $c_{\beta,u}$ and $c_{\alpha,t}$
lie in different molecules of~$\Gamma$ they lie in different
$K$-submolecules; that is, $\delta\ne\gamma$.
By the definition of the $W_K$-restriction of~$\Gamma$ we
have $\D(v)\setminus\D(x)=(\D(u)\setminus\D(t))\cap[1,n-2]$, and
so we see that $k\in\D(v)\setminus\D(x)$. Hence
$\mu_K(c'_{\delta,v},c'_{\gamma,x})=\mu(c_{\beta,u},c_{\alpha,t})\neq 0$.
Since $\Gamma_K$ is ordered and $\gamma\neq\delta$
it follows that $v<x$, and hence $\sigma\leqslant\theta$,
by Definition~\ref{ExtDominance Order Tableaux}.

If $\col_u(n)\leqslant\col_t(n)$ then, since $\sigma\leqslant \theta$,
Lemma~\ref{nsbeforent} yields $\pi\leqslant\lambda$. Since
$\pi\neq\lambda$, the desired conclusion $\pi<\lambda$ holds in this case.
So we may assume henceforth that $\col_u(n)>\col_t(n)=l$.

Observe that $\col_t(k)\leqslant l=\col_t(n)$, by Eq.~\eqref{kfirst}. Suppose
first that $\col_t(k)=l$. Then it follows from Eqs \eqref{kfirst}~and~\eqref{incseq} that
$l\leqslant\col_t(k+2)\leqslant\col_t(k+3)\leqslant\cdots\leqslant\col_t(n-1)\leqslant l$,
and so $\col_t(i)=l$ for all $i\in[k+2,n]$. Hence $i\in\WD(t)$ for all $i\in[k+2,n-1]$,
and $k+1\in\SD(t)$ since $\col_t(k+1)=l+1$ by Eq.~\eqref{colk+1}.
Since $\D(t)\subsetneqq\D(u)$ and also $k\in\D(u)\setminus\D(t)$, it follows that
$i\in\D(u)$ for all $i\in[k,n-1]$. Thus
$\col_u(n)\leqslant\col_u(n-1)\leqslant\cdots\leqslant\col_u(k)$. But
since $(u,t)$ is $k$-restricted, $\col_u(k)=\col_t(k)=l$, and
we deduce that $\col_u(n)\leqslant l=\col_t(n)$. Since we have already shown that
this gives $\pi<\lambda$, for the rest of the proof we may (and do) assume that
\begin{equation}\label{columninequalities}
\col_t(k)<\col_t(n)=l<\col_u(n).
\end{equation}
In fact we shall show that Eq.~\eqref{columninequalities} leads to a contradiction,
by showing that $(u,t)$ satisfies the conditions of Lemma~\ref{k-minimalonerow}.

Since $\sigma\leqslant\theta$ (equivalent to $\sigma\trianglerighteq\theta$ by Definition~\ref{BruhatforC(n)}) we have
\begin{equation}\label{etalethetaeq}
\sum_{m=1}^l\theta_m\leqslant\sum_{m=1}^l\sigma_m.
\end{equation}
It follows from Eqs \eqref{colk+1}~and~\eqref{columninequalities} that $\col_t(k+1)\leqslant l$,
and by Eq.~\eqref{incseq} we deduce that
\begin{equation}\label{coltileql}
[k+1,n-1]\subseteq\{\,i\mid\col_t(i)\leqslant l\,\}.
\end{equation}
So all elements of $[k+1,n-1]$ appear in the first $l$ columns of
the tableau $\morethan{(\Lessthan t{n-1})}k=\morethan xk$, and since
$\morethan xk\in\STD_k(\theta/\nu)$ it follows that $\sum_{i=1}^l(\theta_m-\nu_m)=n-1-k$.
So Eq.~\eqref{etalethetaeq} can be expressed in the form
\begin{equation*}
n-1-k + \sum_{m=1}^l\nu_m \leqslant\sum_{m=1}^l (\sigma_m-\nu_m) + \sum_{m=1}^l\nu_m.
\end{equation*}
Thus $\sum_{m=1}^l (\sigma_m-\nu_m)\geqslant n-1-k$. Now since
$\sigma/\nu\vdash(n-1-k)$ it follows that $\sigma_m-\nu_m=0$ for
all $m>l$ and $\sum_{m=1}^l (\sigma_m-\nu_m)= n-1-k$.
Since $\morethan{(\Lessthan u{n-1})}k\in\STD_k(\sigma/\nu)$ it follows that
\begin{equation}\label{colsileql}
[k+1,n-1]\subseteq\{\,i\mid\col_u(i)\leqslant l\,\}.
\end{equation}
Eq.~\eqref{coltileql} and $\col_t(n)=l$ combined give
$[k+1,n]\subseteq\{\,i\mid\col_t(i)\leqslant l\,\}$, and we see that
column~$l$ is the last nonempty column of $\morethan tk$. Since
$\shape(\morethan tk)=\lambda/\nu\vdash (n-k)$,
it follows that $\sum_{m=1}^l(\lambda_m-\nu_m)=n-k$. Now
$\shape(\morethan uk)=\pi/\nu\vdash (n-k)$, but although
the first $l$ columns of $\morethan{(\Lessthan u{n-1})}k$ include
all elements of $[k+1,n-1]$ (by Eq.~\eqref{colsileql}) they do not include
all elements of $[k+1,n]$, since $\col_u(n)>l$ (by Eq.~\eqref{columninequalities}).
Thus $[\pi/\nu]$ is the disjoint union of $[\sigma/\nu]$ and $\{u^{-1}(n)\}$, and
$\sum_{m=1}^l(\pi_m-\nu_m)=n-k-1$. We deduce that
\begin{equation}\label{smuonemorethansnu}
\sum_{m=1}^l\pi_m = n-k-1+\sum_{m=1}^l\nu_m=
\left(\sum_{m=1}^l\lambda_m\right)-1.
\end{equation}

Looking now at $\Gamma_L$, we put $w=\jdt(\morethan u1)$ and $y=\jdt(\morethan t1)$,
and define $\rho=\shape(w)$ and $\xi=\shape (y)$.
By Eq.~\eqref{verticesofGammaL} we have $c_{\beta,u}=c''_{\zeta,w}$ and
$c_{\alpha,t}=c''_{\epsilon,y}$ for some $\zeta\in\mathcal{I}_{L,\rho}$ and
$\epsilon\in\mathcal{I}_{L,\xi}$. Since $c_{\beta,u}$ and $c_{\alpha,t}$
lie in different molecules of~$\Gamma$ they must also lie in different
$L$-submolecules; that is, $\zeta\ne\epsilon$. By
Proposition~\ref{glambdapiconnectedcomps} and the definition of the $W_L$-restriction
of~$\Gamma$ we have $\D(w)\setminus\D(y)=\D(\morethan u1)\setminus\D(\morethan t1)=
(\D(u))\setminus\D(t))\cap[2,n-1]$, and so it follows that $k\in\D(w)\setminus\D(y)$.
Hence $\mu_L(c''_{\zeta,w},c''_{\epsilon,y})=\mu(c_{\beta,u},c_{\alpha,t})\neq 0$.
Since $\Gamma_L$ is ordered and $\zeta\neq\epsilon$
it follows that $w<y$, and hence $\rho\leqslant\xi$,
by Definition~\ref{ExtDominance Order Tableaux}.

Let $(g,p)$ and $(h,r)$ be the boxes vacated by $\jdt((1,1),\morethan u1)$
and $\jdt((1,1),\morethan t1)$. Note that $\rho_p=\pi_p-1$ and $\rho_m=\pi_m$
for all $m\ne p$. Similarly, $\xi_r=\lambda_r-1$ and $\xi_m=\lambda_m$ for
all $m\ne r$. We claim that
\begin{equation}\label{qleqcoltnlp}
r\leqslant l<p.
\end{equation}
If $p\leqslant l$ then
$\sum_{m=1}^l \rho_m =\bigl(\sum_{m=1}^l \pi_m\bigr)-1<\sum_{m=1}^l \pi_m$,
and since $\bigl(\sum_{m=1}^l \lambda_m\bigr) - 1\leqslant\sum_{m=1}^l\xi_m$,
it follows by Eq.~(\ref{smuonemorethansnu}) that
$\sum_{m=1}^l \rho_m <\sum_{m=1}^l\xi_m$, which contradicts $\rho\leqslant\xi$.
Similarly, if $l<r$ then $\sum_{m=1}^l \xi_m = \sum_{m=1}^l \lambda_m$,
and since $\sum_{m=1}^l \pi_m \geqslant \sum_{m=1}^l \rho_m$, it again follows by
Eq.(\ref{smuonemorethansnu}) that $\sum_{m=1}^l \xi_m > \sum_{m=1}^l \rho_m$,
contradicting $\rho\leqslant\xi$. Hence $r\leqslant l<p$, as claimed.

We shall now show that $u(g,p)=n$.
If $k+1\leqslant u(g,p)<n$ then $p=\col_u(u(g,p))\leqslant l$ by Eq.~(\ref{colsileql}),
contradicting Eq.~(\ref{qleqcoltnlp}). Now suppose that $u(g,p)\leqslant k$. Then
clearly $u(b)\leqslant k$ for all boxes $b$ in the slide path of $\jdt((1,1),\morethan u1)$
and hence the slide path of $\jdt((1,1),\morethan u1)$ coincides with the
slide path of $\jdt((1,1),\morethan{(\Lessthan uk)}1)$. Furthermore, since
$\Lessthan uk=\Lessthan tk$, and the slide path of $\jdt((1,1),\morethan t1)$
extends (or equals) that of $\jdt((1,1),\morethan{(\Lessthan tk)}1)$, it follows that
$(g,p)$ is in the slide path of $\jdt((1,1),\morethan t1)$. By Lemma~\ref{slidepath}
it follows that $h\geqslant g$ and $r\geqslant p$. Since the latter inequality
contradicts Eq.(~\ref{qleqcoltnlp}) we conclude that $u(g,p)<n$ is impossible, and
so $u(g,p)=n$, as claimed.

We claim that
\begin{equation}\label{colsiltpminus1}
[k+1,n-1] \subseteq \{\,i\mid\col_u(i)<p-1\,\}.
\end{equation}
By Eqs~\eqref{colsileql}~and~\eqref{columninequalities} we see that
$\col_u(n-1)<\col_u(n)$. Thus $n-1\in\A(u)$, and hence $n-1\in\A(t)$, since
$\D(t)\subsetneqq\D(u)$. So $\col_t(n-1)<\col_t(n)=l$.
Note that, by Eq.~(\ref{coltileql}), we have $\col_t(k+1)\leq l$.

Suppose first that $\col_t(k+1)=l$. Then $\col_t(n-1)<\col_t(k+1)=\col_t(k)+1$,
by Eq.~\eqref{colk+1}, and since Eq.~\eqref{kfirst} gives $\col_t(k)\leqslant\col_t(n-1)$
it follows that $\col_t(n-1)=\col_t(k)$. By Eqs \eqref{kfirst}~and~\eqref{incseq}
we deduce that $\col_t(k)=\col_t(k+2)=\col_t(k+3)=\cdots=\col_t(n-1)$.
Thus $[k+2,n-2]\subseteq\WD(t)$. Moreover, $k+1\in\SD(t)$, since
$\col_t(k+1)>\col_t(k)=\col_t(k+2)$. Hence $[k+1,n-2]\subseteq\D(t)$. Since
$\D(t)\subsetneqq\D(u)$, and we
also have $k=\min(\D(u)\setminus\D(t))$, it follows that $[k,n-2]\subseteq\D(u)$.
Thus $\col_u(n-1)\leqslant\col_u(n-2)\leqslant\cdots\leqslant\col_u(k+1)\leqslant\col_u(k)$.
But $\col_u(k)=\col_t(k)=l-1$, since $(u,t)$ is $k$-restricted and
$l=\col_t(k+1)=\col_t(k)+1$.
So $\col_u(k)<p-1$ by Eq.~\eqref{qleqcoltnlp}, and so Eq.~\eqref{colsiltpminus1}
holds in this case.

Suppose now that $\col_t(k+1)<l$. Since we have shown above that $\col_t(n-1)<l$,
and since we also have $\col_t(k+2)\leqslant\col_t(k+3)\leqslant\cdots\leqslant\col_t(n-1)$ by
Eq.~\eqref{incseq}, it follows that $[k+1,n-1]\subseteq\{\,i\mid\col_t(i)\leqslant l-1\,\}$
(a strengthening of Eq.~\eqref{coltileql}). So all elements of $[k+1,n-1]$ appear in
the first $l-1$ columns of $\morethan{(\Lessthan t{(n-1)})}k$, and so
$\sum_{m=1}^{l-1}(\theta_m-\nu_m)=n-k-1$. Since $\sigma\leqslant\theta$, we have
$\sum_{m=1}^{l-1}\theta_m\leqslant\sum_{m=1}^{l-1}\sigma_m$, which can be written as
\begin{equation*}
n-1-k + \sum_{m=1}^{l-1}\nu_m \leqslant\sum_{m=1}^{l-1} (\sigma_m-\nu_m) + \sum_{m=1}^{l-1}\nu_m,
\end{equation*}
giving $n-1-k\leqslant\sum_{m=1}^{l-1} (\sigma_m-\nu_m)$. Since
$\sigma/\nu\vdash(n-1-k)$ it follows that $\sigma_m-\nu_m=0$ for
all $m>l-1$ and $\sum_{m=1}^{l-1} (\sigma_m-\nu_m)= n-1-k$. Since
$\morethan{(\Lessthan u{(n-1)})}k\in\STD_k(\sigma/\nu)$ it follows that
$[k+1,n-1]\subseteq\{\,i\mid\col_u(i)\leqslant l-1\,\}$,
and since $l<p$ by Eq.~\eqref{qleqcoltnlp} it follows that
Eq.~\eqref{colsiltpminus1} also holds in this case. This completes the proof of our claim.

Recall that $(g,p)$ is vacated by $\jdt((1,1),\morethan u1)$ and that $u(g,p)=n$.
Let $b$ be the box that $n$ slides into, so that either $b=(g-1,p)$ or
$b=(g,p-1)$. Obviously $u(b)\notin\{\,i\mid\col_u(i)<p-1\}$, and so
Eq.~(\ref{colsiltpminus1}) shows that $u(b)\notin[k+1,n]$. So $b$
is in the diagram of $\shape(\Lessthan uk)$, and must be the box vacated by
$\jdt((1,1),\morethan{(\Lessthan uk)}1)$. The slide path of
$\jdt((1,1),\morethan u1)$ is the slide path of
$\jdt((1,1),\morethan{(\Lessthan uk)}1)$ extended by the additional box $(g,p)$.
Since $\Lessthan uk=\Lessthan tk$, the slide path of $\jdt((1,1),\morethan t1)$
also extends the slide path of $\jdt((1,1),\morethan{(\Lessthan uk)}1)$, and,
in particular, includes~$b$. If $b=(g-1,p)$ then Lemma~\ref{slidepath}
gives $p\leqslant r$, contradicting Eq.~(\ref{qleqcoltnlp}). So $b=(g,p-1)$,
and Lemma~\ref{slidepath} gives $p-1\leqslant r$.
But since $r\leqslant l<p$ by Eq.~(\ref{qleqcoltnlp}),
this shows that
\begin{equation}\label{columnequalities}
p-1=r=l=\col_t(n).
\end{equation}

Recall that $(u,t)$ was chosen as an arbitrary element of $A(u',t_\lambda)$, and so
all that we have proved thus far applies for all elements of $A(u',t_\lambda)$.
However, since $n>k$ and $\morethan uk=\morethan{(u')}k$ for all
$(u,t)\in A(u',t_\lambda)$, it follows that $(g,p)=u^{-1}(n)=(u')^{-1}(n)$
is independent of the choice of~$(u,t)$. So $(g,p-1)$ in independent of the
choice of~$(u,t)$. Since we have just shown that $(g,p-1)$ is the box vacated
by $\jdt((1,1),\morethan{(\Lessthan uk)}1)$, this must be true for all choices
of~$(u,t)$. Furthermore, $(g,p-1)\ne u^{-1}(k)$, by Eqs
\eqref{columnequalities}~and~\eqref{columninequalities}, and so $(g,p-1)$ is in the
diagram of $\shape(\lessthan uk)$, and is the box vacated by
$\jdt((1,1),\morethan{(\lessthan uk)}1)$, for all choices of $(u,t)$.

Recall from Remark~\ref{k-1subclassofk} that there exists $\kappa\in P(k-1)$
and a bijection $A(u',t_\lambda)\to\STD(\kappa)$ given by
$(u,t)\mapsto w=\lessthan uk=\lessthan tk$. So the previous paragraph
tells us that for every $w\in\STD(\kappa)$ the jeu-de-taquin slide
$\jdt((1,1),\morethan w1)$ vacates the box $(g,p-1)$. Considering in particular
the cases $w=\tau_\kappa$ and $w=\tau^\kappa$ and applying Lemma~\ref{jofminandmaxt},
we deduce that $\kappa$ has exactly $p-1$ parts, all equal to~$g$.

Observe that since $\kappa=(g^{p-1})$ there are only two $\kappa$-addable boxes,
namely $(g+1,1)$ and $(1,p)$. But since
$[\shape(\Lessthan tk)]=\{t^{-1}(k)\}\cup[\shape(\lessthan tk)]=\{t^{-1}(k)\}\cup[\kappa]$,
it is clear that $t^{-1}(k)$ is $\kappa$-addable. And $t^{-1}(k)\ne(1,p)$, by
Eq.~\eqref{columninequalities} and the fact that $\col_u(n)=p$. So $t^{-1}(k)=(g+1,1)$,
and since $(u,t)$ is $k$-restricted it follows that $u^{-1}(k)=(g+1,1)$ also.

Since $\shape(\lessthan uk)=(g^{p-1})$ it follows that $\col_u(i)<p$ for
all $i\in[1,k-1]$. But we have just shown that $\col_u(k)=1$, and
Eq.~\eqref{colsiltpminus1} gives $\col_u(i)<p-1$ for all $i\in[k+1,n-1]$.
So $n$ is the only number in column~$p$ of~$u$, and so clearly it
must be in the first row. That is, $g=1$.

Since $\kappa=\shape(\lessthan tk)$ is a partition of $k-1$, and since $\kappa=(1^{p-1})$,
it follows that $k=p$. It follows from Eq.~\eqref{columninequalities} that
$1=\col_t(k)<\col_t(n)<\col_u(n)=k$, and so $k\geqslant 3$. Moreover,
since $\row_u(i)=1$ for $i\in[1,k-1]\cup\{n\}$, while $\col_u(k)=1$
and $\col_u(i)<k-1$ for $i\in[k+1,n-1]$ (by Eq.~\eqref{colsiltpminus1}),
we see that $(2,k-1)\notin [\pi]$. We have now shown that $(u,t)$
satisfies all the conditions in Lemma~\ref{k-minimalonerow}, contradicting
$(u,t)\notin A(u',t')$, and completing our proof.
\end{proof}

\begin{lemm}\label{samelambda}
Suppose that $\Gamma$ is strongly connected. Then $\Lambda$ has only one element.
\end{lemm}

\begin{proof}
Assume to the contrary that $\Lambda$ has more than one element, and
choose $\lambda\in\Lambda$ to be minimal, in the sense that there is no
$\pi\in\Lambda\setminus\{\lambda\}$ such that $\pi<\lambda$. Note that
$\Ini_{\lambda}(\Gamma)\neq\emptyset$, since
$\Gamma$ is strongly connected and $\Lambda\setminus\{\lambda\}\ne\emptyset$.
Hence there exists $\pi\in\Lambda\setminus\{\lambda\}$ such that
$\mu(c_{\beta,u},c_{\alpha,t_{\lambda}})\neq 0$,
for some $(\beta,u)\in\mathcal{I}_{\pi}\times\STD(\pi)$ and
$\alpha\in\mathcal{I}_{\lambda}$. But now Lemma~\ref{kminimimalcompared}
gives $\pi<\lambda$, contradicting the choice of $\lambda$.
\end{proof}
Lemma~\ref{samelambda} says that if $\Gamma$ is an admissible $W_n$-cell
then all molecules of $\Gamma$ have the same type~$\lambda$, for some
$\lambda\in P(n)$. In general whenever $\Lambda=\{\lambda\}$ we may
call $\lambda$ the type of~$\Gamma$.

\begin{lemm}\label{muandnucells}
Suppose that $n\geqslant 2$. Let $D$ and $D'$ be cells of $\Gamma$,
of types $\lambda$ and $\pi$ respectively. Let $\leqslant_\Gamma$
be the partial order on the set of cells of $\Gamma$ (as defined in
Section~\ref{sec:3} above).
Then $D'\leqslant_{\Gamma} D$ implies $\pi\leqslant\lambda$. In particular,
$\pi\leqslant\lambda$ holds if there exist vertices $c\in D$ and $c'\in D'$
such that $\mu(c',c) \neq 0$.
\end{lemm}

\begin{proof}
We can assume that $\pi\neq\lambda$, since the result is trivial otherwise.
Write $\mathcal{C}$ for the set of cells of $\Gamma$, and observe that
$|\mathcal{C}|\geqslant 2$, since $\pi\neq\lambda$ implies that $D\ne D'$.

Suppose first that $\mathcal{C}=\{D,D'\}$. Since $D'\leqslant_{\Gamma} D$
it is immediate from the definition of the partial order $\leqslant_{\Gamma}$
that there exist $c\in D$ and $c'\in D'$ with $\mu(c',c) \neq 0$, and
since $\pi\ne\lambda$ it follows that $\Ini_\Gamma(\lambda)\ne\emptyset$.
So $t_\lambda$ exists, and by the definition of $\Ini_\Gamma(\lambda)$
there exist $\alpha\in\mathcal{I}_\lambda$ and
$(\beta,u)\in\mathcal{I}_\pi\times\STD(\pi)$ with
$\mu(c_{\beta,u},c_{\alpha,t_\lambda})\ne 0$.
So Lemma~\ref{kminimimalcompared} gives $\pi<\lambda$, as required.

Proceeding inductively, suppose now that $|\mathcal{C}|>2$ and that the result
holds for all admissible $W_n$-graphs with fewer than $|\mathcal{C}|$ cells.
Let $C_0,\,C_1\in\mathcal{C}$ be such that $C_0$ is minimal and $C_1$
is maximal in the partial order~$\leqslant_{\Gamma}$. Then $C_0$ and
$C\setminus C_1$ are subsets of $C$ that are closed in the sense defined
in Section~\ref{sec:3}. Writing $\Gamma_0$ and $\Gamma_1$ for the full
subgraphs of $\Gamma$ spanned by $C\setminus C_0$ and $C\setminus C_1$,
we see that $\Gamma_0$ and $\Gamma_1$ are admissible $W_n$-graphs, with
arc weights and vertex colours inherited from $\Gamma$, and with cells
that are cells of~$\Gamma$. It follows that if $C_0$ and $C_1$ can be
chosen so that $D$ and $D'$ are both contained in $C\setminus C_0$ or both
contained in $C\setminus C_1$ then the result follows from the inductive
hypothesis. Since $D'\leqslant_{\Gamma} D$ by assumption, it remains to
consider the possibility that $C_0=D'$ is the unique minimal cell and
$C_1=D$ is the unique maximal cell.

Since $|\mathcal{C}|>2$ we can choose $C'\in\mathcal{C}\setminus\{C_0,C_1\}$.
By Lemma~\ref{samelambda} the cell $C'$ has type $\{\nu\}$ for some
$\nu\in\Lambda$. Moreover, $C_0\leqslant_{\Gamma} C'$ since $C_0$ is the
unique minimal cell, and $C'\leqslant_{\Gamma}C_1$ since $C_1$ is the
unique maximal cell. Since $D'$ and $C'$ are cells of $\Gamma_1$, we have
$\pi\leqslant\nu$ by the inductive hypothesis, and since
$C'$ and $D$ are cells of $\Gamma_0$, we have $\nu\leqslant\lambda$
by the inductive hypothesis. So it follows that $\pi\leq\lambda$, as required.

Recall that if $\mu(c',c)\ne0$ then it follows as a consequence that
$\tau(c')\nsubseteq\tau(c)$ (since $\Gamma$ is reduced, in the
terminology of Section~\ref{sec:3}). So, by the definition of~$\leqslant_{\Gamma}$,
if $\mu(c',c)\ne0$ then $c'\leqslant_{\Gamma}c$, and so the last
assertion of the lemma follows from the rest.
\end{proof}

We are now able to complete the proof of Theorem~\ref{inductivestep}.

\begin{proof}
Suppose that $(\alpha,t)\in\mathcal{I}_{\lambda}\times\STD(\lambda)$ and
$(\beta,u)\in\mathcal{I}_{\pi}$ satisfy
$\mu(c_{\beta,u},c_{\alpha,t}) \neq 0$. It follows from
Lemma~\ref{muandnucells} that $\pi\leqslant\lambda$. Now
Proposition~\ref{kldominancesamenolecule} says that
$u < t$ unless $\alpha=\beta$ and $u=s_i t>t$ for
some $i\in[1,n-1]$. That is, $\Gamma$ is ordered.
\end{proof}

\begin{rema}\label{orderedregular}
In particular, it follows from Theorem~\ref{inductivestep} that the
Kazhdan--Lusztig $W_n$-graph corresponding to the regular representation
of~$\mathcal{H}(W_n)$ is ordered in the sense of Definition~\ref{orderedWgraphdef}.
In this case the vertex set of $\Gamma=(C,\mu,\tau)$ is $C=W_n$, the set
of molecule types is $\Lambda=P(n)$, for each $\lambda\in P(n)$ the set of
molecules of type $\lambda$ is indexed by~$\mathcal{I}_\lambda=\STD(\lambda)$,
and for each $\lambda\in\Lambda$ and $x\in\mathcal{I}_\lambda$ the
set $C_{x,\lambda}$ consists of those $w\in W_n$ such that
$\RSQ(w)=x$, where $\RSQ(w)$ is the recording tableau in the Robinson--Schensted
process.

Now let $y,\,w\in W_n$ and put $\RS(w)=(t,x)\in\STD(\lambda)^2$ and
$\RS(y)=(u,v)\in\STD(\pi)^2$, where $\lambda,\,\pi\in P(n)$.
The conclusion of Theorem~\ref{inductivestep}, applied in this case, is that
if $\mu(y,w)\ne 0$ and $\tau(y)\nsubseteq\tau(w)$ then either $u<t$ or
else $\pi=\lambda$ and $(u,v)=(st,x)$ for some $s\in S_n$.

If $\Gamma$ is replaced by $\Gamma\opp=(C,\mu,\tau\opp)$, then since
$\RS(w^{-1})=(x,t)$ and $\RS(y^{-1})=(v,u)$ by Theorem~\ref{rswinv},
the conclusion of Theorem~\ref{inductivestep} is that if $\mu(y,w)\ne 0$ and
$\tau\opp(y)\nsubseteq\tau\opp(w)$ then either $v<x$ or
else $\pi=\lambda$ and $(u,v)=(t,sx)$ for some $s\in S_n$.

Thus, in particular, if $\mu(y,w)\neq 0$ and $\tau(y)\nsubseteq\tau(w)$ or
$\tau\opp(y)\nsubseteq\tau\opp(w)$ then $\pi\leq\lambda$.

It follows from the definition of the preorder~$\preceq\leftright$ (in
Section~\ref{sec:3} above) that if $y\preceq\leftright w$
then there is a sequence of elements
 $y=z_0,\, z_1,\, \ldots,\, z_{m-1},\, z_m=w$ such that $\mu(z_{i-1},z_i)\ne 0$
and $\bar\tau(z_{i-1})\nsubseteq\bar\tau(z_i)$ for
each $i\in[1,m]$. Since $\bar\tau(z_{i-1})\nsubseteq\bar \tau(z_i)$ is equivalent
to $\tau(z_{i-1})\nsubseteq\tau(z_i)$ or $\tau\opp(z_{i-1})\nsubseteq\tau\opp(z_i)$,
it follows that $\pi\leqslant\lambda$ whenever $y\preceq\leftright w$.
\end{rema}

\begin{rema}\label{orderedtwosided}
Let $y,w \in W_n$, and let $RS(y)=(u,v)$ and $RS(w)=(t,x)$. Remark~\ref{orderedregular}
says that if $y\preceq\leftright w$ then
$\pi\leqslant\lambda$, where $\pi=\shape(x)=\shape(u)$ and
$\lambda=\shape(y)=\shape(v)$. This gives an alternative approach to
the ``only if'' part of the following well-known result. (See, for example,
\cite[Theorem~5.1]{geck:klMurphy}.)
\begin{theo}\label{Wn-two-sidedcell}
Let $y,w\in W_n$ and $\pi,\lambda\in P(n)$, and suppose that $\RS(y)\in\STD(\pi)\times\STD(\pi)$ and
$\RS(w)\in\STD(\lambda)\times\STD(\lambda)$. Then
$y\preceq\leftright w$ if and only if $\pi\leqslant\lambda$. In particular,
the sets $D(\lambda)=\{w\in W_n \mid \RS(w)\in\STD(\lambda)^2\}$, where
$\lambda\in P(n)$, are precisely the Kazhdan--Lusztig two-sided cells.
\end{theo}
Let $\lambda\in P(n)$. Recall that for each $t\in\STD(\lambda)$ the set
$C(t)=\{w\in W_n\mid \RSQ(w)=t\}$ is a left cell, and $\Gamma(C(t))$ is
isomorphic to $\Gamma_{\lambda}$ (defined in \ref{std-tableau-W-graph} above).
Thus $C(t)$ is of type~$\lambda$, by Remark~\ref{Gamma-lambdaIsAdmissible},
and the set $D(\lambda)=\bigsqcup_{t\in\STD(\lambda)}C(t)$
is a union of $\lvert \STD(\lambda) \rvert$
left cells of type~$\lambda$.
\end{rema}

\section{\texorpdfstring{$W$}{W}-graphs for admissible cells in type \texorpdfstring{$A$}{A}}
\label{sec:9}

\begin{defi}\label{probablepair}
Let $\lambda\in P(n)$. A pair of standard $\lambda$-tableaux $(u,t)$ is said to
be a \emph{probable pair} if $u<t$ and $\D(t)\subsetneqq\D(u)$.
\end{defi}
It is readily checked that there are no probable pairs in $P(n)$ unless $n\geqslant 5$.
Note that if $(u,t)$ is a probable pair then $u\ne t$, and so the set
$F(u,t)$ is defined and nonempty.

Recall Definition~\ref{k-reduced}: if $\pi,\,\lambda\in P(n)$ then a pair
$(u,t)\in\STD(\pi)\times\STD(\lambda)$ is said to be favourable if the restriction number
of $(u,t)$ lies in $\D(u)\oplus\D(t)$.
\begin{lemm}\label{rnltmaxsd}
Let $\lambda\in P(n)$ and $u,t\in\STD(\lambda)$, and let $i$ be the
restriction number of~$(u,t)$. Suppose that $(u,t)$ is favourable and
$\D(t)\subsetneqq\D(u)$. Then there exists $j\in \SD(t)$ such that $j>i$.
\end{lemm}

\begin{proof}
Suppose that there is no such~$j$, so that
$\D(\morethan ti)\cap [i+1,n-1]=\WD(\morethan ti)\cap [i+1,n-1]$.
Thus $\morethan ti$ is the minimal tableau of its shape, by Remark~\ref{minimal-tableau},
and so $\col_t(k)\geqslant\col_t(i+1)$ for all $k\in[i+1,n]$.
Furthermore, $\D(u)\oplus\D(t)=\D(u)\setminus\D(t)$, since $\D(t)\subsetneqq\D(u)$,
and so $i\in\D(u)\setminus\D(t)$, since $(u,t)$ is favourable.
Thus $\col_t(i+1)>\col_t(i)$, since $i\notin\D(t)$, and so
$\col_t(k)>\col_t(i)$ for all $k>i$. In other words, for all
$m\in[1,\col_t(i)]$, column~$m$ of $t$ is entirely filled by numbers
from the set $[1,i]$. But now since $u$~and~$t$ have the same shape and
$\Lessthan ui=\Lessthan ti$, it follows that the same holds for~$u$: for all
$m\in[1,\col_u(i)]$, column~$m$ of $u$ is entirely filled by numbers
from the set $[1,i]$. In particular, $\col_u(i+1)>\col_u(i)$, contradicting $i\in\D(u)$.
\end{proof}

\begin{lemm}\label{existencef}
Let $\lambda\in P(n)$ and $u,t\in\STD(\lambda)$ with $\D(t)\subsetneqq\D(u)$,
and suppose that $(u,t)$ is favourable.
Let $i$ be the restriction number of~$(u,t)$, and suppose that $i+1=\max(\SD(t))$. Then
$\col_{t}(i+2)\neq\col_t(i)$.
\end{lemm}

\begin{proof}
Suppose to the contrary that $\col_{t}(i+2)=\col_t(i)$. We have that
$\Lessthan ui=\Lessthan ti$, since $(u,t)$ is $i$-restricted.
Since $i+1=\max(\SD(t))$, it follows
$\SD(\morethan t{(i+1)})=\emptyset$, and thus $\morethan t{(i+1)}$ is
minimal, by Remark~\ref{minimal-tableau}. Hence
$\col_t(i)=\col_t(i+2)\leqslant\col_t(k)$ for all $k>i+2$.
Furthermore, $i\in\D(u)\setminus\D(t)$, since $(u,t)$ is favourable,
and so $\col_t(i)<\col_t(i+1)$. So $\col_t(k)\geqslant\col_t(i)$ for
all $k>i$. Now since $u$ and $t$ are of the same shape and
$\Lessthan ui=\Lessthan ti$, we deduce that $\col_u(k)\geqslant\col_u(i)$
for all $k>i$. In particular, $\col_u(i+1)\geqslant\col_u(i)$.
But $\col_u(i+1)\leqslant\col_u(i)$, since $i\in\D(u)\setminus\D(t)$,
and so $\col_u(i+1)=\col_u(i)$.

Let $(g,p)=t^{-1}(i)=u^{-1}(i)$. Since $i+1$ is in the same column of~$u$
as~$i$, it follows that $u(g+1,p)=i+1$. Since $i+2$ is in the same column
of $t$ as $i$, and $i+1$ is in a different column, it follows that
$t(g+1,p)=i+2$. Now let $m$ be the maximal positive integer such that
$\col_u(i+l)=p$ for all $l\in[1,m]$, so that
$u(g+l,p)=i+l$ for all $l\in[1,m]$. Since $t$ and $u$ have the
same shape and $\morethan t{(i+1)}$ is minimal, we see that
$t(g+l,p)=i+l+1$ for all $l\in[1,m]$.
If $m>1$ then $\col_t(i+m+1)=p=\col_t(i+m)$,
and it follows that $i+m\in\WD(t)$, while if $m=1$ then $i+m=i+1\in\SD(t)$
(since $i+1=\max(\SD(t))$ is given). In either case it follows that
$i+m\in\D(u)$, since $\D(t)\subsetneqq\D(u)$. So
$\col_u(i+m+1)\leqslant\col_u(i+m)=p=\col_u(i)$, whence $\col_u(i+m+1)=p$,
since it was shown above that $\col_u(k)\geqslant\col_u(i)$ for
all $k>i$. But this contradicts the choice of~$m$.
\end{proof}

Suppose that $\Gamma = \Gamma(C,\mu,\tau)$ is an admissible $W_n$-graph
whose molecules are all of type $\lambda$, for some $\lambda\in P(n)$,
and let $\mathcal{I}$ index the molecules.
By Remark~\ref{vertexsetofWngraph}, the vertex set of $\Gamma$
can be written as $C=\bigsqcup_{\alpha\in\mathcal{I}}C_{\alpha}$,
where for each $\alpha\in\mathcal{I}$ the set
$C_{\alpha}=\{c_{\alpha,t}\mid t\in\STD(\lambda)\}$ spans a molecule,
$\tau(c_{\alpha,t})=\{\,s_j\mid j\in\D(t)\,\}$ for all $\alpha\in\mathcal{I}$ and $t\in\STD(\lambda)$,
and the simple edges of $\Gamma$ are the pairs $\{c_{\beta,u},c_{\alpha,t}\}$ such that
$\alpha=\beta$ and $u$ and $t$ are related by a DKM.

\begin{lemm}\label{linktominimalCor}
Let $u,t\in\STD(\lambda)$, and suppose that $(u,t)$ is a probable pair.
Let $(v,x)$ be an arbitrary element of the set $F(u,t)$. Then $(v,x)$ is a
probable pair, $\,\max(\SD(x))=\max(\SD(t))$, and
$\mu(c_{\beta,v},c_{\alpha,x})=\mu(c_{\beta,u},c_{\alpha,t})$.
\end{lemm}

\begin{proof}
\looseness-1
Note that $t$ is not minimal, since $u<t$. Hence $\SD(t)\ne\emptyset$.
Let $i$ be the restriction number of $(u,t)$. Recall that, by the
way $F(u,t)$ is defined, $(v,x)$ is $i$-restricted and favourable.

Suppose first that $(u,t)$ is favourable. Then $i<\max(\SD(t))$, by
Lemma~\ref{rnltmaxsd}, and so $\max(\SD(t))=\max(\SD(\morethan ti))$. But
$\max(\SD(\morethan ti))=\max(\SD(\morethan xi))$,
since $\morethan ti=\morethan xi$, and so
$\max(\SD(t))=\max(\SD(x))$. Moreover,
$\col_u(i+1)\leqslant\col_u(i)=\col_t(i)<\col_t(i+1)$, since
$i\in\D(u)\oplus\D(t)=\D(u)\setminus\D(t)$, and hence
$\col_v(i+1)\leqslant\col_v(i)=\col_x(i)<\col_x(i+1)$.
Thus $i\in\D(v)\setminus\D(x)$. Clearly $\D(v)\cap[1,i-1]=\D(x)\cap[1,i-1]$,
since $\Lessthan vi=\Lessthan xi$, and since $\D(v)\cap[i+1,n-1]=\D(u)\cap[i+1,n-1]$
and $\D(x)\cap[i+1,n-1]=\D(t)\cap[i+1,n-1]$, it follows from
$\D(t)\subsetneqq\D(u)$ that $\D(x)\subsetneqq\D(v)$. Since $u<t$,
Proposition~\ref{bruhatorderpreserved} gives $v<x$, and it follows that
$(v,x)$ is a probable pair. Finally, by Lemma~\ref{noname0}
applied with $j=\max(\SD(t))$, we see that
$\mu(c_{\beta,v},c_{\alpha,x})=\mu(c_{\beta,u},c_{\alpha,t})$,
as required.

Now suppose that $(u,t)$ is not favourable. Again it
follows from Proposition~\ref{bruhatorderpreserved} that $v<x$.
Moreover, $u<t$ gives $\Lessthan u{(i+1)}\leqslant\Lessthan t{(i+1)}$,
so that $\Lessthan u{(i+1)} < \Lessthan t{(i+1)}$ (since
$\Lessthan u{(i+1)}\neq\Lessthan t{(i+1)}$). It follows by
Remark~\ref{resdomrem} that $\col_u(i+1)<\col_t(i+1)$.
Since $\D(u)\oplus\D(t)=\D(u)\setminus\D(t)$,
we have $i<\min(\D(u)\setminus\D(t))$ by Remark~\ref{k-reducedrem}.
Thus $(u,t)$ satisfies the hypothesis of Lemma~\ref{noname2},
and the conclusion is that
$\D(v)\setminus\D(x)=\{i\}\cup(\D(u)\setminus\D(t))$ and $\D(x)\setminus\D(v)=\emptyset$.
Since $\D(x)\setminus\D(v)=\emptyset$ while $\D(v)\setminus\D(x)\neq\emptyset$,
we have $\D(x)\subsetneqq\D(v)$. Hence $(v,x)$ is probable.

Let $j=\max(\SD(x))$, and note that $j>i$ by Lemma~\ref{rnltmaxsd},
Thus $j=\max(\SD(\morethan xi))$, and since $\morethan ti = \morethan xi$,
it follows that $j=\max(\SD(\morethan ti))=\max(\SD(t))$.
Finally, since it was shown above that $i<\min(\D(u)\setminus\D(t))$,
Lemma~\ref{noname0} again gives
$\mu(c_{\beta,v},c_{\alpha,x})=\mu(c_{\beta,u},c_{\alpha,t})$.
\end{proof}

\begin{prop}\label{monomolecularadmcellsareKL}
Monomolecular admissible cells of type $A_{n-1}$ are Kazhdan--Lusztig.
\end{prop}

\begin{proof}
Suppose that $\Gamma=\Gamma(C,\mu,\tau)$ is a monomolecular admissible $W_n\,$-cell.
Then there is a partition $\lambda$ of $n$ such that
$C=\{c_t\mid t\in\STD(\lambda)\}$, and
$\{c_u,c_t\}$ is a simple edge of $\Gamma$
if and only if $u,\,t\in\STD(\lambda)$ are related by a DKM.
In view of Corollary~\ref{leftcelllambda}, our task is to show that
$\Gamma\cong \Gamma_\lambda=\Gamma(\STD(\lambda),\mu^{(\lambda)},\tau^{(\lambda)})$.
Recall that $\Gamma_\lambda$ is an admissible $W_n$-graph consisting
of a single molecule of type~$\lambda$ (by Remark~\ref{Gamma-lambdaIsAdmissible}).
Clearly $t\mapsto c_t$ is a bijection from the vertex set of $\Gamma_\lambda$
to the vertex set of $\Gamma$.
Since it follows from Remark~\ref{vertexsetofWngraph} that
$\tau(c_t) = \{\,s_j\mid j\in\D(t)\,\} = \tau^{(\lambda)}(t)$ for all $t\in\STD(\lambda)$, it remains
to show that $\mu(c_u,c_t)=\mu^{(\lambda)}(u,t)$ for all $u,\,t\in\STD(\lambda)$.
Note that, by Theorem~\ref{combinatorialCharacterisation}, both $\Gamma$ and
$\Gamma_\lambda$ satisfy the
Compatibility Rule, the Simplicity Rule,
the Bonding Rule and the Polygon Rule.

We have shown in Theorem~\ref{inductivestep} that $\Gamma$ and
$\Gamma_\lambda$ are both ordered. Thus if $u,t\in\STD(\lambda)$
then $\mu(c_u,c_t)=\mu^{(\lambda)}(u,t)=0$
unless $u<t$ or $u=s_it>t$ for some $i\in[1,n-1]$.
If $u=s_it>t$ for some $i\in[1,n-1]$ then we have
$\mu(c_u,c_t)=\mu^{(\lambda)}(u,t)=1$ by Corollary~\ref{arweightone}.
Now suppose that $u<t$ and $\D(t)\nsubseteq\D(u)$. If one
or other of $\mu(c_u,c_t)$ and $\mu^{(\lambda)}(u,t)$ is nonzero then,
by the Simplicity Rule, one or other of
$\{c_u,c_t\}$ and $\{u,t\}$ is a simple edge, whence $u$ and $t$ are
related by a DKM
(by Remark~\ref{vertexsetofWngraph}), and
both $\{c_u,c_t\}$ and $\{u,t\}$ are simple edges. So
$\mu(c_u,c_t)=\mu^{(\lambda)}(u,t)=1$ in this case. Obviously there
is nothing to show if $\mu(c_u,c_t)$ are $\mu^{(\lambda)}(u,t)$ both zero,
and so all that remains is to show that $\mu(c_u,c_t)=\mu^{(\lambda)}(u,t)$
whenever $u<t$ and $\D(t)\subsetneqq\D(u)$. That is, it remains to show
that $\mu(c_u,c_t)=\mu^{(\lambda)}(u,t)$ for all probable pairs of
standard $\lambda$-tableaux. In other words we must show that for all
$t\in\STD(\lambda)$ we have $\mu(c_u,c_t)=\mu^{(\lambda)}(u,t)$ whenever
$u\in\STD(\lambda)$ and $(u,t)$ is a probable pair.

If $t=\tau_{\lambda}$ then there is nothing to prove. Proceeding inductively on
the lexicographic order, let $t'\in\STD(\lambda)\setminus\{\tau_\lambda\}$, and assume
that the result holds for all $t\in\STD(\lambda)$ such that $t<_{\lex} t'$.

Suppose that $(u',t')$ is a probable pair. Let $i$ be the restriction number
of $(u',t')$, and let $j=\max(\SD(t'))$. Let $(u,t)\in F(u',t')$, and note
that $(u,t)$ is $i$-restricted and favourable, and satisfies
$\morethan ti=\morethan {t'}i$
and $\morethan ui=\morethan {u'}i$. Moreover, Lemma~\ref{linktominimalCor} shows that
$(u,t)$ is probable, $\mu(c_{u},c_{t})=\mu(c_{u'},c_{t'})$ and
$\mu^{(\lambda)}(u,t)=\mu^{(\lambda)}(u',t')$, and $\max(\SD(t))=\max(\SD(t'))=j$.
Thus the conclusion $\mu(c_{u'},c_{t'})=\mu^{(\lambda)}(u',t')$ will
follow if we can show that $\mu(c_{u},c_{t})=\mu^{(\lambda)}(u,t)$.

Note that $i\in\A(t)$ and $i<j$, by Lemma~\ref{rnltmaxsd}. Let $v=s_jt$. Then
$v\in\STD(\lambda)$ and $v<t$, since $j\in\SD(t)$, and
$v<_{\lex}t'$ by Lemma~\ref{xlexlessthant}.
Since $(u,t)$ is favourable, and
since $\D(t)\subsetneqq\D(u)$ (since $(u,t)$ is probable),
we have $i\in\D(u)\setminus\D(t)$ and $j\in D(u)\cap\D(t)$. That is,
$i\notin\D(t)$ and $j\in\D(t)$, and $i,\,j\in\D(u)$.
Since $i<j$ we have either $j-i>1$ or $j-i=1$.

\begin{Case}{1.} Suppose that $j-i>1$, so that $m(s_i,s_j)=2$.
Lemma~\ref{variousaltpaths}~\ref{lemma7.29_i} tells us that $i,\,j\notin\D(v)$.
We set $X=\{\,x\in \STD(\lambda)\mid i\in\D(x)\text{ and }j\notin\D(x)\,\}$ and
$Y=\{\,y\in \STD(\lambda)\mid j\in\D(y)\text{ and }i\notin\D(y)\,\}$.

If $(c_v,c_{y_1},c_{u})$ is any alternating directed path of type $(j,i)$, then, since
$\Gamma$ is ordered, it follows that either $y_1=s_jv=t>v$ or $y_1<v$. Similarly, if
$(c_v,c_{x_1},c_{u})$ is any alternating directed path of type $(i,j)$, then
it follows that either $x_1=s_iv>v$ or $x_1<v$. Note that
if $x_1=s_iv>v$, then since $x_1\in\STD(\lambda)$, it follows that
$i\in\SA(v)$. Thus, if $x_1=s_iv>v$, then
$i\in\D(s_iv)$ and $j\notin\D(s_iv)$ by Lemma~\ref{variousaltpaths}~\ref{lemma7.29_i}. That is,
$s_iv\in X$. Now since $\Gamma$ satisfies the
Polygon Rule,
we have
$N^{2}_{j,i}(c_{v},c_{u})=N^{2}_{i,j}(c_{v},c_{u})$, and it follows that
\begin{multline}\label{sumijorder2}
\mu(c_{u},c_{t})\mu(c_{t},c_v)\,\, +
\!\!\!\!\sum_{y_1\in Y,\ y_1<v}\!\!\!\!\mu(c_{u},c_{y_{1}})\mu(c_{y_{1}},c_v)\\[-5pt]
=\mu(c_{u},c_{s_iv})\mu(c_{s_iv},c_v)\,\, +
\!\!\!\!\sum_{x_1\in X,\ x_1<v}\!\!\!\!\mu(c_{u},c_{x_1})\mu(c_{x_1},c_v),
\end{multline}
where the term $\mu(c_{u},c_{s_iv})\mu(c_{s_iv},c_v)$ on the right hand side of
Eq.~\eqref{sumijorder2} should be omitted if $i\notin\SA(v)$. Note that if $i\in\SA(v)$
then $(c_v,c_{s_iv},c_{u})$ is not necessarily a directed path, since there need not
be an arc from $s_iv$ to~$u$, but in this case
$\mu(c_{u},c_{s_iv})\mu(c_{s_iv},c_v)=0$ since $\mu(c_{u},c_{s_iv})=0$. Similarly,
$(c_v,c_{t},c_{u})$ is not necessarily a directed path, since there need not be an arc
from $t$ to~$u$, but $\mu(c_{u},c_{t})\mu(c_{t},c_v)=0$ in this case.
So Eq.~\eqref{sumijorder2} still holds in these cases.

Since Corollary~\ref{arweightone} gives $\mu(c_{t},c_v)= 1$, and $\mu(c_{s_iv},c_v)= 1$
if $i\in\SA(v)$,
Eq.~\eqref{sumijorder2} yields the following formula for $\mu(c_{u},c_{t})$:
\begin{equation*}
\mu(c_{u},c_{t}) =\mu(c_{u},c_{s_iv})\,\, +
\!\!\!\sum_{x_1\in X,\ x_1<v}\!\!\!\mu(c_{u},c_{x_1})\mu(c_{x_1},c_v)\,\,-
\!\!\!\sum_{y_1\in Y,\ y_1<v}\!\!\!\mu(c_{u},c_{y_1})\mu(c_{y_1},c_v),
\end{equation*}
where $\mu(c_{u},c_{s_iv})$ should be interpreted as 0 if $s_iv\notin\STD(\lambda)$.

Working similarly on $\Gamma_{\lambda}$ yields the following formula for
$\mu^{(\lambda)}(u,t)$:
\begin{equation*}
\mu^{(\lambda)}(u,t) \Mk=\Mk\mu^{(\lambda)}(u,s_iv)
+\!\!\!\sum_{x_1\in X,\ x_1<v}\!\!\!\mu^{(\lambda)}(u,x_1)\mu^{(\lambda)}(x_1,v)
-\!\!\!\sum_{y_1\in Y,\ y_1<v}\!\!\!\mu^{(\lambda)}(u,y_1)\mu^{(\lambda)}(y_1,v).
\end{equation*}

Now $v<_{\lex}t'$ by Lemma~\ref{xlexlessthant} and $s_iv<_{\lex}t'$ (if $i\in\SA(v)$)
by Lemma~\ref{xlexlessthant}~\ref{lemma7.30_i}. Moreover, all $x_1$ and $y_1$ appearing above
satisfy $x_1<_{\lex}t'$ and $y_1 <_{\lex} t'$, by Lemma~\ref{xlexlessthant}~\ref{lemma7.30_ii}. Hence
it follows by the inductive hypothesis that the corresponding arc weights
in the two formulae above are the same. Thus $\mu(c_{u},c_{t})=\mu^{(\lambda)}(u,t)$,
and $\mu(c_{u'},c_{t'})=\mu^{(\lambda)}(u',t')$, as desired.
\end{Case}

\begin{Case}{2.} Suppose that $i=j-1$, so that $m(s_i,s_j)=3$.
By Lemma~\ref{existencef}, $\col_{t}(j-1)\neq\col_{t}(j+1)$, and so either
$\col_{t}(j-1)<\col_{t}(j+1)$ or $\col_{t}(j-1)>\col_{t}(j+1)$.

If $\col_{t}(j-1)<\col_{t}(j+1)$, then the result follows by the same argument as
used in Case~1, with $j-1$ replacing $i$ and using Lemma~\ref{variousaltpaths}~\ref{lemma7.29_ii}
in place of Lemma~\ref{variousaltpaths}~\ref{lemma7.29_i}.

Suppose that $\col_{t}(j-1)>\col_{t}(j+1)$. Since
$j-1\in\SD(v)$ by Lemma~\ref{variousaltpaths}~\ref{lemma7.29_iii}, we have
$s_{j-1}v\in\STD(\lambda)$ and $s_{j-1}v<v$. Let
$w=s_{j-1}v$. It follows by Lemma~\ref{variousaltpaths}~\ref{lemma7.29_iii} that
$j-1,\,j\notin\D(w)$, whereas
$j-1,\,j\in\D(u)$ (as we have already seen).

We consider directed paths from $c_w$ to~$c_{u}$ that have length three
and are alternating of type $(j-1,j)$ or type $(j,j-1)$.
We have $j\in\D(t)$ and $j-1\notin\D(t)$ (as we have already seen),
while $j-1\in\D(v)$ and $j\notin\D(v)$ by Lemma~\ref{variousaltpaths}~\ref{lemma7.29_iii}.

If $(c_w,c_{x_1},c_{x_2},c_{u})$ is any alternating directed path of
type~$(j-1,j)$, then, since $\Gamma$ is ordered, it follows that
either $x_1=s_{j-1}w=v>w$, or else $x_1<w$. Moreover, since $\Gamma$ satisfies the
Simply-Laced Bonding Rule,
the fact that
$j-1\in\D(x_1)$ and $j\notin\D(x_1)$
shows that $c_{x_{2}}$ is the unique vertex adjacent to $c_{x_{1}}$ satisfying $j-1\notin\D(x_2)$
and $j\in\D(x_2)$. That is, $x_2$ is the $(j-1)$-neighbour of $x_1$. Thus it follows that
either $x_1=v$ and $x_2=s_{j}v=t$, or else $x_1<w$ and
either $x_2=s_jx_1>x_1$ or $x_2=s_{j-1}x_1<x_1$.

Similarly, if $(c_{w},c_{y_1},c_{y_2},c_{u})$ is any alternating directed path
of type~$(j,j-1)$, then it follows that either
$y_1=s_jw>w$ or $y_1<w$, and $y_2$ is the $(j-1)$-neighbour of~$y_1$.
Note that if $y_1=s_jw>w$, then since $y_1\in\STD(\lambda)$,
it follows that $j\in\SA(w)$. Thus, if $y_1=s_jw>w$ then Lemma~\ref{variousaltpaths}~\ref{lemma7.29_iii}
tells us that $y_2=s_{j-1}y_1=s_{j-1}s_jw>s_jw=y_1$, and also that
$j\in\D(s_jw)$ and $j-1\notin\D(s_jw)$, and that
$j-1\in\D(s_{j-1}s_jw)$ and $j\notin\D(s_{j-1}s_jw)$. On the other hand
if $y_1<w$ then either $y_2=s_{j-1}y_1>y_1$ or $y_2=s_jy_1<y_1$.

Now since $\Gamma$ satisfies the
Polygon Rule,
we have
$N^{3}_{j-1,j}(c_{w},c_{u})=N^{3}_{j,j-1}(c_{w},c_{u})$, and it follows that
\begin{multline}\label{sumijorder3}
\mu(c_{u},c_{t})\mu(c_{t},c_v)\mu(c_v,c_w)\,\,\, +\mskip -25mu\sum_{\substack{x_1\in X,\ x_1<w\\x_2=(j-1)\neb(x_1)}}\mskip -25mu
\mu(c_{u},c_{x_{2}})\mu(c_{x_{2}},c_{x_{1}})\mu(c_{x_{1}},c_w)\\
=\mu(c_{u},c_{s_{j-1}s_jw})\mu(c_{s_{j-1}s_jw},c_{s_jw})\mu(c_{s_jw},c_w)\\
+\sum_{\substack{y_1\in Y,\ y_1<w\\y_2=(j-1)\neb(y_1)}}\mskip -25mu
\mu(c_{u},c_{y_{2}})\mu(c_{y_{2}},c_{y_{1}})\mu(c_{y_{1}},c_{w}),
\end{multline}
where the term $\mu(c_{u},c_{s_{j-1}s_jw})\mu(c_{s_{j-1}s_jw},c_{s_jw})\mu(c_{s_jw},c_w)$
on the right hand side of Eq.~\eqref{sumijorder3} should be omitted if $j\notin\SA(w)$.
Note that if $j\in\SA(w)$ then $(c_w,c_{s_jw},c_{s_{j-1}s_jw},c_u)$ is not necessarily a
directed path, since there need not be an arc from $c_{s_{j-1}s_jw}$ to $c_u$, but
in this case $\mu(c_{u},c_{s_{j-1}s_jw})\mu(c_{s_{j-1}s_jw},c_{s_jw})\mu(c_{s_jw},c_w)=0$
since $\mu(c_u,c_{s_{j-1}s_jw})=0$. Similarly, $(c_w,c_v,c_t,c_u)$ is not necessarily a
directed path, since there need not be an arc from $c_t$ to $c_u$, but
$\mu(c_u,c_t)\mu(c_t,c_v)\mu(c_v,c_w)=0$ in this case.
So Eq.~(\ref{sumijorder3}) still holds in these cases.

Since $\mu(c_v,c_w) = \mu(c_{s_jw},c_w)= 1$ and
$\mu(c_{t},c_v)= \mu(c_{s_{j-1}s_jw},c_{s_jw}) = 1$, by Corollary~\ref{arweightone},
and since $\mu(c_{x_{2}},c_{x_{1}})=\mu(c_{y_{2}},c_{y_{1}})=1$, since
$\{c_{x_{1}},c_{x_{2}}\}$ and $\{c_{y_{1}},c_{y_{2}}\}$ are simple edges,
Eq.~(\ref{sumijorder3}) yields the following formula for $\mu(c_{u},c_{t})$:
\begin{multline*}
\quad\mu(c_{u},c_{t})=\mu(c_{u},c_{s_{j-1}s_{j}w})\,\,+\mskip -25mu
\sum_{\substack{y_1\in Y,\ y_1<w\\y_2=(j-1)\neb(y_1)}}\mskip -25mu
\mu(c_{u},c_{y_{2}})\mu(c_{y_{1}},c_{w})\\[-10pt]
-\mskip -25mu\sum_{\substack{x_1\in X,\ x_1<w\\x_2=(j-1)\neb(x_1)}}\mskip -25mu
\mu(c_{u},c_{x_{2}})\mu(c_{x_{1}},c_w),\quad
\end{multline*}
where $\mu(c_{u},c_{s_{j-1}s_{j}w})$ should be interpreted as $0$ if $s_jw\notin\STD(\lambda)$.

Working similarly on $\Gamma_{\lambda}$ yields the following formula for $\mu^{(\lambda)}(u,t)$:
\begin{multline*}
\mu^{(\lambda)}(u,t)=\mu^{(\lambda)}(u,s_{j-1}s_{j}w)\,\,+\mskip -25mu
\sum_{\substack{y_1\in Y,\ y_1<w\\y_2=(j-1)\neb(y_1)}}\mskip -25mu
\mu^{(\lambda)}(u,y_{2})\mu^{(\lambda)}(y_{1},w)\\[-10pt]
-\mskip -25mu\sum_{\substack{x_1\in X,\ x_1<w\\x_2=(j-1)\neb(x_1)}}\mskip -25mu
\mu^{(\lambda)}(u,x_{2})\mu^{(\lambda)}(x_{1},w).
\end{multline*}
Now $w<_{\lex} t'$ and $s_{j-1}s_jw<_{\lex}t'$ (if $j\in\SA(w)$) by Lemma~\ref{xlexlessthant}~\ref{lemma7.30_iii}.
Moreover, $x_2,\,y_2<_{\lex}t'$ by Lemma~\ref{xlexlessthant}~\ref{lemma7.30_iv}. Hence
it follows by the inductive hypothesis that the corresponding arc weights in the two formulae
above are the same. Thus $\mu(c_{u'},c_{t'})=\mu^{(\lambda)}(u',t')$, as desired.
\qedhere
\end{Case}
\end{proof}

\begin{prop}\label{admcellaremonomol}
Let $\Gamma = \Gamma(C,\mu,\tau)$ be an admissible $W_n$-graph
whose molecules are all of type $\lambda$, for some $\lambda\in P(n)$.
Then there are no arcs between distinct molecules, and each of the molecules
is isomorphic to $\Gamma_\lambda$.
\end{prop}

\begin{proof}
Slightly modifying the notation used in Theorem~\ref{orderedWgraphdef},
let $\mathcal{I}$ be a set that indexes the molecules of $\Gamma$, and for
each $\alpha\in\mathcal{I}$ let $C_\alpha$ be the vertex set of the corresponding molecule.
We also write $C_\alpha=\{\,c_{\alpha,t}\mid t\in\STD(\lambda)\,\}$, and let
$\Gamma_\alpha=\Gamma(C_\alpha)$
denote the molecule spanned by~$C_\alpha$.
If $\Gamma = \bigoplus_{\alpha\in\mathcal{I}}\Gamma_\alpha$, the direct sum of
the $\Gamma_\alpha$, then each $\Gamma_\alpha$ is a monomolecular admissible
$W_n$-cell of type~$\lambda$, and hence isomorphic to
$\Gamma_{\lambda}$, by
Proposition~\ref{monomolecularadmcellsareKL}. Hence it will suffice to show
that $\Gamma = \bigoplus_{\alpha\in\mathcal{I}}\Gamma_\alpha$.

Suppose otherwise. Then there exists $\alpha\in\mathcal{I}$ such that
$\Ini_\Gamma(\alpha)\ne\emptyset$, where
\begin{equation*}
\Ini_{\Gamma}(\alpha)=\{\,t\in\STD(\lambda)\mid\mu(c_{\beta,u},c_{\alpha,t})\neq 0
\text{ for some }(\beta,u)\in(\mathcal{I}\setminus\{\alpha\})\times\STD(\lambda)\}.
\end{equation*}
Fix such an $\alpha$, and let $t'$ be the element of $\Ini_{\alpha}(\Gamma)$
that is minimal in the lexicographic order on $\STD(\lambda)$. Choose
$(\beta,u')\in (\mathcal{I}\setminus \{\alpha\})\times\STD(\lambda)$ with
$\mu(c_{\beta,u'},c_{\alpha,t'})\neq 0$. Since
$\Gamma$ satisfies the Simplicity Rule (by
Theorem~\ref{combinatorialCharacterisation}), the assumption that
$\alpha\neq\beta$ and $\mu(c_{\beta,u'},c_{\alpha,t'})\neq 0$ implies that
$\D(t')\subsetneqq\D(u')$. Moreover, since $\Gamma$ is ordered (by Theorem~\ref{inductivestep}),
$\alpha\neq\beta$ implies that $u'<t'$. Hence $(u',t')$ is a probable pair.

Let $i$ be the restriction number of $(u',t')$ and
$j=\max(\SD(t'))$. Let $(u,t)\in F(u',t')$, so that $(u,t)$ is
$i$-restricted and favourable, and Lemma~\ref{linktominimalCor} shows that
$(u,t)$ is a probable pair, $\max(\SD(t))=\max(\SD(t'))=j$,
and $\mu(c_{\beta,u},c_{\alpha,t})=\mu(c_{\beta,u'},c_{\alpha,t'})\neq 0$.
Since $\D(t)\subsetneqq\D(u)$ (since $(u,t)$ is probable) and $i\in\D(u)\oplus\D(t)$
(since $(u,t)$ is favourable), $i\in\D(u)\setminus\D(t)$.
Thus, since $j\in\D(t)\subsetneqq\D(u)$, it follows that
$j\in\D(t)$ and $i\notin\D(t)$, and
$i,\,j\in\D(u)$. Let $v=s_jt$,
and note that $v\in\STD(\lambda)$ and $v<t$, by Remark~\ref{equalsubsetsS}.
Lemma~\ref{rnltmaxsd} yields that $i<j$.

Suppose first that $i<j-1$. Then
Lemma~\ref{variousaltpaths}~\ref{lemma7.29_i} yields that $i,\,j\notin\D(v)$.
Moreover, since $\mu(c_{\alpha,t},c_{\alpha,v})=1$
by Corollary~\ref{arweightone}, and since $\mu(c_{\beta,u},c_{\alpha,t})\neq 0$,
it follows that $(c_{\alpha,v},c_{\alpha,t},c_{\beta,u})$ is
an alternating directed path of type $(j,i)$. Thus since $\Gamma$ is admissible
it follows that $N^{2}_{j,i}(c_{\alpha,v},c_{\beta,u})>0$, and so
$N^{2}_{i,j}(c_{\alpha,v},c_{\beta,u})>0$, since $\Gamma$ satisfies the
Polygon Rule.
Thus there exists at least one $(\delta,x_1)\in\mathcal{I}\times\STD(\lambda)$
such that $(c_{\alpha,v},c_{\delta,x_1},c_{\beta,u})$
is an alternating directed path of type $(i,j)$. If $\delta\neq\alpha$ then
$v\in\Ini_{\Gamma}(\alpha)$, and since $\morethan {(t')}i=\morethan ti$ we have
$v<_{\lex}t'$, by Lemma~\ref{xlexlessthant}, contradicting the definition
of~$t'$. So we must have $\delta=\alpha$, whence $x_1\in\Ini_{\Gamma}(\alpha)$. Now
Theorem~\ref{inductivestep} shows that either $x_1=s_iv$ and $i\in\SA(v)$, or else
$x_1<v$. But $x_1<_{\lex}t'$ by Lemma~\ref{xlexlessthant}~\ref{lemma7.30_i}
in the former case, and $x_1<_{\lex}t'$ by Lemma~\ref{xlexlessthant}~\ref{lemma7.30_ii} in
the latter case. Both alternatives contradict the definition of~$t'$,
and we conclude that $i<j-1$ is impossible.

Thus $i=j-1$, and it now follows from Lemma~\ref{existencef} that
$\col_{t}(j-1)\neq\col_{t}(j+1)$. But if $\col_{t}(j-1)<\col_{t}(j+1)$ then
we obtain a contradiction by the same reasoning as in the $i<j-1$ case,
using Lemma~\ref{variousaltpaths}~\ref{lemma7.29_ii} in place of Lemma~\ref{variousaltpaths}~\ref{lemma7.29_i}.
So $\col_{t}(j-1)>\col_{t}(j+1)$.

Let $w=s_{j-1}v$. Lemma~\ref{variousaltpaths}~\ref{lemma7.29_iii} gives
$j-1\in\SD(v)$, whence $w\in\STD(\lambda)$ and $w<v$.
Moreover, Lemma~\ref{variousaltpaths}~\ref{lemma7.29_iii} also gives
$j-1,\,j\notin\D(w)$, as well as $j-1\in\D(v)$ and $j\notin\D(v)$.
Next, since $\mu(c_{\alpha,v},c_{\alpha,w})=\mu(c_{\alpha,t},c_{\alpha,v})=1$
by Corollary~\ref{arweightone} and since $\mu(c_{\beta,u},c_{\alpha,t})\neq 0$,
it follows that
$(c_{\alpha,w},c_{\alpha,v},c_{\alpha,t},c_{\beta,u})$ is
an alternating directed path of type $(j-1,j)$.
Thus $N^{3}_{j-1,j}(c_{\alpha,w},c_{\beta,u})>0$, and so
$N^{3}_{j,j-1}(c_{\alpha,w},c_{\beta,u})>0$, by the Polygon Rule.

It follows that there exists at least one $(\delta,x_1)\in\mathcal{I}\times\STD(\lambda)$
and one $(\gamma,x_2)\in\mathcal{I}\times\STD(\lambda)$
such that $(c_{\alpha,w},c_{\delta,x_1},c_{\gamma,x_2},c_{\beta,u})$
is an alternating directed path of type $(j,j-1)$.
If $\delta\neq\alpha$ then $w\in\Ini_{\alpha}(\Gamma)$, and
since $\morethan t{(j-1)}=\morethan {t'}{(j-1)}$, we have $w<_{\lex}t'$ by
Lemma~\ref{xlexlessthant}~\ref{lemma7.30_iii}. But this contradicts the definition of~$t'$,
and so $\delta=\alpha$.

Since
$\D(x_1)\cap\{j-1,j\}=\{j\}$ and $\D(x_2)\cap\{j-1,j\}=\{j-1\}$,
and $\mu(c_{\gamma,x_2},c_{\delta,x_1})\neq 0$, it follows from the
Simplicity Rule
that $\{c_{\delta,x_1},c_{\gamma,x_2}\}$ is a simple edge.
Thus $\gamma=\delta$, and $x_1$ and $x_2$ are related by a DKM.
Thus $x_2$ is the $(j-1)$-neighbour of $x_1$, and $x_2\in\Ini_{\alpha}(\Gamma)$.
It will suffice to show that $x_2<_{\lex}t'$, contradicting the definition of~$t'$.

By Theorem~\ref{inductivestep} either $x_1=s_jw>w$ or $x_1<w$. If
$x_1<w$ then since $\morethan t{(j-1)}=\morethan {(t')}{(j-1)}$, the conclusion
$x_2<_{\lex}t'$ follows from Lemma~\ref{xlexlessthant}~\ref{lemma7.30_iv}.
We are left with the case $x_1=s_jw>w$. This gives $j\in\SA(w)$, and we see
that the conditions of Lemma~\ref{variousaltpaths}~\ref{lemma7.29_iii} are satisfied: we
have $v=s_jt$ with $j\in\SD(t)$ and $\col_t(j+1)<\col_t(j-1)$, and $w=s_{j-1}v$.
Since $j\in\SA(w)$ it follows that $j-1\in\SA(x_1)$, and $s_{j-1}x_1$
is the $(j-1)$-neighbour of~$x_1$. Thus $x_2=s_{j-1}x_1=s_{j-1}s_jw$, and
since $\morethan t{(j-1)}=\morethan {(t')}{(j-1)}$, we have
$x_2<_{\lex}t'$ by Lemma~\ref{xlexlessthant}~\ref{lemma7.30_iii}.
\end{proof}

\begin{rema}
In the above proof, after noting that
$\mu(c_{\beta,u'},c_{\alpha,t'})\neq 0$ and $\alpha\neq\beta$,
Remark~\ref{nonemptyAkst} can be used to deduce that $A(u',t')\neq\emptyset$, and
then, choosing $(u,t)\in A(u',t')$, a proof similar to that of
Proposition~\ref{cellorder} (with $t'$ replacing $t_\lambda$) shows
that $\Morethan ti$ is $i$-critical.
But then Lemma~\ref{rnltmaxsd} and Definition~\ref{m-critical} combined show that
$i+1=\max(\SD(t))$ and $\col_t(i+2)=\col_t(i)$, and this contradicts Lemma~\ref{existencef}.
\end{rema}

We are now in a position to state and prove the main result of the paper.
\begin{theo}\label{mainresult}
Admissible cells of type $A_{n-1}$ are Kazhdan--Lusztig.
\end{theo}

\begin{proof}
Let $\Gamma = \Gamma(C,\mu,\tau)$ be an admissible $W_n$-cell, and let
$\Lambda$ be the set of molecule types for $\Gamma$. By Lemma~\ref{samelambda},
$\Lambda=\{\lambda\}$ for some $\lambda\in P(n)$.
Let $\mathcal{I}$ index
the molecules of $\Gamma$, and for each $\gamma\in\mathcal{I}$ let
$C_{\gamma}= \{c_{\gamma,w}\mid w\in\STD(\lambda)\}$
be the vertex set of the corresponding molecule. Then Proposition~\ref{admcellaremonomol}
says that $\Gamma = \bigoplus_{\gamma\in\mathcal{I}}\Gamma(C_{\gamma})$,
with each $\Gamma(C_{\gamma})$
isomorphic to $\Gamma_{\lambda}$. But $\Gamma$ must be connected, since it
is a cell, and so $\mathcal{I}$ has only one element.
Thus $\Gamma$ is isomorphic to~$\Gamma_{\lambda}$.
Since $\Gamma_{\lambda}$ is isomorphic to $\Gamma(C(\tau_{\lambda}))$, it follows from
Corollary~\ref{leftcelllambda} that $\Gamma$ is isomorphic to a Kazhdan--Lusztig left cell.
\end{proof}

\begin{rema}\label{remontwosidedcells}
Let $\lambda\in P(n)$ and let $D(\lambda)=\bigsqcup_{t\in\STD(\lambda)}C(t)$,
the Kazhdan--Lusztig two-sided cell corresponding to $\lambda$.
By Remark~\ref{orderedtwosided}, the singleton set $\{\lambda\}$ is the set
of molecule types of the admissible $W_n$-graph $\Gamma(D(\lambda))$. It follows
from Proposition~\ref{admcellaremonomol} that $\Gamma(D(\lambda))$ is a disjoint union of
the Kazhdan--Lusztig left cells $\Gamma(C(t))$. This implies the following well known
result (see, for example, \cite[Theorem~5.3]{gecpfei:charHecke}).
\begin{theo}\label{xyinsameleftcell}
Let $\lambda\in P(n)$ and $y,w\in D(\lambda)$. If $y\preceq\lside w$
then $y,w\in C(t)$ for some $t\in\STD(\lambda)$.
\end{theo}
\end{rema}

\longthanks
I thank the referees for valuable suggestions, and Professor Robert B. Howlett for many improvements to the
exposition of this paper.

\nocite{*}
\bibliographystyle{amsplain-ac}
\bibliography{ALCO_Nguyen_205}

\end{document}

